% ----------------------
\subsection{Mosiah}
% ----------------------
	\label{MOSIAH}

	% -----
	\subsubsection{Mosiah 1} % *MOSIAH-1
	% -----
		\label{MOSIAH-1}
	
		% ---
		\paragraph{Mosiah 1:1} % *MOSIAH-1-1 
		% ---
			\label{MOSIAH-1-1}

			And now there was no more contention in all the land of Zarahemla
			among all the people which belonged to king Benjamin,
			so that king Benjamin had continual peace all the remainder of his days.

		% ---
		\paragraph{Mosiah 1:2} % *MOSIAH-1-2 
		% ---
			\label{MOSIAH-1-2}

			And it came to pass that he had three sons,
			and he called their names Mosiah and Helorum and Helaman.
			And he caused that they should be taught in all the language of his fathers,
			that thereby they might become men of understanding
			and that they might know concerning the prophecies
			which had been spoken by the mouths of their fathers,
			which was delivered them by the hand of the Lord.

		% ---
		\paragraph{Mosiah 1:3} % *MOSIAH-1-3 
		% ---
			\label{MOSIAH-1-3}

			And he also taught them concerning the records
			which were engraven on the plates of brass, saying:
			My sons, I would that ye should remember
			that were it not for these plates
			which contain these records and these commandments,
			we must have suffered in ignorance, even at this present time,
			not knowing the mysteries of God.

		% ---
		\paragraph{Mosiah 1:4} % *MOSIAH-1-4 
		% ---
			\label{MOSIAH-1-4}

			For it were not possible
			that our father Lehi could have remembered all these things,
			to have taught them to his children,
			except it were for the help of these plates;
			for he having been taught in the language of the Egyptians,
			therefore he could read these engravings and teach them to his children,
			that thereby they could teach them to their children,
			and so fulfilling the commandments of God, even down to this present time.

		% ---
		\paragraph{Mosiah 1:5} % *MOSIAH-1-5 
		% ---
			\label{MOSIAH-1-5}

			I say unto you my sons:
			Were it not for these things
			which have been kept and preserved by the hand of God,
			that we might read and understand of his mysteries
			and have his commandments always before our eyes,
			that even our fathers would have dwindled in unbelief,
			and we should have been like unto our brethren the Lamanites,
			which know nothing concerning these things,
			or even do not believe them when they are taught them
			because of the traditions of their fathers, which are not correct.

		% ---
		\paragraph{Mosiah 1:6} % *MOSIAH-1-6 
		% ---
			\label{MOSIAH-1-6}

			O my sons, I would that ye should remember
			that these sayings are true,
			and also that these records are true.
			And behold also the plates of Nephi
			which contain the records and the sayings of our fathers
			from the time they left Jerusalem until now,
			and they are true;
			and we can know of their surety
			because we have them before our eyes.

		% ---
		\paragraph{Mosiah 1:7} % *MOSIAH-1-7 
		% ---
			\label{MOSIAH-1-7}

			And now my sons, I would that ye should remember to search them diligently,
			that ye may profit thereby.
			And I would that ye should keep the commandments of God,
			that ye may prosper in the land
			according to the promises which the Lord made unto our fathers.

		% ---
		\paragraph{Mosiah 1:8} % *MOSIAH-1-8 
		% ---
			\label{MOSIAH-1-8}

			And many more things did king Benjamin teach his sons
			which are not written in this book.

		% ---
		\paragraph{Mosiah 1:9} % *MOSIAH-1-9 
		% ---
			\label{MOSIAH-1-9}

			And it came to pass that
			after king Benjamin had made an end of teaching his sons
			that he waxed old,
			and he saw that he must very soon go the way of all the earth.
			Therefore he thought it expedient
			that he should confer the kingdom upon one of his sons;

		% ---
		\paragraph{Mosiah 1:10} % *MOSIAH-1-10 
		% ---
			\label{MOSIAH-1-10}

			therefore he had Mosiah brought before him,
			and these are the words which he spake unto him, saying:
			My son, I would that ye should make a proclamation
			throughout all this land among all this people,
			or the people of Zarahemla and the people of Mosiah which dwell in this land,
			that thereby they may be gathered together,
			for on the morrow I shall proclaim unto this my people out of mine own mouth
			that thou art a king and a ruler over this people,
			which the Lord our God hath given us.

		% ---
		\paragraph{Mosiah 1:11} % *MOSIAH-1-11 
		% ---
			\label{MOSIAH-1-11}

			And moreover I shall give this people a name
			that thereby they may be distinguished above all the people
			which the Lord God hath brought out of the land of Jerusalem.
			And this I do because they have been a diligent people
			in keeping the commandments of the Lord.

		% ---
		\paragraph{Mosiah 1:12} % *MOSIAH-1-12 
		% ---
			\label{MOSIAH-1-12}

			And I give unto them a name that never shall be blotted out
			except it be through transgression.

		% ---
		\paragraph{Mosiah 1:13} % *MOSIAH-1-13 
		% ---
			\label{MOSIAH-1-13}

			Yea, and moreover I say unto you that
			if this highly favored people of the Lord should fall into transgression
			and become a wicked and an adulterous people
			that the Lord will deliver them up,
			that thereby they become weak like unto their brethren,
			and he will no more preserve them by his matchless and marvelous power
			as he hath hitherto preserved our fathers.

		% ---
		\paragraph{Mosiah 1:14} % *MOSIAH-1-14 
		% ---
			\label{MOSIAH-1-14}

			For I say unto you that
			if he had not extended his arm in the preservation of our fathers,
			they must have fallen into the hands of the Lamanites
			and become victims to their hatred.

		% ---
		\paragraph{Mosiah 1:15} % *MOSIAH-1-15 
		% ---
			\label{MOSIAH-1-15}

			And it came to pass that
			after king Benjamin had made an end of these sayings to his son
			that he gave him charge concerning all the affairs of the kingdom.

		% ---
		\paragraph{Mosiah 1:16} % *MOSIAH-1-16 
		% ---
			\label{MOSIAH-1-16}

			And moreover he also gave him charge concerning the records
			which were engraven on the plates of brass,
			and also the plates of Nephi,
			and also the sword of Laban,
			and the ball or director which led our fathers through the wilderness,
			which was prepared by the hand of the Lord
			that thereby they might be led, every one,
			according to the heed and diligence which they gave unto him.

		% ---
		\paragraph{Mosiah 1:17} % *MOSIAH-1-17 
		% ---
			\label{MOSIAH-1-17}

			Therefore as they were unfaithful,
			they did not prosper nor progress in their journey,
			but were driven back and incurred the displeasure of God upon them;
			and therefore they were smitten with famine and sore afflictions,
			to stir them up in remembrance of their duty.

		% ---
		\paragraph{Mosiah 1:18} % *MOSIAH-1-18 
		% ---
			\label{MOSIAH-1-18}

			And now it came to pass that
			Mosiah went and did as his father had commanded him
			and proclaimed unto all the people which were in the land of Zarahemla
			that thereby they might gather themselves together
			to go up to the temple to hear the words which his father should speak unto them.
			
	% -----
	\subsubsection{Mosiah 2} % *MOSIAH-2
	% -----
		\label{MOSIAH-2}
	
		% --
		\paragraph{Mosiah 2:1} % *MOSIAH-2-1
		% ---
			\label{MOSIAH-2-1}
			
			And it came to pass that after Mosiah had done as his father had commanded him
			and had made a proclamation throughout all the land
			that the people gathered themselves together throughout all the land,
			that they might go up to the temple to hear the words
			which king Benjamin should speak unto them.
			
		% --
		\paragraph{Mosiah 2:2} % *MOSIAH-2-2
		% ---
			\label{MOSIAH-2-2}
			
			And there were a great number,
			even so many that they did not number them,
			for they had multiplied exceedingly and waxed great in the land.
			
		% --
		\paragraph{Mosiah 2:3} % *MOSIAH-2-3
		% ---
			\label{MOSIAH-2-3}
			
			And they also took of the firstlings of their flocks,
			that they might offer sacrifice and burnt offerings according to the law of Moses.%
				\footnote{%
					\label{MOSIAH-2-3-NOTE-offerings}
					Note that the people went up ready to participate in temple ordinances as part of their worship. It was more for them than just listening to the prophet.
					}%
			
		% --
		\paragraph{Mosiah 2:4} % *MOSIAH-2-4
		% ---
			\label{MOSIAH-2-4}
			
			And also that they might give thanks to the Lord their God,
			who had brought them out of the land of Jerusalem,
			and who had delivered them out of the hands of their enemies
			and had appointed just men to be their teachers
			and also a just man to be their king,
			who had established peace in the land of Zarahemla,
			and who had taught them to keep the commandments of God,
			that thereby they might rejoice and be filled with love towards God and all men.
			
		% --
		\paragraph{Mosiah 2:5} % *MOSIAH-2-5
		% ---
			\label{MOSIAH-2-5}
			
			And it came to pass that when they came up to the temple,
			they pitched their tents round about,
			every man according to his family,
			consisting of his wife and his sons and his daughters%
				\footnote{%
					\label{MOSIAH-2-5-NOTE-daughters}%
					Maybe ``daughters'' here refers to daughters-in-law? Not sure how it would work otherwise, since the text is clear that each family is separate from each other.
					}
			and their sons and their daughters,
			from the eldest down to the youngest,
			every family being separate one from another.
			
		% --
		\paragraph{Mosiah 2:6} % *MOSIAH-2-6
		% ---
			\label{MOSIAH-2-6}
			
			And they pitched their tents round about the temple,
			every man having his tent with the door thereof towards the temple,
			that thereby they might remain in their tents
			and hear the words which king Benjamin should speak unto them;
			
		% --
		\paragraph{Mosiah 2:7} % *MOSIAH-2-7
		% ---
			\label{MOSIAH-2-7}
			
			for the multitude being so great
			that king Benjamin could not teach them all within the walls of the temple,
			therefore he caused a tower to be erected,
			that thereby his people might hear the words which he should speak unto them.
			
		% --
		\paragraph{Mosiah 2:8} % *MOSIAH-2-8
		% ---
			\label{MOSIAH-2-8}
			
			And it came to pass that
			he began to speak to his people from the tower,
			and they could not all hear his words because of the greatness of the multitude.
			Therefore he caused that the words which he spake should be written
			and sent forth among those that were not under the sound of his voice,
			that they might also receive his words.
			
		% --
		\paragraph{Mosiah 2:9} % *MOSIAH-2-9
		% ---
			\label{MOSIAH-2-9}
			
			And these are the words which he spake and caused to be written, saying:
			My brethren, all ye that have assembled yourselves together,
			you that can hear my words which I shall speak unto you this day,
			for I have not commanded you to come up hither
			to trifle%
				\footnote{%
					% \label{}%
					For ``trifle,'' see \hyperref[DC32-5]{DC 32:5} and notes there.
					} 
			with the words which I shall speak,
			but that you should hearken unto me
			and open your ears that ye may hear
			and your hearts that ye may understand
			and your minds that the mysteries of God may be unfolded to your
				\footnote{%
					\label{MOSIAH-2-9-NOTE-view}%
					Here Benjamin says that the mysteries of God will be unfolded to the ``view'' of those listening. The word ``view'' then appears in three other places in this story---
						\begin{compactenum}
							\item Mosiah 3:25: ``And if they [a person's works] be evil, they are consigned to an awful \textbf{view} of their own guilt and abominations, which doth cause them to shrink from the presence of the Lord into a state of misery and endless torment, from whence they can no more return;''
							\item Mosiah 4:2: ``And they had \textbf{viewed} themselves in their own carnal state, even less than the dust of the earth, and they all cried aloud with one voice, saying: O have mercy and apply the atoning blood of Christ that we may receive forgiveness of our sins and our hearts may be purified, for we believe in Jesus Christ the Son of God, who created heaven and earth and all things, who shall come down among the children of men.''
							\item Mosiah 5:3: ``And we ourselves also through the infinite goodness of God and the manifestations of his Spirit have great \textbf{views} of that which is to come. And were it expedient, we could prophesy of all things.''
						\end{compactenum}
					All three of the things that the people experience and learn here are things that are revealed to us by God, and so count as the ``mysteries'' that Benjamin is referring to. In Mosiah 3:25, Benjamin is talking about the time at or near the judgment at which a person remembers all their evil works (2 Nephi 9:13--16). In Mosiah 4:2, the people learn how badly off they are in their carnal state, something that wouldn't occur to you, much less matter, unless God was revealing it. And in Mosiah 5:3, the people are looking forward to good things in the future, which have been revealed to them through the Spirit (the text explicitly states this).
					So in the ``mysteries'' that are unfolded to this group of people we see a progression, from a knowledge of their need for repentance---which they respond to by repenting---to much better things in the future. I think we will go through the same progression in seeking the new birth. Think of how the ``unfolding'' process occurs, where some parts that are visible are then no longer visible once a larger portion has unfolded; maybe some of those parts that aren't visible any longer, once we are ``viewing'' the good things in the future, are our sins and carnal nature.
					}%
			view.
			
		% --
		\paragraph{Mosiah 2:10} % *MOSIAH-2-10
		% ---
			\label{MOSIAH-2-10}
			
			I have not commanded you to come up hither that ye should fear me,
			or that ye should think that I of myself am more than a mortal man.
			
		% --
		\paragraph{Mosiah 2:11} % *MOSIAH-2-11
		% --
			\label{MOSIAH-2-11}
			
			But I am like as yourselves,
			subject to all manner of infirmities in body and mind.
			Yet as I have been chosen by this people
			and was consecrated by my father
			and was suffered by the hand of the Lord
			that I should be a ruler and a king over this people
			and have been kept and preserved by his matchless power,
			to serve thee with all the might, mind, and strength
			which the Lord hath granted unto me---
			
		% --
		\paragraph{Mosiah 2:12} % *MOSIAH-2-12
		% --
			\label{MOSIAH-2-12}
			
			I say unto you that as I have been suffered
			to spend my days in your service, even up to this time,
			and have not sought gold nor silver
			nor no manner of riches of you,
			
		% --
		\paragraph{Mosiah 2:13} % *MOSIAH-2-13
		% --
			\label{MOSIAH-2-13}
			
			neither have I suffered that ye should be confined in dungeons
			nor that ye should make slaves one of another
			or that ye should murder or plunder or steal or commit adultery,
			or even I have not suffered that ye should commit any manner of wickedness,
			and have taught you that ye should keep the commandments of the Lord
			in all things which he hath commanded you---
			
		% --
		\paragraph{Mosiah 2:14} % *MOSIAH-2-14
		% --
			\label{MOSIAH-2-14}
			
			and even I myself have labored with mine own hands that I might serve you
			and that ye should not be laden with taxes
			and that there should nothing come upon you which was grievious to be borne---
			and of all these things which I have spoken,
			ye yourselves are witnesses this day.
			
		% --
		\paragraph{Mosiah 2:15} % *MOSIAH-2-15
		% --
			\label{MOSIAH-2-15}
			
			Yet my brethren, I have not done these things that I might boast,
			neither do I tell these things that thereby I might accuse you.
			But I tell you these things that ye may know
			that I can answer a clear conscience before God this day.
			
		% --
		\paragraph{Mosiah 2:16} % *MOSIAH-2-16
		% --
			\label{MOSIAH-2-16}
			
			Behold, I say unto you that
			because I said unto you that I had spent my days in your service,
			I do not desire to boast,
			for I have only been in the service of God.
			
		% --
		\paragraph{Mosiah 2:17} % *MOSIAH-2-17
		% --
			\label{MOSIAH-2-17}
			
			And behold, I tell you these things that ye may learn wisdom,
			that ye may learn that when ye are in the service of your fellow beings,
			ye are only in the service of your God.
			
		% --
		\paragraph{Mosiah 2:18} % *MOSIAH-2-18
		% --
			\label{MOSIAH-2-18}
			
			Behold, ye have called me your king.
			And if I, whom ye call your king, do labor to serve you,
			then had not ye ought to labor to serve one another?
			
		% --
		\paragraph{Mosiah 2:19} % *MOSIAH-2-19
		% --
			\label{MOSIAH-2-19}
			
			And behold also, if I, who ye call your king,
			who has spent his days in your service
			and yet hath been in the service of God,
			doth merit any thanks from you,
			O how had you ought to thank your heavenly King!
			
		% --
		\paragraph{Mosiah 2:20} % *MOSIAH-2-20
		% --
			\label{MOSIAH-2-20}
			
			I say unto you my brethren
			that if you should render all the thanks and praise
			which your whole souls hath power to possess
			to that God who hath created you
			and hath kept and preserved you
			and hath caused that ye should rejoice
			and hath granted that ye should live in peace one with another---
			
		% --
		\paragraph{Mosiah 2:21} % *MOSIAH-2-21
		% --
			\label{MOSIAH-2-21}
			
			I say unto you that
			if ye should serve him who hath created you from the beginning
			and art preserving you from day to day by lending you breath
			that ye may live and move and do according to your own will,
			and even supporting you from one moment to another---%
				\footnote{%
					\label{MOSIAH-2-21-NOTE-supporting}%
					This passage here, more than any other in the restoration scriptures that I've found yet, reminds me of medieval and Cartesian views about God's active role in maintaining how things are in the world. I think that later thoughts from JS and others push back against this; I think I'm inclined to explain in terms of the laws that govern things like respiration, which then fit into a larger picture of laws of the priesthood.
					}
			I say, if ye should serve him with all your whole soul,
			and yet ye would be unprofitable servants.
			
		% --
		\paragraph{Mosiah 2:22} % *MOSIAH-2-22
		% --
			\label{MOSIAH-2-22}
			
			And behold, all that he requires of you is to keep his commandments.
			And he hath promised you that if ye would keep his commandments,
			ye should prosper in the land.
			And he never doth vary from that which he hath said;
			therefore if ye do keep his commandments,
			he doth bless you and prosper you.
			
		% --
		\paragraph{Mosiah 2:23} % *MOSIAH-2-23
		% --
			\label{MOSIAH-2-23}
			
			And now in the first place, he hath created you
			and granted unto you your lives,
			for which ye are indebted unto him.
			
		% --
		\paragraph{Mosiah 2:24} % *MOSIAH-2-24
		% --
			\label{MOSIAH-2-24}
			
			And secondly, he doth require
			that ye should do as he hath commanded you,
			for which if ye do,
			he doth \hyperref[IMMEDIATE]{immediately}%
				\footnote{%
					% \label{MOSIAH-2-24-NOTE-immediately}%
					See \hyperref[IMMEDIATE]{here} for the 1828 definitions of ``immediate''; note that the third one is the one that we normally use nowadays, talking about time, but the first definition in the 1828 (and the one that was no doubt much more common during JS's time) is that of ``acting directly'' or ``acting without a medium.'' Compare also \hyperref[JOHN-3-18]{John 3:18} and ``already'' in that verse, and see \hyperref[MOSIAH-25-10]{Mosiah 25:10}, \hyperref[ALMA-15-4]{Alma 15:4--5} (where ``immediately'' seems to be used the way we use it). See also Nelson, SAM 170: ``Now the soul that has thus repented---howso­ever contradictory it may seem---is absolved by the very fact of its repentance. It has complied with the law, and must therefore inevitably secure the reward. No power, not even God Himself, can rob it of the security of its newly-won relationship to the universe. It is inevitable in the nature of things that obedience to law carries with it absolu­tion from former opposition to law. And this equilibrium is instant: we need fear no locked-up vengeance to be sprung upon us---no reprisals to satisfy an offended God.''
					}
			bless you,
			and therefore he hath paid you.
			And ye are still indebted unto him
			and are and will be forever and ever.
			Therefore of what have ye to boast?
			
		% --
		\paragraph{Mosiah 2:25} % *MOSIAH-2-25
		% --
			\label{MOSIAH-2-25}
			
			And now I ask:
			Can ye say \hyperref[AUGHT-1828]{aught} of yourselves?
			I answer you: Nay.
			Ye cannot say that thou art even as much as the dust of the earth,
			yet thou wast created of the dust of the earth;
			but behold, it belongeth to him who created you.
			
		% --
		\paragraph{Mosiah 2:26} % *MOSIAH-2-26
		% --
			\label{MOSIAH-2-26}
			
			And I, even I, whom ye call your king, am no better than ye yourselves are,
			for I am also of the dust.
			And thou beholdest that I am old
			and am about to yield up this mortal frame to its mother earth.
			
		% --
		\paragraph{Mosiah 2:27} % *MOSIAH-2-27
		% --
			\label{MOSIAH-2-27}
			
			Therefore as I said unto you
			that I had served you, walking with a clear conscience before God,
			even so I at this time have caused that ye should assemble yourselves together,
			that I might be found blameless
			and that your blood should not come upon me
			when I shall stand to be judged of God of the things
			whereof he hath commanded me concerning you.
			
		% --
		\paragraph{Mosiah 2:28} % *MOSIAH-2-28
		% --
			\label{MOSIAH-2-28}
			
			 I say unto you that I have caused that ye should assemble yourselves together
			that I might rid my garments of your blood,%
				\footnote{%
					\label{MOSIAH-2-28-NOTE-assemble}%
					Am I the only one who thinks this is a weird reason to get everyone together? So that he can totally exonerate himself of all responsibility for their actions before he dies? As he notes in the next verse, he does have multiple reasons for getting people together, but this one just sounds a little strange to me.
					}
			at this period of time when I am about to go down to my grave,
			that I might go down in peace
			and my immortal spirit may join the choirs above
			in singing the praises of a just God.
		
		% --
		\paragraph{Mosiah 2:29} % *MOSIAH-2-29
		% --
			\label{MOSIAH-2-29}
			
			And moreover I say unto you
			that I have caused that ye should assemble yourselves together
			that I might declare unto you
			that I can no longer be your teacher nor your king.
		
		% --
		\paragraph{Mosiah 2:30} % *MOSIAH-2-30
		% --
			\label{MOSIAH-2-30}
			
			For even at this time my whole frame doth tremble exceedingly
			while attempting to speak unto you,
			but the Lord God doth support me
			and hath suffered me that I should speak unto you
			and hath commanded me that I should declare unto you this day
			that my son Mosiah is a king and a ruler over you.
			
		% --
		\paragraph{Mosiah 2:31} % *MOSIAH-2-31 
		% ---
			\label{MOSIAH-2-31}

			And now my brethren, I would that ye should do as ye hath hitherto done;
			as ye have kept my commandments,
			and also the commandments of my father,
			and have prospered and have been kept
			from falling into the hands of your enemies,
			even so if ye shall keep the commandments of my son,
			or the commandments of God which shall be delivered unto you by him,
			ye shall prosper in the land,
			and your enemies shall have no power over you.

		% --
		\paragraph{Mosiah 2:32} % *MOSIAH-2-32 
		% ---
			\label{MOSIAH-2-32}

			But O my people, beware lest there shall arise contentions among you,
			and ye \hyperref[LIST-1828]{list}%
				\footnote{%
					\label{MOSIAH-2-32-NOTE-list}%
					This unusual use of the word ``list'' is the first time we see this meaning in the BOM. The word with this meaning occurs seven times in the \hyperref[LIST-BOM]{BOM}, and three times in the \hyperref[LIST-DC]{DC}.
					}
			to obey the evil spirit which was spoken of by my father Mosiah.

		% --
		\paragraph{Mosiah 2:33} % *MOSIAH-2-33 
		% ---
			\label{MOSIAH-2-33}

			For behold, there is a woe pronounced upon him who \hyperref[LIST-1828]{listeth} to obey that spirit;
			for if he \hyperref[LIST-1828]{listeth} to obey him and remaineth and dieth in his sins,
			the same drinketh damnation to his own soul,
			for he receiveth for his wages an everlasting punishment,
			having transgressed the law of God contrary to his own knowledge.

		% --
		\paragraph{Mosiah 2:34} % *MOSIAH-2-34 
		% ---
			\label{MOSIAH-2-34}

			I say unto you that there are not one among you
			---except it be your little children---
			that have not been taught concerning these things
			but what knoweth that ye are eternally indebted to your heavenly Father,
			to render to him all that you have and are,
			and also have been taught concerning the records,
			which contain the prophecies which hath been spoken by the holy prophets,
			even down to the time our father Lehi left Jerusalem,

		% --
		\paragraph{Mosiah 2:35} % *MOSIAH-2-35 
		% ---
			\label{MOSIAH-2-35}

			and also all that hath been spoken by our fathers until now.
			And behold also, they spake that which was commanded them of the Lord;
			therefore they are just and true.

		% --
		\paragraph{Mosiah 2:36} % *MOSIAH-2-36 
		% ---
			\label{MOSIAH-2-36}

			And now I say unto you my brethren that
			after ye have known and have been taught all these things,
			if ye should transgress and go contrary to that which hath been spoken,
			that ye do \hyperref[WITHDRAW]{withdraw} yourselves from the Spirit of the Lord,%
				\footnote{%
					% \label{MOSIAH-2-36-NOTE-withdraw}%
					I think this verse is helpful in developing the correct way of thinking about ``having'' the Spirit. It isn't like an object that you can either possess or not possess at any given time. I think a more correct way of thinking about it, and a more useful image, is that of standing closer or farther away from a source of light and warmth, like a fire or lantern. When you stand close, you feel the warmth and can see better because there is more light, but you could also stand farther away and feel less warmth and see more poorly. The Spirit (or Christ, or whatever you want to plug into this image, it doesn't really matter) is the source of light and warmth, especially for our spirits and minds. The closer we want to approach---and how close we approach is almost entirely up to us---then the more light and warmth we will be able to receive. Or we can withdraw, as this verse suggests, and get less. But it isn't the case that we are ever completely out of range of the light and warmth; we just sort of banish ourselves to a position very far from the source, and as a result are naturally able to receive very little of it.
					
					Consider also \hyperref[DC88-63]{DC 88:63} in connection with this.
					}
			that it may have no place in you to guide you in wisdom's paths
			that ye may be blessed, prospered, and preserved---

		% --
		\paragraph{Mosiah 2:37} % *MOSIAH-2-37 
		% ---
			\label{MOSIAH-2-37}

			I say unto you that the man that doeth this,
			the same cometh out in open rebellion against God;
			therefore he \hyperref[LIST-1828]{listeth} to obey the evil spirit
			and becometh an enemy to all righteousness.%
				\footnote{%
					\label{MOSIAH-2-37-NOTE-enemy}%
					See \hyperref[MORONI-9-6-NOTE-enemy]{Moroni 9:6} for this interesting phrase, ``enemy to all righteousness.''
					}
			Therefore the Lord hath no place in him,%
				\footnote{%
					\label{MOISAH-2-37-NOTE-place}%
					Compare \hyperref[ALMA-34-34]{Alma 34:34--36}, which also uses the term ``withdraw'' that we see earlier in verse 36 of this chapter. We are the ones who control what occupies this place inside ourselves; see the ``give place'' construction in \hyperref[ALMA-32-27]{Alma 32:27--28}.
					}
			for he dwelleth not in unholy temples.

		% --
		\paragraph{Mosiah 2:38} % *MOSIAH-2-38 
		% ---
			\label{MOSIAH-2-38}

			Therefore if that man repenteth not
			and remaineth and dieth an enemy to God,
			the demands of divine justice doth awaken his immortal soul
			to a lively sense%
				\footnote{%
					% \label{MOSIAH-2-38-NOTE-lively-sense}%
					See \hyperref[LIVELY-1828]{here} for definitions of ``lively''; I think the most apparent meanings are \#1, \#3, and \#5. The word here in ``lively sense'' (which is the only use of this phrase in the scriptures, by the way) means ``strong'' and ``vigorous'' (meanings \#5 and \#1), but also ``living'' or ``vivid'' in the sense of being powerful and present, as something alive (meaning \#3). It's interesting to see this use of the word and compare it to what is said in \hyperref[2NEPHI-9-13]{2 Nephi 9:13--14}, where we learn about the ``perfect knowledge'' we'll have of our guilt; in verse 13, immediately before referring to the ``perfect knowledge,'' it says that ``they are living souls'' and that the knowledge will be ``like unto us in the flesh.'' ``Living'' and ``flesh'' have either direct or indirect connections to life, and so there is something \textit{alive} and \textit{powerful} about this recollection that we will have at that time---depending on how we have lived.
					}
			of his own guilt,
			which doth cause him to shrink from the presence of the Lord
			and doth fill his breast with guilt and pain and anguish,
			which is like an unquenchable fire whose flames ascendeth up forever and ever.%
				\footnote{%
					% \label{MOSIAH-2-38-NOTE-unquenchable}%
					The punishment is psychological, and self-inflicted in the sense that it seems to be about us turning inward to the hell of our own choices and guilt. See also \hyperref[DC19-20]{DC 19:20}, where the Lord talks about Martin Harris experiencing a bit of this when he (the Lord) withdrew the Spirit from him.
					
					See also JS, \hyperref[EMS-1843-JS-6-8-APR-1843-WC-B]{6--8 April 1843} (``There is no pain so awful as the pain of suspense. This is the condemnation of the wicked; Their doubt and anxiety and suspense causes weeping, wailing and gnashing of teeth''); 
					}

		% --
		\paragraph{Mosiah 2:39} % *MOSIAH-2-39 
		% ---
			\label{MOSIAH-2-39}

			And now I say unto you that mercy hath no claim on that man;
			therefore his final doom is to endure a never-ending torment.

		% --
		\paragraph{Mosiah 2:40} % *MOSIAH-2-40 
		% ---
			\label{MOSIAH-2-40}

			O all ye old men and also ye young men
			and you little children which can understand my words
			---for I have spoken plain unto you that ye might understand---
			I pray that ye should awake to a remembrance of the awful situation
			of those that have fallen into transgression.

		% --
		\paragraph{Mosiah 2:41} % *MOSIAH-2-41 
		% ---
			\label{MOSIAH-2-41}

			And moreover I would desire
			that ye should consider on the blessed and happy state
			of those that keep the commandments of God;
			for behold, they are blessed in all things, both temporal and spiritual.
			And if they hold out faithful to the end,
			they are received into heaven,
			that thereby they may dwell with God in a state of never-ending happiness.
			O remember, remember that these things are true,
			for the Lord God hath spoken it.
			
	% -----
	\subsubsection{Mosiah 3} % *MOSIAH-3
	% -----
		\label{MOSIAH-3}
		
		% --
		\paragraph{Mosiah 3:1} % *MOSIAH-3-1 
		% ---
			\label{MOSIAH-3-1}

			And again my brethren, I would call your attention,
			for I have somewhat more to speak unto you;
			for behold, I have things to tell you concerning that which is to come.

		% --
		\paragraph{Mosiah 3:2} % *MOSIAH-3-2 
		% ---
			\label{MOSIAH-3-2}

			And the things which I shall tell you
			are made known%
				\footnote{%
					\label{MOSIAH-3-2-NOTE-made-known}%
					This is an unusual way of denoting past time. It's clearly an event or series of events that took place in the past, as we see in just a moment when Benjamin says that the angel ``said'' (not ``is saying'' or something like that). And the text refers to this past time with a present form of ``to be'' and then a past participle. It's almost like a more Latin way of doing it, where the past participle really has its past force, so the meaning is something like, ``the things are in a state of having been made known unto me.''
					}
			unto me by an angel from God.
			And he said unto me:
			Awake!
			And I awoke;
			and behold, he stood before me.%
				\footnote{%
					\label{MOSIAH-3-2-NOTE-angel}%
					Reminiscent of JS's own experiences, but also of the experiences of others with heavenly messengers.
					}

		% --
		\paragraph{Mosiah 3:3} % *MOSIAH-3-3 
		% ---
			\label{MOSIAH-3-3}

			And he said unto me:
			Awake%
				\footnote{%
					\label{MOSIAH-3-3-NOTE-awake}%
					Benjamin is told to ``awake'' again after he's already awake. He needed to wake up physically, and then wake up mentally or spiritually in order to pay attention to what the angel was saying. At many days and at many times, we are physically awake, but spiritually and mentally asleep.
					}
			and hear the words which I shall tell thee;
			for behold, I am come to declare unto thee glad tidings of great joy.

		% --
		\paragraph{Mosiah 3:4} % *MOSIAH-3-4 
		% ---
			\label{MOSIAH-3-4}

			For the Lord hath heard thy prayers and hath judged of thy righteousness%
				\footnote{%
					\label{MOSIAH-3-4-NOTE-judged}%
					Very important point about prayer and receiving answers here---the angel specifically mentions that the Lord heard Benjamin's prayers, then \textit{judged} his level of righteousness, and only after (apparently) making a satisfactory judgment did he send the angel to share information with him. Benjamin's righteousness was a factor in receiving an answer to his prayer. This seems like it ought to be a fairly obvious point but for some reason it's not. If we want greater answers to prayers, or maybe answers in the first place, then we need to enhance our personal righteousness.
					}
			and hath sent me to declare unto thee,
			that thou mayest rejoice,
			and that thou mayest declare unto thy people,
			that they may also be filled with joy.

		% --
		\paragraph{Mosiah 3:5} % *MOSIAH-3-5 
		% ---
			\label{MOSIAH-3-5}

			For behold, the time cometh and is not far distant
			that with power the Lord Omnipotent who reigneth,
			which was and is from all eternity to all eternity,
			shall come down from heaven among the children of men
			and shall dwell in a tabernacle of clay
			and shall go forth amongst men, working mighty miracles,
			such as healing the sick,
			raising the dead,
			causing the lame to walk,
			the blind to receive their sight,
			and the deaf to hear,
			and curing all manner of diseases.

		% --
		\paragraph{Mosiah 3:6} % *MOSIAH-3-6 
		% ---
			\label{MOSIAH-3-6}

			And he shall cast out devils,
			or%
				\footnote{%
					\label{MOSIAH-3-6-NOTE-or}%
					Clearly an explanatory use of ``or.'' See \hyperref[OR-COMMENTARY-explanatory]{here}.
					}
			the evil spirits which dwelleth in the hearts of the children of men.

		% --
		\paragraph{Mosiah 3:7} % *MOSIAH-3-7 
		% ---
			\label{MOSIAH-3-7}

			And lo, he shall suffer temptations
			and pain of body, hunger, thirst, and fatigue,
			even more than man can suffer except it be unto death;%
			\footnote{%
				\label{MOSIAH-3-7-NOTE-except}%
				See \hyperref[EXCEPT-1828-participle]{here} for this meaning of ``except.'' The idea is that a normal person could not suffer all these things, could not experience all these things, unless they died, or without dying. What Christ experienced would have caused a normal person to die.
				} 
			for behold, blood cometh from every pore,
			so great shall be his anguish%
				\footnote{%
					\label{MOSIAH-3-7-NOTE-anguish}%
					This word ``anguish'' appears 11 times in the BOM. Note its occurrence above, in \hyperref[MOSIAH-2-38]{Mosiah 2:38}, in the context of what we will feel if we don't repent. It's part of the self-inflicted psychological punishment.
					}
			for the wickedness
			and the abominations of his people.%
				\footnote{%
					% \label{MOSIAH-3-7-NOTE-pore}%
					President Openshaw used to refer to this scripture as, first of all, evidence that Christ did in fact bleed from every pore, but also as an explanation of why that happened---``so great'' was the anguish he felt for what his people were doing. I have felt great emotional pain at times; probably the strongest was when I once talked to Mirzhi near the end of my mission, and found out that she was unhappy and going through some difficult times. I cried for almost a day straight and it was like actual physical pain, like a pressure on my chest. Maybe this is in the direction of what Christ was experiencing. At least, through that experience I can discern the connection between emotional/spiritual and physical pain.
					}

		% --
		\paragraph{Mosiah 3:8} % *MOSIAH-3-8 
		% ---
			\label{MOSIAH-3-8}

			And he shall be called Jesus Christ the Son of God,
			the Father of heaven and of earth,
			the Creator of all things from the beginning;
			and his mother shall be called Mary.

		% --
		\paragraph{Mosiah 3:9} % *MOSIAH-3-9 
		% ---
			\label{MOSIAH-3-9}

			And lo, he cometh unto his own
			that salvation might come unto the children of men,
			even through faith on his name.
			And even after all this,
			they shall consider him as a man and say that he hath a devil,
			and shall \hyperref[SCOURGE-1828-transverb]{scourge} him%
				\footnote{%
					\label{MOSIAH-3-9-NOTE-scourge}%
					There are at least two meanings in play with ``scourge'' here, and maybe three, based on the 1828 dictionary. See \hyperref[SCOURGE-1828-transverb]{here} for the meanings. This use of the term by Benjamin is has much in common with other important verses about the atonement and Christ's role, including \hyperref[ISAIAH-53-5]{Isaiah 53:5} (see also \hyperref[MOSIAH-14-5]{Mosiah 14:5}), \hyperref[1NEPHI-19-9]{1 Nephi 19:9}, and \hyperref[2NEPHI-6-9]{2 Nephi 6:9}, at least. 
					}
			and shall crucify him.

		% --
		\paragraph{Mosiah 3:10} % *MOSIAH-3-10 
		% ---
			\label{MOSIAH-3-10}

			And he shall rise the third day from the dead;
			and behold, he standeth to judge the world.
			And behold, all these things are done
			that a righteous judgment might come upon the children of men.

		% --
		\paragraph{Mosiah 3:11} % *MOSIAH-3-11 
		% ---
			\label{MOSIAH-3-11}

			For behold, and also his blood atoneth for the sins of those
			who have fallen by the transgression of Adam
			who hath died not knowing the will of God concerning them
			or who have ignorantly sinned.%
				\footnote{%
					\label{MOSIAH-3-11-NOTE-categories}%
					Now this is an interesting verse to try to parse. As we can see here, the original text lacks commas that would be useful in helping us grasp exactly what is meant.
					
					I can easily see at least three different readings for this verse---here they are (for all three readings, there could be other groups of people, I think):
					
						\begin{compactdesc}
							\item[Reading \#1] ``His blood atoneth for the sins of people who belong to at least one of the following three groups: (1) those who have fallen by the transgression of Adam, (2) those who have died not knowing the will of God concerning them, or (3) those who have ignorantly sinned.''
							\item[Reading \#2] ``His blood atoneth for the sins of people who belong to at least one of the following two groups: (1) those who have fallen by the transgression of Adam, or (2) those who have died not knowing the will of God concerning them, or in other words, those who have ignorantly sinned'' (see the \hyperref[OR-COMMENTARY-explanatory]{explanatory or}).
							\item[Reading \#3] ``His blood atoneth for the sins of people who have fallen by the transgression of Adam---that is, those who have died not knowing the will of God concerning them, or in other words, those who have ignorantly sinned.''
						\end{compactdesc}
						
					The commas added in the current edition of the BOM suggest reading \#1 or \#2. I think we can rule out reading \#3 on the grounds that sinning ignorantly is not the same as falling by the transgression of Adam, and reading \#3 says that they are basically equivalent.
					
					So we're deciding between readings \#1 and \#2. If reading \#1 is right, then not knowing the will of God concerning you would be a different matter from sinning ignorantly. Knowing or not knowing the will of God concerning you is actually an important property from the perspective of the BOM; see, for example, \hyperref[ALMA-32-16]{Alma 32:16--19}. Being able to equate not knowing the will of God and sinning ignorantly would be helpful in understanding both properties, and from the outset they seem similar in meaning.
					
					I would say I'm inclined to reading \#1, though, because it's more expansive---it gives us more categories to work with, and so more ways of thinking about different people and the situations they are in.
					
					See \hyperref[ATONEMENT-categories]{here} for more discussion of these categories.
					}

		% --
		\paragraph{Mosiah 3:12} % *MOSIAH-3-12 
		% ---
			\label{MOSIAH-3-12}

			But woe woe unto him who knoweth that he rebelleth against God;
			for salvation cometh to none such
			except it be through repentance and faith on the Lord Jesus Christ.

		% --
		\paragraph{Mosiah 3:13} % *MOSIAH-3-13 
		% ---
			\label{MOSIAH-3-13}

			And the Lord God hath sent his holy prophets among all the children of men
			to declare these things to every kindred, nation, and tongue,
			that thereby whosoever should believe that Christ should come,
			the same might receive remission of their sins
			and rejoice with exceeding great joy,
			even as though he had already come among them.%
				\footnote{%
					\label{MOSIAH-3-13-already-come}%
					This last part of the verse tells us something important about the power of the atonement: it was effective in providing a remission of sin even before Christ came to the earth. Somehow there was power in the atonement before the actual event took place. I think this is a really unusual feature of the atonement, and I'm not really sure what it compares to. I can't really think of anything else that works the same way. Maybe part of it is in understanding what a ``remission'' of sin is; if that's a purely or even mostly psychological phenomenon, then it's easier to see how it could happen before the event that the belief or psychological states are based on. Maybe there's the remission, which is more psychological, and the actual satisfying of the demands of the broken law, which (maybe) couldn't happen until the atonement took place in reality. 
					
					The atonement did happen at a particular time and place, and it seems like there were real times before the atonement, yet it still had some kind of effect on people, their sins, or both.
					}

		% --
		\paragraph{Mosiah 3:14} % *MOSIAH-3-14 
		% ---
			\label{MOSIAH-3-14}

			Yet the Lord God saw that his people were a stiffnecked people,
			and he appointed unto them a law, even the law of Moses.%
				\footnote{%
					\label{MOSIAH-3-14-NOTE-law}%
					Is it that he gave them a law because they were stiffnecked? Or is it that he was going to give them a law anyway, and because they were stiffnecked, he decided that it would be the law of Moses? From \hyperref[DC84-23]{DC 84:23--26}, possibly the latter.
					}

		% --
		\paragraph{Mosiah 3:15} % *MOSIAH-3-15 
		% ---
			\label{MOSIAH-3-15}

			And many signs and wonders and types and shadows
			shewed he unto them concerning his coming;
			and also holy prophets spake unto them concerning his coming.
			And yet they hardened their hearts and understood not
			that the law of Moses availeth nothing
			except it were through the atonement of his blood.

		% --
		\paragraph{Mosiah 3:16} % *MOSIAH-3-16 
		% ---
			\label{MOSIAH-3-16}

			And even if it were possible that little children could sin,
			they could not be saved.
			But I say unto you:
			They are blessed;
			for behold, as in Adam or by nature%
				\footnote{%
					% \label{}%
					For ``by nature,'' see \hyperref[EPHESIANS-2-3]{Ephesians 2:3}, 
					}
			they fall,%
				\footnote{%
					% \label{MOSIAH-3-16-NOTE-nature}%
					For many years I have read this part of the verse as saying that we can consider the fall to be something that happens not necessarily through Adam, but instead \textit{by nature}, or as the result of a natural process. This interpretation gives you some room to think about alternative understandings of the fall, possibly not seeing the Adam and Eve story as a series of real events, but instead as a metaphor that is suppoesd to communicate something about the natural state we live in, the difference between us and God, and what will be required of us to make up that difference and eventually gain the kind of life he has.
					
					The meaning of ``or'' is extremely important in this passage as well. Is it disjunctive, or explanatory? To me this one is not as clear-cut as some others. If it's got a disjunctive meaning then I don't think it has exclusive force; if it's got explanatory meaning then we learn something very important about the fall which may help support the interpretation I gave above.
					
					See also \hyperref[DC29-43]{DC 29:43--44},
					}
			even so the blood of Christ atoneth for their sins.

		% --
		\paragraph{Mosiah 3:17} % *MOSIAH-3-17 
		% ---
			\label{MOSIAH-3-17}

			And moreover I say unto you
			that there shall be no other name given nor no other way nor means
			whereby salvation can come unto the children of men,
			only in and through the name of Christ the Lord Omnipotent.

		% --
		\paragraph{Mosiah 3:18} % *MOSIAH-3-18 
		% ---
			\label{MOSIAH-3-18}

			For behold, he judgeth and his judgment is just.
			And the infant perisheth not that dieth in his infancy,
			but men drinketh damnation%
				\footnote{%
					% \label{}%
					The ``drink damnation'' construction comes from \hyperref[1CORINTHIANS-11-29]{1 Corinthians 11:29}.
					}
			to their own souls%
				\footnote{%
					\label{MOSIAH-3-18-NOTE-damnation}%
					It's interesting that we damn ourselves (``drink damnation to our own souls'') unless we humble ourselves and become as little children. One of the things little children to, all the time and in basically every respect, is \textit{learn} and \textit{grow}. All of their experiences teach them new and interesting things, and they grow physically, mentally, spiritually, emotionally, and in other ways. We as adults can adopt this growth mindset, which requires us to be humble and admit what we don't know, and then move forward to make changes based on what we learn. (Make sure to also note the connection here to \hyperref[1CORINTHIANS-11-29]{1 Corinthians 11:29}. See also the notes on damnation in \hyperref[DAMNATION-GREEK]{Greek}.
					} 
			except they humble themselves and become as little children%
				\footnote{%
					\label{MOSIAH-3-18-NOTE-children}%
					Mosiah 3:19 is the famous ``become as a child'' verse, but I hadn't ever noticed that the three verses before it, 3:16--18, all talk about ``children'' in some way, as though they're ramping up to the important point of verse 19 (and 18, and 21 mentions children as well). Here is what we read:
						\begin{compactdesc}
							\item[Mosiah 3:16] Little children cannot sin; the blood of Christ atones for their sins.
							\item[Mosiah 3:17] No other way to salvation than through Christ for the children of men.
							\item[Mosiah 3:18] Infants don't perish if they die; people damn themselves unless they're humble and become like little children.
							\item[Mosiah 3:21] Only little children are blameless before God---and we can become like that through faith and repentance.
						\end{compactdesc}
					One thing we can understand about these verses that mention children, 16--18 and 21, is that if we're supposed to become as a child, then the things said about children in the surrounding verses will then apply to us as well. 
					}
			and believeth that salvation was and is and is to come
			in and through the atoning blood of Christ the Lord Omnipotent.

		% --
		\paragraph{Mosiah 3:19} % *MOSIAH-3-19
		% ---
			\label{MOSIAH-3-19}

				\footnote{%
					% \label{}%
					I wrote a Rainbows essay on what I think is the correct way to interpret what the ``natural man'' is in this verse, and it's along epistemological lines. Read it \href{https://climbingtherainbows.substack.com/p/knowledge-and-the-natural-man}{here}.
					}%
			For the natural man%
				\footnote{%
					% \label{MOSIAH-3-19-NOTE-natural}%
					The biblical anchor verse for the phrase ``natural man'' is \hyperref[1CORINTHIANS-2-14]{1 Corinthians 2:14} (see notes there). The construction in Greek is \gr{ψυχικὸς ἄνθρωπος}. For \gr{ψυχικός}, see \hyperref[PSUCHIKOS]{here}; for other uses of the phrase ``natural man,'' see \hyperref[NATURAL-PHRASES]{here}.
					
					``Natural'' in this verse does not mean ``naturally occuring,'' or ``as we are by ourselves or in nature.'' This is how we use the word ``natural'' and that meaning does not apply here. Meaning \hyperref[PSUCHIKOS]{1a} in BDAG is ``\textit{an unspiritual person}, one who merely functions bodily, without being touched by the Spirit of God.'' This is not the same as just a person in their natural state of being whoever they are. For more, see \hyperref[NATURAL-man]{this study}.
					
					Note that the term ``natural'' points back to 3:16, ``as in Adam or by nature they fall.'' Compare \hyperref[ETHER-3-3]{Ether 3:3}.
					}
			is an \hyperref[ENEMY-1828]{enemy}%
				\footnote{%
					% \label{MOSIAH-3-19-NOTE-enemy}%
					I now read this ``enemy to God'' bit as, the natural man makes himself an enemy to God, but it doesn't say that God is \textit{his} enemy. I don't think God sees his relationship with us in ``enemy'' terms. See \hyperref[ROMANS-8-6]{Romans 8:6--7} and \hyperref[ROMANS-8-27]{27}, which are all the occurrences of \gr{φρόνημα}.
					}
			to God
			and has been from the fall of Adam and will be forever and ever
			but if he yieldeth to the enticings%
				\footnote{%
					% \label{}%
					For the idea that the Holy Spirit can ``entice'' us, see \hyperref[JOHN-12-32]{John 12:32} and ``draw,'' along with meaning \#2 of the Greek word \gr{ἕλκω}, which is used there. It is also described as ``striving'' with us (\hyperref[1NEPHI-7-14]{1 Nephi 7:14}). See also \hyperref[2CORINTHIANS-2-14]{2 Corinthians 2:14}, \hyperref[GALATIANS-5-17]{Galatians 5:17}, \hyperref[2NEPHI-2-16]{2 Nephi 2:16}, \hyperref[2NEPHI-2-28]{2 Nephi 2:28--29}, \hyperlink{EMS-1830-JS-CIR-EARLY-1830-d}{an unpublished 1830 revelation}, 
					}
			of the Holy Spirit
			and putteth off%
				\footnote{%
					% \label{}%
					Notice how this ``put off'' language also appears in Alma's description of Korihor at \hyperref[ALMA-30-42]{Alma 30:42}. You can put off either one, and we're being enticed by both.
					}
			the natural man
			and becometh a saint through the atonement of Christ the Lord
			and
				\footnote{%
					% \label{}%
					I gave a talk on this verse and some of these ideas on \hyperref[TALKS-22-01-2023]{22 January 2023}.
					}%
			becometh as a
				\footnote{%
					% \label{}%
					See \hyperref[NEWBIRTH]{New Birth} but also \hyperref[MATTHEW-18-3]{Matthew 18:3--6}, \hyperref[MARK-10-14]{Mark 10:14--15},
					}%
			child,%
				\footnote{%
					% \label{MOSIAH-3-19-NOTE-child}%
					Normally we read this passage in such a way that the list of qualities following ``child'' are taken as explanations or elaborations of what is meant by ``child.'' But it's also possible to read ``as a child'' as just one among many things which we must ``become'' in order to not be an enemy to God. On this reading becoming as a child would just be one item in the list, and so would becoming submissive, becoming meek, and so on. That way you don't need to take the qualities as explaining what it means to become like a child, because they clearly don't do a very good job of that, I think. Note the connection to \hyperref[ALMA-13-28]{Alma 13:28}, which gives again almost this exact list of qualities, but does not mention the bit about becoming ``as a child,'' which in my mind strengthens the view that the list of qualities isn't meant epexegetically.
					}
			submissive, meek, humble, patient, full of love,%
				\footnote{%
					% \label{}%
					This list of qualities is similar to that found in \hyperref[ALMA-13-28]{Alma 13:28}; Alma seems very clearly to be drawing from his knowledge of what Benjamin said on this occasion.
					}
			willing to submit to all things which the Lord seeth fit to inflict upon him,
			even as a child doth submit%
				\footnote{%
					% \label{}%
					This idea of ``submitting'' to someone is common in the NT, and the primary word is \hyperref[HUPOTASSO]{\grl{ὑποτάσσω}} (I haven't checked all the uses of ``submit,'' but the ones I've seen are \gr{ὑποτάσσω}). The meaning of that word which most sticks out to me is \#1.b.$\beta$, ``\textit{subject oneself, be subjected or subordinated}...of submission involving recognition of an ordered structure, with dative of the entity to whom/which appropriate respect is shown.'' Christ is going to subdue, or cause to submit, all things to himself anyway (\hyperref[DC19-2]{DC 19:2}); he's just inviting us to make this submission to himself voluntary, so that we can experience the benefits of it, instead of the frustrations and fears.
					}
			to his father.%
				\footnote{%
					% \label{MOSIAH-3-19-NOTE-verse}%
					I've always more or less recognized the importance of this verse, but in the past year or so (I'm writing this note Christmas morning 2020) I've really come to appreciate it in a new light. The main impetus for my new thoughts on it came from MRP and from Pook's essay on the fountain of youth. See \hyperref[NEWBIRTH-child]{here} in my notes on the \hyperref[NEWBIRTH]{new birth}.
					}
					
			\begin{framed}
			\scriptsize
				\begin{compactdesc}
					\item[PM] for the natural man is an enemy to God \& has been from the fall of Adam \& will be forever \& ever but if he yieldeth to the enticeings of the holy spirit \& puteth off the natural man \& he becometh a Saint through the atonement of Christ the Lord \& becometh as a child submissive meek humble patient full of love willing to submit to all things which the Lord seeth fit to inflict upon him even as a child doth submit to his father
					\item[1830] for the natural man is an enemy to God, and has been, from the fall of Adam, and will be, forever and ever; but if he yieldeth to the enticings of the Holy Spirit, and putteth off the natural man, and becometh a saint, through the atonement of Christ, the Lord, and becometh as a child, submissive, meek, humble, patient, full of love, willing to submit to all things which the Lord seeth fit to inflict upon him, even as a child doth submit to his father.
					\item[1837] for the natural man is an enemy to God, and has been, from the fall of Adam, and will be, forever and ever; but if he yields to the enticings of the Holy Spirit, and putteth off the natural man, and becometh a saint, through the atonement of Christ, the Lord, and becometh as a child, submissive, meek, humble, patient, full of love, willing to submit to all things which the Lord seeth fit to inflict upon him, even as a child doth submit to his father.
					\item[1840] for the natural man is an enemy to God, and has been, from the fall of Adam, and will be, forever and ever; but if he yields to the enticings of the Holy Spirit, and putteth off the natural man, and becometh a saint, through the atonement of Christ, the Lord, and becometh as a child, submissive, meek, humble, patient, full of love, willing to submit to all things which the Lord seeth fit to inflict upon him, even as a child doth submit to his father.
				\end{compactdesc}
			\end{framed}

		% --
		\paragraph{Mosiah 3:20} % *MOSIAH-3-20 
		% ---
			\label{MOSIAH-3-20}

			And moreover I say unto you
			that the time shall come when the knowledge of a Savior shall spread
			throughout every nation, kindred, tongue, and people.

		% --
		\paragraph{Mosiah 3:21} % *MOSIAH-3-21 
		% ---
			\label{MOSIAH-3-21}

			And behold, when that time cometh,
			none shall be found blameless before God,
			except it be little children,
			only through repentance and faith on the name of the Lord God Omnipotent.

		% --
		\paragraph{Mosiah 3:22} % *MOSIAH-3-22 
		% ---
			\label{MOSIAH-3-22}

			And even at this time when thou shalt have taught thy people the things
			which the Lord thy God hath commanded thee,
			even then are they found no more blameless in the sight of God,
			only according to the words which I have spoken unto thee.

		% --
		\paragraph{Mosiah 3:23} % *MOSIAH-3-23 
		% ---
			\label{MOSIAH-3-23}

			And now I have spoken the words
			which the Lord God hath commanded me.

		% --
		\paragraph{Mosiah 3:24} % *MOSIAH-3-24 
		% ---
			\label{MOSIAH-3-24}

			And thus saith the Lord:
			They shall stand as a bright testimony against this people at the judgment day;
			whereof they shall be judged every man according to his works,
			whether they be good or whether they be evil.

		% --
		\paragraph{Mosiah 3:25} % *MOSIAH-3-25 
		% ---
			\label{MOSIAH-3-25}

			And if they be evil,
			they are \hyperref[CONSIGNATION-1828]{consigned} to an awful view of their own guilt and abominations,
			which doth cause them to shrink from the presence of the Lord
			into a state of misery and endless torment,
			from whence they can no more return;
			therefore they have drunk damnation to their own souls.

		% --
		\paragraph{Mosiah 3:26} % *MOSIAH-3-26 
		% ---
			\label{MOSIAH-3-26}

			Therefore they have drunk out of the cup of the wrath of God,
			which justice could no more deny unto them than it could deny
			that Adam should fall because of his partaking of the forbidden fruit;
			therefore mercy could have claim on them no more forever.

		% --
		\paragraph{Mosiah 3:27} % *MOSIAH-3-27 
		% ---
			\label{MOSIAH-3-27}

			And their torment is as a lake of fire and brimstone
			whose flames are unquenchable
			and whose smoke ascendeth up forever and ever.
			Thus hath the Lord commanded me.
			Amen.
			
	% -----
	\subsubsection{Mosiah 4} % *MOSIAH-4
	% -----
		\label{MOSIAH-4}
		
		% --
		\paragraph{Mosiah 4:1} % *MOSIAH-4-1 
		% ---
			\label{MOSIAH-4-1}

			And now it came to pass that
			when king Benjamin had made an end of speaking the words
			which had been delivered unto him by the angel of the Lord
			that he cast his eyes round about on the multitude,
			and behold, they had fell to the earth,
			for the fear of the Lord had come upon them.

		% --
		\paragraph{Mosiah 4:2} % *MOSIAH-4-2 
		% ---
			\label{MOSIAH-4-2}

			And they had viewed themselves in their own carnal state,
			even less than the dust of the earth,
			and they all cried aloud with one voice, saying:
			O have mercy and apply the atoning blood of Christ
			that we may receive forgiveness of our sins
			and our hearts may be purified,
			for we believe in Jesus Christ the Son of God,
			who created heaven and earth and all things,
			who shall come down among the children of men.

		% --
		\paragraph{Mosiah 4:3} % *MOSIAH-4-3 
		% ---
			\label{MOSIAH-4-3}

			And it came to pass that after they had spoken these words,
			the Spirit of the Lord came upon them,
			and they were filled with joy,
			having received a remission of their sins
			and having peace of conscience because of the exceeding faith
			which they had in Jesus Christ, which should come,
			according to the words which king Benjamin had spoken unto them.%
				\footnote{%
					\label{MOSIAH-4-3-NOTE-newbirth}%
					An excellent example of the \hyperref[NEWBIRTH]{new birth}.
					}

		% --
		\paragraph{Mosiah 4:4} % *MOSIAH-4-4 
		% ---
			\label{MOSIAH-4-4}

			And king Benjamin again opened his mouth
			and began to speak unto them, saying:
			My friends and my brethren, my kindred and my people,
			I would again call your attention,
			that ye may hear and understand the remainder of my words
			which I shall speak unto you.

		% --
		\paragraph{Mosiah 4:5} % *MOSIAH-4-5 
		% ---
			\label{MOSIAH-4-5}

			For behold that if the knowledge of the goodness of God
			at this time hath awakened you
			to a sense of your nothingness and your worthless and fallen state---%
				\footnote{%
					\label{MOSIAH-4-5-NOTE-nothingness}%
					I take this realization more as an intellectual apprehension and not an emotional reaction. It's more like just the awareness of the fact of the difference between us and God, rather than something intended to make us feel depressed or bad about ourselves. God went through what we are experiencing now, and became holy, so we can do it too, and he is there to help us. There's no reason to feel sad or frustrated or miserable for the difference between us and God. We are living out a process which as repeated itself many times, which will repeat itself again after we are done, and which seemingly won't ever end, so we can't blame ourselves for the fact that we currently find ourselves at a particular point in the process.
					}

		% --
		\paragraph{Mosiah 4:6} % *MOSIAH-4-6 
		% ---
			\label{MOSIAH-4-6}

			I say unto you that
			if ye have come to a knowledge of the goodness of God
			and his matchless power and his wisdom and his patience
			and his long-suffering towards the children of men,
			and also the atonement which hath been prepared
			from the foundation of the world,
			that thereby salvation might come to him
			that should put his trust in the Lord
			and should be diligent in keeping his commandments
			and continue in the faith, even unto the end of his life
			---I mean the life of the mortal body---

		% --
		\paragraph{Mosiah 4:7} % *MOSIAH-4-7 
		% ---
			\label{MOSIAH-4-7}

			I say that this is the man that receiveth salvation
			through the atonement which was prepared from the foundation of the world
			for all mankind which ever was, ever since the fall of Adam,
			or which is or which ever shall be, even unto the end of the world.

		% --
		\paragraph{Mosiah 4:8} % *MOSIAH-4-8 
		% ---
			\label{MOSIAH-4-8}

			And this is the means whereby salvation cometh.
			And there is none other salvation save this which hath been spoken of;
			neither is there any conditions whereby man can be saved
			except the conditions%
				\footnote{%
					\label{MOSIAH-4-8-NOTE-conditions}%
					Reminds me of the ``conditions'' mentioned in \hyperref[DC88-38]{DC 88:38--39}; see \hyperref[CONDITION]{Condition}.
					}
			which I have told you.

		% --
		\paragraph{Mosiah 4:9} % *MOSIAH-4-9 
		% ---
			\label{MOSIAH-4-9}

			Believe in God!
			Believe that he is and that he created all things, both in heaven and in earth;
			believe that he hath all wisdom and all power, both in heaven and in earth;
			believe that man doth not comprehend all the things
			which the Lord can comprehend.

		% --
		\paragraph{Mosiah 4:10} % *MOSIAH-4-10 
		% ---
			\label{MOSIAH-4-10}

			And again, believe that ye must repent of your sins and forsake them
			and humble yourselves before God and ask in sincerity of heart
			that he would forgive you.
			And now if you believe all these things,
			see that ye do them.

		% --
		\paragraph{Mosiah 4:11} % *MOSIAH-4-11 
		% ---
			\label{MOSIAH-4-11}

			And again I say unto you, as I have said before,
			that as ye have come to the knowledge of the glory of God
			or%
				\footnote{%
					\label{MOSIAH-4-11-NOTE-or}%
					Explanatory or disjunctive? Very interesting. See \hyperref[OR-COMMENTARY-explanatory]{here}.
					}
			if ye have known of his goodness and have tasted of his love
			and have received a remission of your sins,
			which causeth such exceeding great joy in your souls,
			even so I would that ye should remember and always retain in remembrance
			the greatness of God and your own nothingness
			and his goodness and long-suffering towards you unworthy creatures,
			and humble yourselves even in the depths of humility,
			calling on the name of the Lord daily
			and standing steadfastly in the faith of that which is to come,
			which was spoken by the mouth of the angel.

		% --
		\paragraph{Mosiah 4:12} % *MOSIAH-4-12 
		% ---
			\label{MOSIAH-4-12}

			And behold, I say unto you that if ye do this,
			ye shall always rejoice and be filled with the love of God
			and always retain a remission of your sins;%
				\footnote{%
					% \label{}%
					I'm wondering if we can identify this ``always retain a remission'' with \hyperref[3NEPHI-27-16]{3 Nephi 27:16}, where Christ says that ``if he endureth to the end, behold, him will I hold guiltless before my Father at that day when I shall stand to judge the world.'' So you could see the things here in this passage of Benjamin's sermon as being about enduring to the end, or as involving things we need to do in order to endure.
					}
			and ye shall grow in the knowledge of the glory of him that created you,
			or%
				\footnote{%
					\label{MOSIAH-4-12-NOTE-or}%
					This is an even more interesting and important use of ``or'' which may be explanatory. See \hyperref[OR-COMMENTARY-explanatory]{the list}.
					}
			in the knowledge of that which is just and true.%
				\footnote{%
					\label{MOSIAH-4-12-NOTE-chiasmus}%
					A long time ago I noticed that verses 11 and 12 here contain a chiasmus, that form of apparently Hebrew poetry (?) with repeating elements in a certain pattern. I'll just give the matching pairs in a list here, starting at the front of verse 11 and the end of verse 12:
						\begin{compactitem}
							\item knowledge of the glory of God | knowledge of the glory of him that created you
							\item tasted of his love | filled with the love of God
							\item received a remission of your sins | retain a remission of your sins
							\item exceeding great joy in your souls | ye shall always rejoice
							\item I would that ye should | I say unto you
							\item remember and always retain in remembrance the greatness of God and your own nothingness and his goodness and long-suffering towards you unworthy creatures, and humble yourselves even in the depths of humility, calling on the name of the Lord daily and standing steadfastly in the faith of that which is to come
						\end{compactitem}
					The last item, the instructions about what actually we should do, does not have a parallel, being right in the center of the chiasmus. That is the emphasis of this passage---that is what Benjamin here is trying to get us to \textit{do}. It's a list of actions that will help us retain a remission of our sins. That goal is further elaborated on in the latter parts of the parallels. Instead of just experiencing joy once, we will \textit{always} rejoice; instead of just tasting of God's love, we will \textit{always be filled} with it; instead of receiving a remission of sin once, we will \textit{always retain} that remission; instead of just coming to that single bit of knowledge about the glory of God, we will \textit{grow} in knowledge. The latter parts of the parallels are all about making permanent the experiences and changes mentioned in the former parts, and about growing in those experiences and changes. It isn't just a one-time thing, it can become a permanent part of our life. The progression we go through is underscored by the form of language, and the language also emphasizes the cyclical nature of all this.
					}

		% --
		\paragraph{Mosiah 4:13} % *MOSIAH-4-13 
		% ---
			\label{MOSIAH-4-13}

			And ye will not have a mind to injure one another but to live peaceably,
			and to render to every man according to that which is his due.

		% --
		\paragraph{Mosiah 4:14} % *MOSIAH-4-14 
		% ---
			\label{MOSIAH-4-14}

			And ye will not suffer your children that they go hungry or naked,
			neither will you suffer that they transgress the laws of God
			and fight and quarrel one with another
			and serve the devil, which is the master of sin,
			or which is the evil spirit, which hath been spoken of by our fathers,
			he being an enemy to all righteousness.

		% --
		\paragraph{Mosiah 4:15} % *MOSIAH-4-15 
		% ---
			\label{MOSIAH-4-15}

			But ye will teach them to walk in the ways of truth and soberness;
			ye will teach them to love one another and to serve one another.

		% --
		\paragraph{Mosiah 4:16} % *MOSIAH-4-16 
		% ---
			\label{MOSIAH-4-16}

			And also, ye yourselves will succor those that stand in need of your succor;
			ye will administer of your substance unto him that standeth in need,
			and ye will not suffer that the beggar putteth up his petition to you in vain
			and turn him out to perish.

		% --
		\paragraph{Mosiah 4:17} % *MOSIAH-4-17 
		% ---
			\label{MOSIAH-4-17}

			Perhaps thou shalt say:
			The man hath brought upon himself his misery;
			therefore I will stay my hand and will not give unto him of my food,
			nor impart unto him of my substance that he may not suffer,
			for his punishments are just.

		% --
		\paragraph{Mosiah 4:18} % *MOSIAH-4-18 
		% ---
			\label{MOSIAH-4-18}

			But I say unto you, O man:
			Whosoever doeth this, the same hath great cause to repent;
			and except he repenteth of that which he hath done,
			he perisheth forever and hath no interest%
				\footnote{%
					\label{MOSIAH-4-18-NOTE-interest}%
					This is the only use of any form of ``interest'' in the BOM. The word has several definitions in the 1828 dictionary, making for some very fascinating readings of this verse; see \hyperref[INTEREST-1828]{here}.
					}
			in the kingdom of God.

		% --
		\paragraph{Mosiah 4:19} % *MOSIAH-4-19 
		% ---
			\label{MOSIAH-4-19}

			For behold, are we not all beggars?
			Do we not all depend upon the same Being, even God,
			for all the substance which we have,
			for both food and raiment and for gold and for silver
			and for all the riches which we have, of every kind?

		% --
		\paragraph{Mosiah 4:20} % *MOSIAH-4-20 
		% ---
			\label{MOSIAH-4-20}

			And behold, even at this time ye have been calling on his name
			and begging for a remission of your sins.
			And hath he suffered that ye have begged in vain?
			Nay, he hath poured out his Spirit upon you
			and hath caused that your hearts should be filled with joy
			and hath caused that your mouths should be stopped,
			that ye could not find utterance,
			so exceeding great was your joy.

		% --
		\paragraph{Mosiah 4:21} % *MOSIAH-4-21 
		% ---
			\label{MOSIAH-4-21}

			And now if God, who hath created you,
			on whom you are dependent for your lives and for all that ye have and are,
			doth grant unto you whatsoever ye ask that is right,
			in faith believing that ye shall receive,
			O then how had ye ought to impart of the substance that ye have one to another!

		% --
		\paragraph{Mosiah 4:22} % *MOSIAH-4-22 
		% ---
			\label{MOSIAH-4-22}

			And if ye judge the man who putteth up his petition to you
			for your substance that he perish not,
			and condemn him,
			how much more just will be your condemnation
			for withholding your substance, which doth not belong to you but to God,
			to whom also your life belongeth,
			and yet ye put up no petition or repenteth not of the thing which thou hast done.

		% --
		\paragraph{Mosiah 4:23} % *MOSIAH-4-23 
		% ---
			\label{MOSIAH-4-23}

			I say unto you:
			Woe be unto that man,
			for his substance shall perish with him.
			And now I say these things unto those which are rich
			as pertaining to the things of this world.

		% --
		\paragraph{Mosiah 4:24} % *MOSIAH-4-24 
		% ---
			\label{MOSIAH-4-24}

			And again I say unto the poor, ye that have not
			and yet hath sufficient that ye remain from day to day
			---I mean all you that deny the beggar because ye have not---
			I would that ye say in your hearts
			that I give not because I have not,
			but if I had, I would give.

		% --
		\paragraph{Mosiah 4:25} % *MOSIAH-4-25 
		% ---
			\label{MOSIAH-4-25}

			And now if ye say this in your hearts,
			ye remain guiltless;
			otherwise ye are condemned,
			and your condemnation is just,
			for ye covet that which ye have not received.

		% --
		\paragraph{Mosiah 4:26} % *MOSIAH-4-26 
		% ---
			\label{MOSIAH-4-26}

			And now for the sake of these things which I have spoken unto you
			---that is, for the sake of retaining a remission of your sins from day to day,
			that ye may walk guiltless before God---
			I would that ye should impart of your substance to the poor,
			every man according to that which he hath,
			such as feeding the hungry,
			clothing the naked,
			visiting the sick and administering to their relief,
			both spiritually and temporally according to their wants.

		% --
		\paragraph{Mosiah 4:27} % *MOSIAH-4-27 
		% ---
			\label{MOSIAH-4-27}

			And see that all these things are done in wisdom and order,
			for it is not requisite that a man should run faster than what he hath strength.
			And again, it is expedient that he should be diligent,
			that thereby he might win the prize.
			Therefore all things must be done in order.

		% --
		\paragraph{Mosiah 4:28} % *MOSIAH-4-28 
		% ---
			\label{MOSIAH-4-28}

			And I would that ye should remember
			that whosoever among you that borroweth of his neighbor
			should return the thing that he borroweth,
			according as he doth agree;
			or else thou shalt commit sin,
			and perhaps thou shalt cause thy neighbor to commit sin also.

		% --
		\paragraph{Mosiah 4:29} % *MOSIAH-4-29 
		% ---
			\label{MOSIAH-4-29}

			And finally, I cannot tell you all the things whereby ye may commit sin;
			for there are divers ways and means,
			even so many that I cannot number them.

		% --
		\paragraph{Mosiah 4:30} % *MOSIAH-4-30 
		% ---
			\label{MOSIAH-4-30}

			But this much I can tell you,
			that if ye do not watch yourselves
			and your thoughts and your words and your deeds
			and observe to keep the commandments of God
			and continue in the faith of what ye have heard
			concerning the coming of our Lord,
			even unto the end of your lives,
			ye must perish.
			And now O man, remember and perish not!%
				\footnote{%
					\label{MOSIAH-4-30-NOTE-exclamation}%
					For the life of me I cannot understand why the current edition of the BOM weakens the text so much. Note the original version here, and compare to the current edition: ``And now, O man, remember, and perish not.'' It isn't just removing the exclamation point, which is bad enough; the new edition also adds two commas, which significantly slows the sentence down. It really makes the passage much weaker than it was intended to be.
					}
					
	% -----
	\subsubsection{Mosiah 5} % *MOSIAH-4
	% -----
		\label{MOSIAH-5}
					
		% --
		\paragraph{Mosiah 5:1} % *MOSIAH-5-1 
		% ---
			\label{MOSIAH-5-1}

			And now it came to pass that
			when king Benjamin had thus spoken to his people,
			he sent among them, desiring to know of his people
			if they believed the words which he had spoken unto them.

		% --
		\paragraph{Mosiah 5:2} % *MOSIAH-5-2 
		% ---
			\label{MOSIAH-5-2}

			And they all cried with one voice, saying:
			Yea, we believe all the words which thou hast spoken unto us.
			And also we know of their surety and truth
			because of the Spirit of the Lord Omnipotent,
			which hath wrought a mighty change in us or in our hearts,
			that we have no more disposition to do evil but to do good continually.

		% --
		\paragraph{Mosiah 5:3} % *MOSIAH-5-3 
		% ---
			\label{MOSIAH-5-3}

			And we ourselves also through the infinite goodness of God
			and the manifestations of his Spirit
			have great views of that which is to come.
			And were it expedient, we could prophesy of all things.

		% --
		\paragraph{Mosiah 5:4} % *MOSIAH-5-4 
		% ---
			\label{MOSIAH-5-4}

			And it is the faith which we have had on the things
			which our king hath spoken unto us
			and hath brought us to this great knowledge,
			whereby we do rejoice with such exceeding great joy.

		% --
		\paragraph{Mosiah 5:5} % *MOSIAH-5-5 
		% ---
			\label{MOSIAH-5-5}

			And we are willing to enter into a covenant with our God
			to do his will and to be obedient to his commandments
			in all things that he shall command us
			all the remainder of our days,
			that we may not bring upon ourselves a never-ending torment
			as has been spoken by the angel,
			that we may not drink out of the cup of the wrath of God.

		% --
		\paragraph{Mosiah 5:6} % *MOSIAH-5-6 
		% ---
			\label{MOSIAH-5-6}

			And now these are the words which king Benjamin desired of them.
			And therefore he said unto them:
			Ye have spoken the words that I desired,
			and the covenant which ye have made is a righteous covenant.%
				\footnote{%
					\label{MOSIAH-5-6-NOTE-covenant}%
					When do the people make the covenant? In this verse Benjamin says ``the covenant which ye have made,'' and in the next he says the very same thing. But it isn't clear when this actually occurs. In verse 5 the people declare that they are ``willing to enter into a covenant with our God,'' but the text doesn't say anything about a ritual or ordinance or anything else. It could be that this part was just left out for some reason, or it may also be that their expression of willingness \textit{does} mark the making of the covenant, and there isn't anything more to it than that.
					}

		% --
		\paragraph{Mosiah 5:7} % *MOSIAH-5-7 
		% ---
			\label{MOSIAH-5-7}

			And now because of the covenant which ye have made,
			ye shall be called the
				\footnote{%
					% \label{}%
					See \hyperref[NEWBIRTH-children-christ]{this study} for more on the ``children of Christ.'' I believe we are called the ``children of Christ'' because we are reborn through his power, through the priesthood, which is the priesthood ``after the order of the son of God.''
					}%
			children of Christ, his sons and his daughters;
			for behold, this day he hath spiritually begotten you,
			for ye say that your hearts are changed through faith on his name;
			therefore ye are born of him and have become his sons and his daughters.

		% --
		\paragraph{Mosiah 5:8} % *MOSIAH-5-8 
		% ---
			\label{MOSIAH-5-8}

			And under this head ye are made free,
			and there is no other head%
				\footnote{%
					% \label{MOSIAH-5-8-NOTE-head}%
					A few possibilities for what this word means, from the OED:
					\begin{compactitem}
						\item[III.31] ``A chief or principal point or division of a discourse, subject, etc.; each of a set or succession of such points or divisions; (more generally) a point, a category, a topic, a matter.''
						\item[IV.34] ``A person or thing holding the senior or most important position; a chief or leader.''
						\item[IV.34a] ``A person to whom others are subordinate; a chief, a ruler, a leader, a commander.''
					\end{compactitem}
					}
			whereby ye can be made free;
			there is no other name given whereby salvation cometh.
			Therefore I would that ye should take upon you the name of Christ,
			all you that have entered into the covenant with God
			that ye should be obedient unto the end of your lives.

		% --
		\paragraph{Mosiah 5:9} % *MOSIAH-5-9 
		% ---
			\label{MOSIAH-5-9}

			And it shall come to pass that
			whosoever doeth this shall be found at the right hand of God,
			for he shall know the name by which he is called;
			for he shall be called by the name of Christ.

		% --
		\paragraph{Mosiah 5:10} % *MOSIAH-5-10 
		% ---
			\label{MOSIAH-5-10}

			And now it shall come to pass that
			whosoever shall not take upon them the name of Christ
			must be called by some other name;
			therefore he findeth himself on the left hand of God.

		% --
		\paragraph{Mosiah 5:11} % *MOSIAH-5-11 
		% ---
			\label{MOSIAH-5-11}

			And I would that ye should remember also
			that this is the name that I said I should give unto you
			that never should be blotted out except it be through transgression;
			therefore take heed that ye do not transgress,
			that the name be not blotted out of your hearts.

		% --
		\paragraph{Mosiah 5:12} % *MOSIAH-5-12 
		% ---
			\label{MOSIAH-5-12}

			I say unto you:
			I would that ye should remember
			to retain the name written
			always in your hearts,
			that ye are not found on the left hand of God,
			but that ye hear and know the voice by which ye shall be called,
			and also the name by which he shall call you.%
				\footnote{%
					% \label{MOSIAH-5-12-NOTE-voice}%
					According to this verse, we are supposed to hear and know two things: the\textit{voice} that will call us, and the \textit{name} which it will use to call us. If we know the voice we'll be aware that God is speaking, and we'll be paying attention, and if we know the particular name, then we'll be ready when the voice calls us. For what exactly, I'm not sure.
					}

		% --
		\paragraph{Mosiah 5:13} % *MOSIAH-5-13 
		% ---
			\label{MOSIAH-5-13}

			For how knoweth a man the master which he hath not served
			and which is a stranger unto him
			and is far from the thoughts and intents of his heart?

		% --
		\paragraph{Mosiah 5:14} % *MOSIAH-5-14 
		% ---
			\label{MOSIAH-5-14}

			And again, doth a man take an ass which belongeth to his neighbor and keep him?
			I say unto you: Nay.
			He will not even suffer that he shall feed among his flocks,
			but will drive him away and cast him out.
			I say unto you that even so shall it be among you
			if ye know not the name by which ye are called.%
				\footnote{%
					\label{MOSIAH-5-14-NOTE-flocks}%
					Until today (1/17/21) I hadn't ever really thought about this verse. It's an unusual image. Is it that we as individuals are the asses? And if we're owned or sealed by the devil, then Christ won't be able to take us and have us feed with his flock; we'll be driven away and cast out. Is that it? Kind of strange.
					}

		% --
		\paragraph{Mosiah 5:15} % *MOSIAH-5-15 
		% ---
			\label{MOSIAH-5-15}

			Therefore I would that ye should be steadfast and immovable,
			always abounding in good works,
			that Christ the Lord God Omnipotent may seal you his,
			that you may be brought to heaven,
			that ye may have everlasting salvation and eternal life
			through the wisdom and power and justice and mercy of him
			who created all things in heaven and in earth,
			who is God above all.
			Amen.
			
	% -----
	\subsubsection{Mosiah 6} % *MOSIAH-6
	% -----
		\label{MOSIAH-6}
		
		% ---
		\paragraph{Mosiah 6:1} % *MOSIAH-6-1 
		% ---
			\label{MOSIAH-6-1}

			And now king Benjamin thought it was expedient,
			after having finished speaking to the people,
			that he should take the names of all those
			who had entered into a covenant with God to keep his commandments.

		% ---
		\paragraph{Mosiah 6:2} % *MOSIAH-6-2 
		% ---
			\label{MOSIAH-6-2}

			And it came to pass that there was not one soul, except it were little children,
			but what had entered into the covenant
			and had taken upon them the name of Christ.

		% ---
		\paragraph{Mosiah 6:3} % *MOSIAH-6-3 
		% ---
			\label{MOSIAH-6-3}

			And again it came to pass that
			when king Benjamin had made an end of all these things
			and had \hyperref[CONSECRATION]{consecrated} his son Mosiah to be a ruler and a king over his people
			and had given him all the charges concerning the kingdom
			and also had appointed priests to teach the people,
			that thereby they might hear and know the commandments of God,
			and to stir them up in remembrance of the oath which they had made,
			he dismissed the multitude,
			and they returned every one according to their families to their own houses.

		% ---
		\paragraph{Mosiah 6:4} % *MOSIAH-6-4 
		% ---
			\label{MOSIAH-6-4}

			And Mosiah began to reign in his father's stead,
			and he began to reign in the thirtieth year of his age,
			making in the whole about four hundred and seventy six years
			from the time that Lehi left Jerusalem.

		% ---
		\paragraph{Mosiah 6:5} % *MOSIAH-6-5 
		% ---
			\label{MOSIAH-6-5}

			And king Benjamin lived three years and he died.

		% ---
		\paragraph{Mosiah 6:6} % *MOSIAH-6-6 
		% ---
			\label{MOSIAH-6-6}

			And it came to pass that king Mosiah did walk in the ways of the Lord
			and did observe his judgments and his statutes
			and did keep his commandments in all things whatsoever he commanded him.

		% ---
		\paragraph{Mosiah 6:7} % *MOSIAH-6-7 
		% ---
			\label{MOSIAH-6-7}

			And king Mosiah did cause his people that they should till the earth,
			and he also himself did till the earth,
			that thereby he might not become burthensome to his people,
			that he might do according to that which his father had done in all things.
			And there was no contention among all his people for the space of three years.
			
	% -----
	\subsubsection{Mosiah 7} % *MOSIAH-7 
	% -----
		\label{MOSIAH-7}
		
		% ---
		\paragraph{Mosiah 7:1} % *MOSIAH-7-1 
		% ---
			\label{MOSIAH-7-1}

			And now it came to pass that
			after king Mosiah had had continual peace for the space of three years,
			he was desirous to know concerning the people
			which went up to dwell in the land of Lehi-Nephi, or in the city of Lehi-Nephi,
			for his people had heard nothing from them
			from the time they left the land of Zarahemla;
			therefore they wearied him with their teasings.

		% ---
		\paragraph{Mosiah 7:2} % *MOSIAH-7-2 
		% ---
			\label{MOSIAH-7-2}

			And it came to pass that king Mosiah granted
			that sixteen of their strong men might go up to the land of Lehi-Nephi
			to inquire concerning their brethren.

		% ---
		\paragraph{Mosiah 7:3} % *MOSIAH-7-3 
		% ---
			\label{MOSIAH-7-3}

			And it came to pass that on the morrow they started to go up,
			having with them one Ammon,
			he being a strong and mighty man and a descendant of Zarahemla;
			and he was also their leader.

		% ---
		\paragraph{Mosiah 7:4} % *MOSIAH-7-4 
		% ---
			\label{MOSIAH-7-4}

			And now they knew not the course they should travel in the wilderness
			to go up to the land of Lehi-Nephi;
			therefore they wandered many days in the wilderness,
			even forty days did they wander.

		% ---
		\paragraph{Mosiah 7:5} % *MOSIAH-7-5 
		% ---
			\label{MOSIAH-7-5}

			And when they had wandered forty days,
			they came to a hill which is north of the land of Shilom,
			and there they pitched their tents.

		% ---
		\paragraph{Mosiah 7:6} % *MOSIAH-7-6 
		% ---
			\label{MOSIAH-7-6}

			And Ammon took three of his brethren
			---and their names were Amaleki, Helem, and Hem---
			and they went down into the land of Nephi.

		% ---
		\paragraph{Mosiah 7:7} % *MOSIAH-7-7 
		% ---
			\label{MOSIAH-7-7}

			And behold, they met the king of the people
			which was in the land of Nephi and in the land of Shilom;
			and they were surrounded by the king's guard
			and was taken and was bound and was committed to prison.

		% ---
		\paragraph{Mosiah 7:8} % *MOSIAH-7-8 
		% ---
			\label{MOSIAH-7-8}

			And it came to pass when they had been in prison two days,
			they were again brought before the king
			and their bands were loosed;
			and they stood before the king
			and was permitted---or rather commanded---
			that they should answer the questions which he should ask them.

		% ---
		\paragraph{Mosiah 7:9} % *MOSIAH-7-9 
		% ---
			\label{MOSIAH-7-9}

			And he saith unto them:
			Behold, I am Limhi the son of Noah, which was the son of Zeniff,
			which came up out of the land of Zarahemla
			to inherit this land which was the land of their fathers,
			which was made a king by the voice of the people.

		% ---
		\paragraph{Mosiah 7:10} % *MOSIAH-7-10 
		% ---
			\label{MOSIAH-7-10}

			And now I desire to know the cause whereby ye were so bold
			as to come near the walls of the city
			when I myself was with my guards without the gate.

		% ---
		\paragraph{Mosiah 7:11} % *MOSIAH-7-11 
		% ---
			\label{MOSIAH-7-11}

			And now for this cause have I suffered that ye should be preserved,
			that I might inquire of you;
			or else I should have caused that my guards should have put you to death.
			Ye are permitted to speak.

		% ---
		\paragraph{Mosiah 7:12} % *MOSIAH-7-12 
		% ---
			\label{MOSIAH-7-12}

			And now when Ammon saw that he was permitted to speak,
			he went forth and bowed himself before the king;
			and rising again, he said:
			O king, I am very thankful before God this day
			that I am yet alive and am permitted to speak.
			And I will endeavor to speak with boldness;

		% ---
		\paragraph{Mosiah 7:13} % *MOSIAH-7-13 
		% ---
			\label{MOSIAH-7-13}

			for I am assured that if ye had known me,
			ye would not have suffered that I should have wore these bands;
			for I am Ammon and am a descendant of Zarahemla
			and have come up out of the land of Zarahemla
			to inquire concerning our brethren which Zeniff brought up out of that land.

		% ---
		\paragraph{Mosiah 7:14} % *MOSIAH-7-14 
		% ---
			\label{MOSIAH-7-14}

			And now it came to pass that
			after Limhi had heard the words of Ammon,
			he was exceeding glad and said:
			Now I know of a surety
			that my brethren which was in the land of Zarahemla are yet alive.
			And now I will rejoice,
			and on the morrow I will cause that my people shall rejoice also.

		% ---
		\paragraph{Mosiah 7:15} % *MOSIAH-7-15 
		% ---
			\label{MOSIAH-7-15}

			For behold, we are in bondage to the Lamanites
			and are taxed with a tax which is grievious to be borne.
			And now behold, our brethren will deliver us out of our bondage,
			or out of the hands of the Lamanites.
			And we will be their slaves,
			for it is better that we be slaves to the Nephites
			than to pay tribute to the king of the Lamanites.

		% ---
		\paragraph{Mosiah 7:16} % *MOSIAH-7-16 
		% ---
			\label{MOSIAH-7-16}

			And now king Limhi commanded his guards
			that they should no more bind Ammon nor his brethren,
			but caused that they should go to the hill which was north of Shilom
			and bring their brethren into the city,
			that thereby they might eat and drink
			and rest themselves from the labors of their journey,
			for they had suffered many things;
			they had suffered hunger, thirst, and fatigue.

		% ---
		\paragraph{Mosiah 7:17} % *MOSIAH-7-17 
		% ---
			\label{MOSIAH-7-17}

			And now it came to pass on the morrow
			that king Limhi sent a proclamation among all his people,
			that thereby they might gather themselves together to the temple
			to hear the words which he should speak unto them.

		% ---
		\paragraph{Mosiah 7:18} % *MOSIAH-7-18 
		% ---
			\label{MOSIAH-7-18}

			And it came to pass that when they had gathered themselves together
			that he spake unto them in this wise, saying:
			O ye my people, lift up your heads and be comforted;
			for behold, the time is at hand---or is not far distant---
			when we shall no longer be in subjection to our enemies;
			notwithstanding our many strugglings which have been in vain,
			yet I trust there remaineth an effectual struggle to be made.

		% ---
		\paragraph{Mosiah 7:19} % *MOSIAH-7-19 
		% ---
			\label{MOSIAH-7-19}

			Therefore lift up your heads and rejoice and put your trust in God,
			in that God who was the God of Abraham and Isaac and Jacob,
			and also that God who brought the children of Israel out of the land of Egypt
			and caused that they should walk through the Red Sea on dry ground
			and fed them with manna that they might not perish in the wilderness;
			and many more things did he do for them.

		% ---
		\paragraph{Mosiah 7:20} % *MOSIAH-7-20 
		% ---
			\label{MOSIAH-7-20}

			And again, that same God hath brought our fathers out of the land of Jerusalem
			and hath kept and preserved his people even until now.
			And behold, it is because of our iniquities and abominations
			that has brought us into bondage.

		% ---
		\paragraph{Mosiah 7:21} % *MOSIAH-7-21 
		% ---
			\label{MOSIAH-7-21}

			And ye all are witnesses this day
			that Zeniff, who was made king over this people,
			he being overzealous to inherit the land of his fathers,
			therefore being deceived by the cunning and craftiness of king Laman,
			who having entered into a treaty with king Zeniff
			and having yielded up into his hands the possessions of a part of the land,
			or even the city of Lehi-Nephi and the city of Shilom and the land round about---

		% ---
		\paragraph{Mosiah 7:22} % *MOSIAH-7-22 
		% ---
			\label{MOSIAH-7-22}

			and all this he done for the sole purpose
			of bringing this people into subjection or into bondage.
			And behold, we at this time do pay tribute to the king of the Lamanites
			to the amount of one half of our corn and our barley,
			and even all our grain of every kind,
			and one half of the increase of our flocks and our herds;
			and even one half of all we have or possess
			the king of the Lamanites doth exact of us---or our lives.

		% ---
		\paragraph{Mosiah 7:23} % *MOSIAH-7-23 
		% ---
			\label{MOSIAH-7-23}

			And now, is not this grievious to be borne?
			And is not this our afflictions great?
			Now behold, how great reason have we to mourn!

		% ---
		\paragraph{Mosiah 7:24} % *MOSIAH-7-24 
		% ---
			\label{MOSIAH-7-24}

			Yea, I say unto you:
			Great are the reasons which we have to mourn;
			for behold, how many of our brethren have been slain
			and their blood hath been spilt in vain---
			and all because of iniquity!

		% ---
		\paragraph{Mosiah 7:25} % *MOSIAH-7-25 
		% ---
			\label{MOSIAH-7-25}

			For if this people had not fallen into transgression,
			the Lord would not have suffered that this great evil should come upon them.
			But behold, they would not hearken unto his words,
			but there arose contentions among them,
			even so much that they did shed blood among themselves.

		% ---
		\paragraph{Mosiah 7:26} % *MOSIAH-7-26 
		% ---
			\label{MOSIAH-7-26}

			And a prophet of the Lord have they slain---
			yea, a chosen man of God who told them of their wickedness and abominations
			and prophesied of many things which is to come,
			yea, even the coming of Christ.

		% ---
		\paragraph{Mosiah 7:27} % *MOSIAH-7-27 
		% ---
			\label{MOSIAH-7-27}

			And because he saith unto them that Christ was the God the Father of all things
			and saith that he should take upon him the image of man
			and it should be the image after which man was created in the beginning
			---or in other words, he said that man was created after the image of God
			and that God should come down among the children of men
			and take upon him flesh and blood
			and go forth upon the face of the earth---

		% ---
		\paragraph{Mosiah 7:28} % *MOSIAH-7-28 
		% ---
			\label{MOSIAH-7-28}

			and now because he said this, they did put him to death.
			And many more things did they do
			which brought down the wrath of God upon them.
			Therefore who wondereth that they are in bondage
			and that they are smitten with sore afflictions!

		% ---
		\paragraph{Mosiah 7:29} % *MOSIAH-7-29 
		% ---
			\label{MOSIAH-7-29}

			For behold, the Lord hath said:
			I will not succor my people in the day of their transgression,
			but I will hedge up their ways, that they prosper not;
			and their doings shall be as a stumbling block before them.

		% ---
		\paragraph{Mosiah 7:30} % *MOSIAH-7-30 
		% ---
			\label{MOSIAH-7-30}

			And again he saith:
			If my people shall sow filthiness,
			they shall reap the chaff thereof in the whirlwind;
			and the effects thereof is poison.

		% ---
		\paragraph{Mosiah 7:31} % *MOSIAH-7-31 
		% ---
			\label{MOSIAH-7-31}

			And again he saith:
			If my people shall sow filthiness,
			they shall reap the east wind, which bringeth immediate destruction.

		% ---
		\paragraph{Mosiah 7:32} % *MOSIAH-7-32 
		% ---
			\label{MOSIAH-7-32}

			And now behold, the promise of the Lord is fulfilled,
			and ye are smitten and afflicted.

		% ---
		\paragraph{Mosiah 7:33} % *MOSIAH-7-33 
		% ---
			\label{MOSIAH-7-33}

			But if ye will turn to the Lord with full purpose of heart
			and put your trust in him and serve him with all diligence of mind---
			and if ye do this,
			he will, according to his own will and pleasure, deliver you out of bondage.

	% -----
	\subsubsection{Mosiah 8} % *MOSIAH-8 
	% -----
		\label{MOSIAH-8}
		
		% ---
		\paragraph{Mosiah 8:1} % *MOSIAH-8-1 
		% ---
			\label{MOSIAH-8-1}

			And it came to pass that
			after king Limhi had made an end of speaking to his people
			---for he spake many things unto them,
			and only a few of them have I written in this book---
			he told his people all the things concerning their brethren
			which were in the land of Zarahemla.

		% ---
		\paragraph{Mosiah 8:2} % *MOSIAH-8-2 
		% ---
			\label{MOSIAH-8-2}

			And he caused that Ammon should stand up before the multitude
			and rehearse unto them all that had happened unto their brethren
			from the time that Zeniff went up out of the land
			even until the time that he himself came up out of the land.

		% ---
		\paragraph{Mosiah 8:3} % *MOSIAH-8-3 
		% ---
			\label{MOSIAH-8-3}

			And he also rehearsed unto them the last words
			which king Benjamin had taught them,
			and explained them to the people of king Limhi,
			so that they might understand all the words which he spake.

		% ---
		\paragraph{Mosiah 8:4} % *MOSIAH-8-4 
		% ---
			\label{MOSIAH-8-4}

			And it came to pass that after he had done all this
			that king Limhi dismissed the multitude
			and caused that they should return every one unto his own house.

		% ---
		\paragraph{Mosiah 8:5} % *MOSIAH-8-5 
		% ---
			\label{MOSIAH-8-5}

			And it came to pass that he caused that the plates
			which contained the record of his people
			from the time that they left the land of Zarahemla
			should be brought before Ammon, that he might read them.

		% ---
		\paragraph{Mosiah 8:6} % *MOSIAH-8-6 
		% ---
			\label{MOSIAH-8-6}

			Now as soon as Ammon had read the record,
			the king inquired of him to know if he could interpret languages.
			And Ammon told him that he could not.

		% ---
		\paragraph{Mosiah 8:7} % *MOSIAH-8-7 
		% ---
			\label{MOSIAH-8-7}

			And the king saith unto him:
			I being grieved for the afflictions of my people,
			I caused that forty and three of my people
			should take a journey into the wilderness,
			that thereby they might find the land of Zarahemla,
			that we might appeal unto our brethren to deliver us out of bondage.

		% ---
		\paragraph{Mosiah 8:8} % *MOSIAH-8-8 
		% ---
			\label{MOSIAH-8-8}

			And they were lost in the wilderness for the space of many days,
			yet they were diligent and found not the land of Zarahemla
			but returned to this land,
			having traveled in a land among many waters,
			having discovered a land which was covered with bones of men and of beasts etc.,
			and was also covered with ruins of buildings of every kind,
			having discovered a land which had been peopled with a people
			which were as numerous as the hosts of Israel.

		% ---
		\paragraph{Mosiah 8:9} % *MOSIAH-8-9 
		% ---
			\label{MOSIAH-8-9}

			And for a testimony that the things that they have said is true,
			they have brought twenty four plates which are filled with engravings;
			and they are of pure gold.

		% ---
		\paragraph{Mosiah 8:10} % *MOSIAH-8-10 
		% ---
			\label{MOSIAH-8-10}

			And behold also, they have brought breastplates which are large,
			and they are of brass and of copper and are perfectly sound.

		% ---
		\paragraph{Mosiah 8:11} % *MOSIAH-8-11 
		% ---
			\label{MOSIAH-8-11}

			And again, they have brought swords;
			the hilts thereof hath perished
			and the blades thereof were cankered with rust.
			And there is no one in the land that is able to interpret the language
			or the engravings that are on the plates.
			Therefore I said unto thee:
			Canst thou translate?

		% ---
		\paragraph{Mosiah 8:12} % *MOSIAH-8-12 
		% ---
			\label{MOSIAH-8-12}

			And I say unto thee again:
			Knowest thou of any one that can translate?
			For I am desirous that these records should be translated into our language,
			for perhaps they will give us a knowledge of a remnant of the people
			which have been destroyed,
			from whence these records came.
			Or perhaps they will give us a knowledge of this very people
			which hath been destroyed.
			And I am desirous to know the cause of their destruction.

		% ---
		\paragraph{Mosiah 8:13} % *MOSIAH-8-13 
		% ---
			\label{MOSIAH-8-13}

			Now Ammon saith unto him:
			I can assuredly tell thee, O king, of a man that can translate the records;
			for he hath wherewith that he can look
			and translate all records that are of ancient date,
			and it is a gift from God.
			And the things are called interpreters,
			and no man can look in them except he be commanded,
			lest he should look for that he had not ought and he should perish.
			And whosoever is commanded to look in them, the same is called seer.

		% ---
		\paragraph{Mosiah 8:14} % *MOSIAH-8-14 
		% ---
			\label{MOSIAH-8-14}

			And behold, the king of the people which is in the land of Zarahemla
			is the man that is commanded to do these things
			and which hath this high gift from God.

		% ---
		\paragraph{Mosiah 8:15} % *MOSIAH-8-15 
		% ---
			\label{MOSIAH-8-15}

			And the king saith that a seer is greater than a prophet.

		% ---
		\paragraph{Mosiah 8:16} % *MOSIAH-8-16 
		% ---
			\label{MOSIAH-8-16}

			And Ammon saith that a seer is a revelator and a prophet also.
			And a gift which is greater can no man have
			except he should possess the power of God, which no man can;
			yet a man may have great power given him from God.

		% ---
		\paragraph{Mosiah 8:17} % *MOSIAH-8-17 
		% ---
			\label{MOSIAH-8-17}

			But a seer can know of things which has passed,
			and also of things which is to come;
			and by them shall all things be revealed
			---or rather shall secret things be made manifest---
			and hidden things shall come to light,
			and things which is not known shall be made known by them,
			and also things shall be made known by them
			which otherwise could not be known.

		% ---
		\paragraph{Mosiah 8:18} % *MOSIAH-8-18 
		% ---
			\label{MOSIAH-8-18}

			Thus God hath provided a means
			that man through faith might work mighty miracles;
			therefore he becometh a great benefit to his fellow beings.

		% ---
		\paragraph{Mosiah 8:19} % *MOSIAH-8-19 
		% ---
			\label{MOSIAH-8-19}

			And now when Ammon had made an end of speaking these words,
			the king rejoiced exceedingly and gave thanks to God, saying:
			Doubtless a great mystery is contained within these plates;
			and these interpreters was doubtless prepared
			for the purpose of unfolding all such mysteries to the children of men.

		% ---
		\paragraph{Mosiah 8:20} % *MOSIAH-8-20 
		% ---
			\label{MOSIAH-8-20}

			O how marvelous are the works of the Lord!
			And how long doth he suffer with his people!
			Yea, and how blind and impenetrable
			are the understandings of the children of men,
			for they will not seek wisdom,
			neither do they desire that she should rule over them.

		% ---
		\paragraph{Mosiah 8:21} % *MOSIAH-8-21 
		% ---
			\label{MOSIAH-8-21}

			Yea, they are as a wild flock which fleeth from the shepherd and scattereth,
			and are driven and are devoured by the beasts of the forest.

	% -----
	\subsubsection{Mosiah 9} % *MOSIAH-9 
	% -----
		\label{MOSIAH-9}
		
		The Record of Zeniff

		An account of his people from the time they left the land of Zarahemla
		until the time that they were delivered out of the hands of the Lamanites.
		
		% ---
		\paragraph{Mosiah 9:1} % *MOSIAH-9-1 
		% ---
			\label{MOSIAH-9-1}

			I Zeniff having been taught in all the language of the Nephites
			and having had a knowledge of the land of Nephi,
			or of the land of our fathers' first inheritance,
			and I having been sent as a spy among the Lamanites
			that I might spy out their forces,
			that our army might come upon them and destroy them---
			but when I saw that which was good among them,
			I was desirous that they should not be destroyed.

		% ---
		\paragraph{Mosiah 9:2} % *MOSIAH-9-2 
		% ---
			\label{MOSIAH-9-2}

			Therefore I contended with my brethren in the wilderness,
			for I would that our ruler should make a treaty with them,
			but he being an austere and a bloodthirsty man commanded that I should be slain.
			But I was rescued by the shedding of much blood,
			for father fought against father and brother against brother
			until the greatest number of our army was destroyed in the wilderness;
			and we returned---those of us that were spared---to the land of Zarahemla
			to relate that tale to their wives and their children.

		% ---
		\paragraph{Mosiah 9:3} % *MOSIAH-9-3 
		% ---
			\label{MOSIAH-9-3}

			And yet I being overzealous to inherit the land of our fathers
			collected as many as were desirous to go up to possess the land
			and started again on our journey into the wilderness to go up to the land;
			but we were smitten with famine and sore affliction,
			for we were slow to remember the Lord our God.

		% ---
		\paragraph{Mosiah 9:4} % *MOSIAH-9-4 
		% ---
			\label{MOSIAH-9-4}

			Nevertheless, after many days wandering in the wilderness,
			we pitched our tents in the place where our brethren were slain,
			which was near to the land of our fathers.

		% ---
		\paragraph{Mosiah 9:5} % *MOSIAH-9-5 
		% ---
			\label{MOSIAH-9-5}

			And it came to pass that
			I went again with four of my men into the city in unto the king
			that I might know of the disposition of the king
			and that I might know if I might go in with my people
			and possess the land in peace.

		% ---
		\paragraph{Mosiah 9:6} % *MOSIAH-9-6 
		% ---
			\label{MOSIAH-9-6}

			And I went in unto the king and he covenanted with me
			that I might possess the land of Lehi-Nephi and the land of Shilom.

		% ---
		\paragraph{Mosiah 9:7} % *MOSIAH-9-7 
		% ---
			\label{MOSIAH-9-7}

			And he also commanded that his people should depart out of that land,
			and I and my people went into the land that we might possess it.

		% ---
		\paragraph{Mosiah 9:8} % *MOSIAH-9-8 
		% ---
			\label{MOSIAH-9-8}

			And we began to build buildings and to repair the walls of the city,
			yea, even the walls of the city of Lehi-Nephi and the city of Shilom.

		% ---
		\paragraph{Mosiah 9:9} % *MOSIAH-9-9 
		% ---
			\label{MOSIAH-9-9}

			And we began to till the ground,
			yea, even with all manner of seeds:
			with seeds of corn and of wheat and of barley
			and with neas and with sheum
			and with seeds of all manner of fruits.
			And we did begin to multiply and prosper in the land.

		% ---
		\paragraph{Mosiah 9:10} % *MOSIAH-9-10 
		% ---
			\label{MOSIAH-9-10}

			Now it was the cunning and the craftiness of king Laman
			to bring my people into bondage
			that he yielded up the land that we might possess it.

		% ---
		\paragraph{Mosiah 9:11} % *MOSIAH-9-11 
		% ---
			\label{MOSIAH-9-11}

			Therefore it came to pass that
			after we had dwelt in the land for the space of twelve years
			that king Laman began to grow uneasy
			lest by any means my people should wax strong in the land
			and that they could not overpower them and bring them into bondage.

		% ---
		\paragraph{Mosiah 9:12} % *MOSIAH-9-12 
		% ---
			\label{MOSIAH-9-12}

			Now they were a lazy and an idolatrous people.
			Therefore they were desirous to bring us into bondage
			that they might glut themselves with the labors of our hands,
			yea, that they might feast themselves upon the flocks of our fields.

		% ---
		\paragraph{Mosiah 9:13} % *MOSIAH-9-13 
		% ---
			\label{MOSIAH-9-13}

			Therefore it came to pass that king Laman began to stir up his people
			that they should contend with my people;
			therefore there began to be wars and contentions in the land.

		% ---
		\paragraph{Mosiah 9:14} % *MOSIAH-9-14 
		% ---
			\label{MOSIAH-9-14}

			For in the thirteenth year of my reign
			in the land of Nephi, away on the south of the land of Shilom,
			when my people were watering and feeding their flocks and tilling their lands,
			a numerous host of Lamanites came upon them and began to slay them
			and to take of their flocks and the corn of their fields.

		% ---
		\paragraph{Mosiah 9:15} % *MOSIAH-9-15 
		% ---
			\label{MOSIAH-9-15}

			Yea, and it came to pass that
			they fled---all that were not overtaken---even into the city of Nephi
			and did call upon me for protection.

		% ---
		\paragraph{Mosiah 9:16} % *MOSIAH-9-16 
		% ---
			\label{MOSIAH-9-16}

			And it came to pass that I did arm them with bows and with arrows,
			with swords and with scimitars and with clubs and with slings,
			and with all manner of weapons which we could invent.
			And I and my people did go forth against the Lamanites to battle;

		% ---
		\paragraph{Mosiah 9:17} % *MOSIAH-9-17 
		% ---
			\label{MOSIAH-9-17}

			yea, in the strength of the Lord did we go forth to battle against the Lamanites.
			For I and my people did cry mightily to the Lord
			that he would deliver us out of the hands of our enemies,
			for we were awakened to a remembrance of the deliverance of our fathers.

		% ---
		\paragraph{Mosiah 9:18} % *MOSIAH-9-18 
		% ---
			\label{MOSIAH-9-18}

			And God did hear our cries and did answer our prayers,
			and we did go forth in his might;
			yea, we did go forth against the Lamanites.
			And in one day and a night we did slay three thousand and forty three;
			we did slay them even until we had driven them out of our land.

		% ---
		\paragraph{Mosiah 9:19} % *MOSIAH-9-19 
		% ---
			\label{MOSIAH-9-19}

			And I myself with mine own hands did help bury their dead.
			And behold, to our great sorrow and lamentation,
			two hundred and seventy nine of our brethren were slain.

	% -----
	\subsubsection{Mosiah 10} % *MOSIAH-10 
	% -----
		\label{MOSIAH-10}
		
		% ---
		\paragraph{Mosiah 10:1} % *MOSIAH-10-1 
		% ---
			\label{MOSIAH-10-1}

			And it came to pass that we again began to establish the kingdom,
			and we again began to possess the land in peace.
			And I caused that there should be weapons of war made of every kind,
			that thereby I might have weapons for my people
			against the time the Lamanites should come up again to war against my people.

		% ---
		\paragraph{Mosiah 10:2} % *MOSIAH-10-2 
		% ---
			\label{MOSIAH-10-2}

			And I sat guards round about the land
			that the Lamanites might not come upon us again unawares and destroy us.
			And thus I did guard my people and my flocks
			and keep them from falling into the hands of our enemies.

		% ---
		\paragraph{Mosiah 10:3} % *MOSIAH-10-3 
		% ---
			\label{MOSIAH-10-3}

			And it came to pass that we did inherit the land of our fathers for many years,
			yea, for the space of twenty and two years.

		% ---
		\paragraph{Mosiah 10:4} % *MOSIAH-10-4 
		% ---
			\label{MOSIAH-10-4}

			And I did cause that the men should till the ground
			and raise all manner of grain and all manner of fruit of every kind.

		% ---
		\paragraph{Mosiah 10:5} % *MOSIAH-10-5 
		% ---
			\label{MOSIAH-10-5}

			And I did cause that the women should spin and toil
			and work all manner of fine linen,
			yea, and cloth of every kind,
			that we might clothe our nakedness.
			And thus we did prosper in the land;
			thus we did have continual peace in the land
			for the space of twenty and two years.

		% ---
		\paragraph{Mosiah 10:6} % *MOSIAH-10-6 
		% ---
			\label{MOSIAH-10-6}

			And it came to pass that king Laman died,
			and his son began to reign in his stead.
			And he began to stir his people up in rebellion against my people.
			Therefore they began to prepare for war
			and to come up to battle against my people,

		% ---
		\paragraph{Mosiah 10:7} % *MOSIAH-10-7 
		% ---
			\label{MOSIAH-10-7}

			but I having sent my spies out round about the land of Shemlon
			that I might discover their preparations,
			that I might guard against them,
			that they might not come upon my people and destroy them.

		% ---
		\paragraph{Mosiah 10:8} % *MOSIAH-10-8 
		% ---
			\label{MOSIAH-10-8}

			And it came to pass that they came up
			upon the north of the land of Shilom with their numerous hosts,
			men armed with bows and with arrows
			and with swords and with scimitars and with stones and with slings.
			And they had their heads shaved that they were naked,
			and they were girded about with a leathern girdle about their loins.

		% ---
		\paragraph{Mosiah 10:9} % *MOSIAH-10-9 
		% ---
			\label{MOSIAH-10-9}

			And it came to pass that I caused
			that the women and children of my people should be hid in the wilderness.
			And I also caused that all my old men that could bear arms
			and also all my young men that were able to bear arms
			should gather themselves together to go to battle against the Lamanites,
			and I did place them in their ranks every man according to his age.

		% ---
		\paragraph{Mosiah 10:10} % *MOSIAH-10-10 
		% ---
			\label{MOSIAH-10-10}

			And it came to pass that we did go up to battle against the Lamanites;
			and I, even I in my old age, did go up to battle against the Lamanites.
			And it came to pass that we did go up in the strength of the Lord to battle.

		% ---
		\paragraph{Mosiah 10:11} % *MOSIAH-10-11 
		% ---
			\label{MOSIAH-10-11}

			Now the Lamanites knew nothing concerning the Lord nor the strength of the Lord;
			therefore they depended upon their own strength,
			yet they were a strong people as to the strength of men.

		% ---
		\paragraph{Mosiah 10:12} % *MOSIAH-10-12 
		% ---
			\label{MOSIAH-10-12}

			They were a wild and ferocious and a bloodthirsty people,
			believing in the tradition of their fathers, which is this:
			believing that they were driven out of the land of Jerusalem
			because of the iniquities of their fathers,
			and that they were wronged in the wilderness by their brethren,
			and they were also wronged while crossing the sea,

		% ---
		\paragraph{Mosiah 10:13} % *MOSIAH-10-13 
		% ---
			\label{MOSIAH-10-13}

			and again, that they were wronged while in the land of their first inheritance
			after they had crossed the sea---
			and all this because that Nephi was more faithful
			in keeping the commandments of the Lord;
			therefore he was favored of the Lord,
			for the Lord heard his prayers and answered them;
			and he took the lead of their journey in the wilderness.

		% ---
		\paragraph{Mosiah 10:14} % *MOSIAH-10-14 
		% ---
			\label{MOSIAH-10-14}

			And his brethren were wroth with him
			because they understood not the dealings of the Lord.
			They were also wroth with him upon the waters
			because they hardened their hearts against the Lord.

		% ---
		\paragraph{Mosiah 10:15} % *MOSIAH-10-15 
		% ---
			\label{MOSIAH-10-15}

			And again, they were wroth with him when they had arriven to the promised land
			because they said that he had taken the ruling of the people out of their hands;
			and they sought to kill him.

		% ---
		\paragraph{Mosiah 10:16} % *MOSIAH-10-16 
		% ---
			\label{MOSIAH-10-16}

			And again, they were wroth with him
			because he departed into the wilderness as the Lord had commanded him
			and took the records which were engraven on the plates of brass,
			for they said that he robbed them.

		% ---
		\paragraph{Mosiah 10:17} % *MOSIAH-10-17 
		% ---
			\label{MOSIAH-10-17}

			And thus they have taught their children that they should hate them,
			and that they should murder them,
			and that they should rob and plunder them
			and do all they could to destroy them.
			Therefore they have an eternal hatred towards the children of Nephi.

		% ---
		\paragraph{Mosiah 10:18} % *MOSIAH-10-18 
		% ---
			\label{MOSIAH-10-18}

			For this very cause hath king Laman
			by his cunning and lying craftiness and his fair promises hath deceived me,
			that I have brought this my people up into this land,
			that they may destroy them;
			yea, and we have suffered this many years in the land.

		% ---
		\paragraph{Mosiah 10:19} % *MOSIAH-10-19 
		% ---
			\label{MOSIAH-10-19}

			And now I Zeniff,
			after having told all these things unto my people concerning the Lamanites,
			I did stimulate them to go to battle with their might,
			putting their trust in the Lord.
			Therefore we did contend with them face to face.

		% ---
		\paragraph{Mosiah 10:20} % *MOSIAH-10-20 
		% ---
			\label{MOSIAH-10-20}

			And it came to pass that we did drive them again out of our land;
			and we slew them with a great slaughter,
			even so many that we did not number them.

		% ---
		\paragraph{Mosiah 10:21} % *MOSIAH-10-21 
		% ---
			\label{MOSIAH-10-21}

			And it came to pass that we returned again to our own land,
			and my people again began to tend their flocks and to till their ground.
			
		% ---
		\paragraph{Mosiah 10:22} % *MOSIAH-10-22
		% ---
			\label{MOSIAH-10-22}
			
			And now I being old did confer the kingdom upon one of my sons;
			therefore I say no more.
			And may the Lord bless my people!
			Amen.

	% -----
	\subsubsection{Mosiah 11} % *MOSIAH-11 
	% -----
		\label{MOSIAH-11}
		
		% ---
		\paragraph{Mosiah 11:1} % *MOSIAH-11-1 
		% ---
			\label{MOSIAH-11-1}

			And now it came to pass that
			Zeniff conferred the kingdom upon Noah, one of his sons;
			therefore Noah began to reign in his stead.
			And he did not walk in the ways of his father;

		% ---
		\paragraph{Mosiah 11:2} % *MOSIAH-11-2 
		% ---
			\label{MOSIAH-11-2}

			for behold, he did not keep the commandments of God,
			but he did walk after the desires of his own heart.
			And he had many wives and concubines.
			And he did cause his people to commit sin
			and do that which was abominable in the sight of the Lord;
			yea, and they did commit whoredoms and all manner of wickedness.

		% ---
		\paragraph{Mosiah 11:3} % *MOSIAH-11-3 
		% ---
			\label{MOSIAH-11-3}

			And he laid a tax of one fifth part of all they possessed:
			a fifth part of their gold and of their silver,
			and a fifth part of their ziff and of their copper and of their brass and their iron,
			and a fifth part of their fatlings,
			and also a fifth part of all their grain.

		% ---
		\paragraph{Mosiah 11:4} % *MOSIAH-11-4 
		% ---
			\label{MOSIAH-11-4}

			And all this did he take to support himself and his wives and his concubines,
			and also his priests and their wives and their concubines.
			Thus he had changed the affairs of the kingdom,

		% ---
		\paragraph{Mosiah 11:5} % *MOSIAH-11-5 
		% ---
			\label{MOSIAH-11-5}

			for he put down all the priests that had been consecrated by his father
			and consecrated new ones in their stead
			such as were lifted up in the pride of their hearts.

		% ---
		\paragraph{Mosiah 11:6} % *MOSIAH-11-6 
		% ---
			\label{MOSIAH-11-6}

			Yea, and thus were they supported in their laziness
			and in their idolatry and in their whoredoms
			by the taxes which king Noah had put upon his people.
			Thus did the people labor exceedingly to support iniquity.

		% ---
		\paragraph{Mosiah 11:7} % *MOSIAH-11-7 
		% ---
			\label{MOSIAH-11-7}

			Yea, and they also became idolatrous
			because they were deceived
			by the vain and flattering words of the king and priests,
			for they did speak flattering things unto them.

		% ---
		\paragraph{Mosiah 11:8} % *MOSIAH-11-8 
		% ---
			\label{MOSIAH-11-8}

			And it came to pass that
			king Noah built many elegant and spacious buildings,
			and he ornamented them with fine work of wood
			and of all manner of precious things,
			of gold and of silver and of iron and of brass and of ziff and of copper.

		% ---
		\paragraph{Mosiah 11:9} % *MOSIAH-11-9 
		% ---
			\label{MOSIAH-11-9}

			And he also built him a spacious palace and a throne in the midst thereof,
			all of which was of fine wood
			and was ornamented with gold and silver and with precious things.

		% ---
		\paragraph{Mosiah 11:10} % *MOSIAH-11-10 
		% ---
			\label{MOSIAH-11-10}

			And he also caused that his workmen should work
			all manner of fine work within the walls of the temple,
			of fine wood and of copper and of brass.

		% ---
		\paragraph{Mosiah 11:11} % *MOSIAH-11-11 
		% ---
			\label{MOSIAH-11-11}

			And the seats which was sat apart for the high priests,
			which was above all the other seats,
			he did ornament with pure gold.
			And he caused a breastwork to be built before them
			that they might rest their bodies and their arms upon
			while they should speak lying and vain words to his people.

		% ---
		\paragraph{Mosiah 11:12} % *MOSIAH-11-12 
		% ---
			\label{MOSIAH-11-12}

			And it came to pass that he built a tower near the temple,
			yea, a very high tower,
			even so high that he could stand upon the top thereof
			and overlook the land of Shilom
			and also the land of Shemlon, which was possessed by the Lamanites;
			and he could even look over all the land round about.

		% ---
		\paragraph{Mosiah 11:13} % *MOSIAH-11-13 
		% ---
			\label{MOSIAH-11-13}

			And it came to pass that he caused many buildings to be built in the land Shilom.
			And he caused a great tower to be built on the hill north of the land Shilom,
			which had been a resort for the children of Nephi
			at the time they fled out of the land.
			And thus he did do with the riches which he obtained by the taxation of his people.

		% ---
		\paragraph{Mosiah 11:14} % *MOSIAH-11-14 
		% ---
			\label{MOSIAH-11-14}

			And it came to pass that he placed his heart upon his riches,
			and he spent his time in riotous living with his wives and his concubines;
			and so did also his priests spend their time with harlots.

		% ---
		\paragraph{Mosiah 11:15} % *MOSIAH-11-15 
		% ---
			\label{MOSIAH-11-15}

			And it came to pass that he planted vineyards round about in the land,
			and he built winepresses and made wine in abundance;
			and therefore he became a winebibber, and also his people.

		% ---
		\paragraph{Mosiah 11:16} % *MOSIAH-11-16 
		% ---
			\label{MOSIAH-11-16}

			And it came to pass that the Lamanites began
			to come in upon his people, upon small numbers,
			and to slay them in their fields and while they were tending their flocks.

		% ---
		\paragraph{Mosiah 11:17} % *MOSIAH-11-17 
		% ---
			\label{MOSIAH-11-17}

			And king Noah sent guards round about the land to keep them off,
			but he did not send a sufficient number;
			and the Lamanites came upon them and killed them
			and drove many of their flocks out of the land.
			Thus the Lamanites began to destroy them and to exercise their hatred upon them.

		% ---
		\paragraph{Mosiah 11:18} % *MOSIAH-11-18 
		% ---
			\label{MOSIAH-11-18}

			And it came to pass that king Noah sent his armies against them;
			and they were driven back,
			or they drove them back for a time.
			Therefore they returned rejoicing in their spoil.

		% ---
		\paragraph{Mosiah 11:19} % *MOSIAH-11-19 
		% ---
			\label{MOSIAH-11-19}

			And now because of this great victory
			they were lifted up in the pride of their hearts;
			they did boast in their own strength,
			saying that their fifty could stand against thousands of the Lamanites.
			And thus they did boast and did delight
			in blood and the shedding of the blood of their brethren---
			and this because of the wickedness of their king and priests.

		% ---
		\paragraph{Mosiah 11:20} % *MOSIAH-11-20 
		% ---
			\label{MOSIAH-11-20}

			And it came to pass that there was a man among them whose name was Abinadi;
			and he went forth among them and began to prophesy, saying:
			Behold, thus saith the Lord and thus hath he commanded me, saying:
			Go forth and say unto this people:
			Thus saith the Lord:
			Woe be unto this people,
			for I have seen their abominations and their wickedness and their whoredoms;
			and except they repent, I will visit them in mine anger.

		% ---
		\paragraph{Mosiah 11:21} % *MOSIAH-11-21 
		% ---
			\label{MOSIAH-11-21}

			And except they repent and turn to the Lord their God,
			behold, I will deliver them into the hands of their enemies;
			yea, and they shall be brought into bondage,
			and they shall be afflicted by the hand of their enemies.

		% ---
		\paragraph{Mosiah 11:22} % *MOSIAH-11-22 
		% ---
			\label{MOSIAH-11-22}

			And it shall come to pass that they shall know
			that I am the Lord their God and am a jealous God,
			visiting the iniquities of my people.

		% ---
		\paragraph{Mosiah 11:23} % *MOSIAH-11-23 
		% ---
			\label{MOSIAH-11-23}

			And it shall come to pass that
			except this people repent and turn to the Lord their God,
			they shall be brought into bondage;
			and none shall deliver them except it be the Lord the Almighty God.

		% ---
		\paragraph{Mosiah 11:24} % *MOSIAH-11-24 
		% ---
			\label{MOSIAH-11-24}

			Yea, and it shall come to pass that
			when they shall cry unto me,
			I will be slow to hear their cries.
			Yea, and I will suffer them that they be smitten by their enemies.

		% ---
		\paragraph{Mosiah 11:25} % *MOSIAH-11-25 
		% ---
			\label{MOSIAH-11-25}

			And except they repent in sackcloth and ashes
			and cry mightily to the Lord their God,
			I will not hear their prayers,
			neither will I deliver them out of their afflictions.
			And thus saith the Lord and thus hath he commanded me.

		% ---
		\paragraph{Mosiah 11:26} % *MOSIAH-11-26 
		% ---
			\label{MOSIAH-11-26}

			Now it came to pass that when Abinadi had spake these words unto them,
			they were wroth with him and sought to take away his life;
			but the Lord delivered him out of their hands.

		% ---
		\paragraph{Mosiah 11:27} % *MOSIAH-11-27 
		% ---
			\label{MOSIAH-11-27}

			Now when king Noah had heard of the words
			which Abinadi had spake unto the people,
			he was also wroth and he saith:
			Who is Abinadi, that I and my people should be judged of him?
			Or who is the Lord that shall bring upon my people such great affliction?

		% ---
		\paragraph{Mosiah 11:28} % *MOSIAH-11-28 
		% ---
			\label{MOSIAH-11-28}

			I command you to bring Abinadi hither that I may slay him,
			for he hath said these things
			that he might stir up my people to anger one with another
			and to raise contentions among my people;
			therefore I will slay him.

		% ---
		\paragraph{Mosiah 11:29} % *MOSIAH-11-29 
		% ---
			\label{MOSIAH-11-29}

			Now the eyes of the people were blinded;
			therefore they hardened their hearts against the words of Abinadi,
			and they sought from that time forward to take him.
			And king Noah hardened his heart against the word of the Lord,
			and he did not repent of his evil doings.

	% -----
	\subsubsection{Mosiah 12} % *MOSIAH-12 
	% -----
		\label{MOSIAH-12}
		
		% ---
		\paragraph{Mosiah 12:1} % *MOSIAH-12-1 
		% ---
			\label{MOSIAH-12-1}

			And it came to pass that after the space of two years
			that Abinadi came among them in disguise,
			that they knew him not,
			and began again to prophesy among them, saying:
			Thus hath the Lord commanded me, saying:
			Abinadi, go and prophesy unto this my people,
			for they have hardened their hearts against my words;
			they have repented not of their evil doings.
			Therefore I will visit them in my anger;
			yea, in my fierce anger will I visit them in their iniquities and abominations;

		% ---
		\paragraph{Mosiah 12:2} % *MOSIAH-12-2 
		% ---
			\label{MOSIAH-12-2}

			yea, woe be unto this generation.
			And the Lord said unto me:
			Stretch forth thy hand and prophesy, saying:
			Thus saith the Lord:
			It shall come to pass that this generation because of their iniquities
			shall be brought into bondage and shall be smitten on the cheek,
			yea, and shall be driven by men and shall be slain.
			And the vultures of the air and the dogs
			---yea, and the wild beasts---
			shall devour their flesh.

		% ---
		\paragraph{Mosiah 12:3} % *MOSIAH-12-3 
		% ---
			\label{MOSIAH-12-3}

			And it shall come to pass that
			the life of king Noah shall be valued even as a garment in a hot furnace,
			for he shall know that I am the Lord.

		% ---
		\paragraph{Mosiah 12:4} % *MOSIAH-12-4 
		% ---
			\label{MOSIAH-12-4}

			And it shall come to pass that
			I will smite this my people with sore afflictions,
			yea, with famine and with pestilence.
			And I will cause that they shall howl all the day long.

		% ---
		\paragraph{Mosiah 12:5} % *MOSIAH-12-5 
		% ---
			\label{MOSIAH-12-5}

			Yea, and I will cause that they shall have burdens lashed upon their backs,
			and they shall be driven before like a dumb ass.

		% ---
		\paragraph{Mosiah 12:6} % *MOSIAH-12-6 
		% ---
			\label{MOSIAH-12-6}

			And it shall come to pass that I will send forth hail among them,
			and it shall smite them;
			and they shall also be smitten with the east wind;
			and insects shall pester their land also and devour their grain;

		% ---
		\paragraph{Mosiah 12:7} % *MOSIAH-12-7 
		% ---
			\label{MOSIAH-12-7}

			and they shall be smitten with a great pestilence.
			And all this will I do because of their iniquities and abominations.

		% ---
		\paragraph{Mosiah 12:8} % *MOSIAH-12-8 
		% ---
			\label{MOSIAH-12-8}

			And it shall come to pass that except they repent,
			I will utterly destroy them from off the face of the earth.
			Yet they shall leave a record behind them,
			and I will preserve them for other nations which shall possess the land.
			Yea, even this will I do
			that I may discover the abominations of this people to other nations.
			And many things did Abinadi prophesy against this people.

		% ---
		\paragraph{Mosiah 12:9} % *MOSIAH-12-9 
		% ---
			\label{MOSIAH-12-9}

			And it came to pass that they were angry with him;
			and they took him and carried him bound before the king
			and saith unto the king:
			Behold, we have brought a man before thee
			which has prophesied evil concerning thy people
			and saith that God will destroy them.

		% ---
		\paragraph{Mosiah 12:10} % *MOSIAH-12-10 
		% ---
			\label{MOSIAH-12-10}

			And he also prophesieth evil concerning thy life
			and saith that thy life shall be as a garment in a furnace of fire.

		% ---
		\paragraph{Mosiah 12:11} % *MOSIAH-12-11 
		% ---
			\label{MOSIAH-12-11}

			And again, he saith that thou shall be as a stalk,
			even as a dry stalk of the field,
			which is ran over by the beasts and trodden under foot.

		% ---
		\paragraph{Mosiah 12:12} % *MOSIAH-12-12 
		% ---
			\label{MOSIAH-12-12}

			And again, he saith thou shalt be as the blossoms of a thistle,
			which when it is fully ripe, if the wind bloweth,
			it is driven forth upon the face of the land.
			And he pretendeth the Lord hath spoken it.
			And he saith all this shall come upon thee except thou repent---
			and this because of thine iniquities.

		% ---
		\paragraph{Mosiah 12:13} % *MOSIAH-12-13 
		% ---
			\label{MOSIAH-12-13}

			And now, O king, what great evil hast thou done?
			Or what great sins has thy people committed
			that we should be condemned of God or judged of this man?

		% ---
		\paragraph{Mosiah 12:14} % *MOSIAH-12-14 
		% ---
			\label{MOSIAH-12-14}

			And now, O king, behold, we are guiltless.
			And thou, O king, hast not sinned.
			Therefore this man hath lied concerning you,
			and he hath prophesied in vain.

		% ---
		\paragraph{Mosiah 12:15} % *MOSIAH-12-15 
		% ---
			\label{MOSIAH-12-15}

			And behold, we are strong;
			we shall not come into bondage or be taken captive by our enemies.
			Yea, and thou hast prospered in the land, and thou shalt also prosper.

		% ---
		\paragraph{Mosiah 12:16} % *MOSIAH-12-16 
		% ---
			\label{MOSIAH-12-16}

			Behold, here is the man.
			We deliver him into thy hands.
			Thou mayest do with him as seemeth thee good.

		% ---
		\paragraph{Mosiah 12:17} % *MOSIAH-12-17 
		% ---
			\label{MOSIAH-12-17}

			And it came to pass that king Noah caused
			that Abinadi should be cast into prison.
			And he commanded that the priests should gather themselves together
			that he might hold a council with them what he should do with him.

		% ---
		\paragraph{Mosiah 12:18} % *MOSIAH-12-18 
		% ---
			\label{MOSIAH-12-18}

			And it came to pass that they saith unto the king:
			Bring him hither that we may question him.
			And the king commanded that he should be brought before them.

		% ---
		\paragraph{Mosiah 12:19} % *MOSIAH-12-19 
		% ---
			\label{MOSIAH-12-19}

			And they began to question him, that they might cross him,
			that thereby they might have wherewith to accuse him.
			But he answered them boldly and withstood all their questions
			---yea, to their astonishment---
			for he did withstand them in all their questions
			and did confound them in all their words.

		% ---
		\paragraph{Mosiah 12:20} % *MOSIAH-12-20 
		% ---
			\label{MOSIAH-12-20}

			And it came to pass that one of them saith unto him:
			What meaneth the words which are written
			and which have been taught by our fathers, saying:

		% ---
		\paragraph{Mosiah 12:21} % *MOSIAH-12-21 
		% ---
			\label{MOSIAH-12-21}

			How beautiful upon the mountains are the feet of him
			that bringeth good tidings,
			that publisheth peace,
			that bringeth good tidings of good,
			that publisheth salvation,
			that saith unto Zion:
			Thy God reigneth!

		% ---
		\paragraph{Mosiah 12:22} % *MOSIAH-12-22 
		% ---
			\label{MOSIAH-12-22}

			Thy watchmen shall lift up the voice;
			with the voice together shall they sing.
			For they shall see eye to eye when the Lord shall bring again Zion.

		% ---
		\paragraph{Mosiah 12:23} % *MOSIAH-12-23 
		% ---
			\label{MOSIAH-12-23}

			Break forth into joy!
			Sing together, ye waste places of Jerusalem!
			For the Lord hath comforted his people;
			he hath redeemed Jerusalem.

		% ---
		\paragraph{Mosiah 12:24} % *MOSIAH-12-24 
		% ---
			\label{MOSIAH-12-24}

			The Lord hath made bare his holy arm in the eyes of all the nations,
			and all the ends of the earth shall see the salvation of our God.

		% ---
		\paragraph{Mosiah 12:25} % *MOSIAH-12-25 
		% ---
			\label{MOSIAH-12-25}

			And now Abinadi saith unto them:
			Are you priests and pretend to teach this people
			and to understand the spirit of prophesying,
			and yet desireth to know of me what these things mean?

		% ---
		\paragraph{Mosiah 12:26} % *MOSIAH-12-26 
		% ---
			\label{MOSIAH-12-26}

			I say unto you:
			Woe be unto you for perverting the ways of the Lord!
			For if ye understand these things, ye have not taught them.
			Therefore ye have perverted the ways of the Lord.

		% ---
		\paragraph{Mosiah 12:27} % *MOSIAH-12-27 
		% ---
			\label{MOSIAH-12-27}

			Ye have not applied your hearts to understanding;
			therefore ye have not been wise.
			Therefore what teachest thou this people?

		% ---
		\paragraph{Mosiah 12:28} % *MOSIAH-12-28 
		% ---
			\label{MOSIAH-12-28}

			And they said:
			We teach the law of Moses.

		% ---
		\paragraph{Mosiah 12:29} % *MOSIAH-12-29 
		% ---
			\label{MOSIAH-12-29}

			And again he saith unto them:
			If ye teach the law of Moses,
			why do ye not keep it?
			Why do ye set your hearts upon riches?
			Why do ye commit whoredoms and spend your strength with harlots,
			yea, and cause this people to commit sin,
			that the Lord hath cause to send me to prophesy against this people?
			---yea, even a great evil against this people.

		% ---
		\paragraph{Mosiah 12:30} % *MOSIAH-12-30 
		% ---
			\label{MOSIAH-12-30}

			Knowest thou not that I speak the truth?
			Yea, thou knowest that I speak the truth;
			and you had ought to tremble before God.

		% ---
		\paragraph{Mosiah 12:31} % *MOSIAH-12-31 
		% ---
			\label{MOSIAH-12-31}

			And it shall come to pass that ye shall be smitten for thine iniquities,
			for ye have said that ye teach the law of Moses.
			And what knowest thou concerning the law of Moses?
			Doth salvation come by the law of Moses?
			What sayest thou?


		% ---
		\paragraph{Mosiah 12:32} % *MOSIAH-12-32 
		% ---
			\label{MOSIAH-12-32}

			And they answered and said
			that salvation did come by the law of Moses.

		% ---
		\paragraph{Mosiah 12:33} % *MOSIAH-12-33 
		% ---
			\label{MOSIAH-12-33}

			But now Abinadi saith unto them:
			I know if ye keep the commandments of God, ye shall be saved---
			yea, if ye keep the commandments
			which the Lord delivered unto Moses in the mount of Sinai, saying:

		% ---
		\paragraph{Mosiah 12:34} % *MOSIAH-12-34 
		% ---
			\label{MOSIAH-12-34}

			I am the Lord thy God,
			which have brought thee out of the land of Egypt,
			out of the house of bondage.

		% ---
		\paragraph{Mosiah 12:35} % *MOSIAH-12-35 
		% ---
			\label{MOSIAH-12-35}

			Thou shalt have no other God before me.

		% ---
		\paragraph{Mosiah 12:36} % *MOSIAH-12-36 
		% ---
			\label{MOSIAH-12-36}

			Thou shalt not make unto thee any graven image,
			or any likeness of any thing in the heaven above,
			or things which is in the earth beneath.

		% ---
		\paragraph{Mosiah 12:37} % *MOSIAH-12-37 
		% ---
			\label{MOSIAH-12-37}

			Now Abinadi saith unto them:
			Have ye done all this?
			I say unto you:
			Nay, ye have not.
			And have ye taught this people that they should do all these things?
			I say unto you:
			Nay, ye have not.

	% -----
	\subsubsection{Mosiah 13} % *MOSIAH-13 
	% -----
		\label{MOSIAH-13}
		
		% ---
		\paragraph{Mosiah 13:1} % *MOSIAH-13-1 
		% ---
			\label{MOSIAH-13-1}

			And now when the king had heard these words,
			he said unto his priests:
			Away with this fellow and slay him!
			For what have we to do with him?
			---for he is mad!

		% ---
		\paragraph{Mosiah 13:2} % *MOSIAH-13-2 
		% ---
			\label{MOSIAH-13-2}

			And they stood forth and attempted to lay their hands on him,
			but he withstood them and said unto them:

		% ---
		\paragraph{Mosiah 13:3} % *MOSIAH-13-3 
		% ---
			\label{MOSIAH-13-3}

			Touch me not!
			For God shall smite you if ye lay your hands upon me,
			for I have not delivered the message which the Lord sent me to deliver,
			neither have I told you that which ye requested that I should tell.
			Therefore God will not suffer that I shall be destroyed at this time.

		% ---
		\paragraph{Mosiah 13:4} % *MOSIAH-13-4 
		% ---
			\label{MOSIAH-13-4}

			But I must fulfill the commandments wherewith God hath commanded me.
			And because I have told you the truth, ye are angry with me.
			And again, because I have spoken the word of God,
			ye have judged me that I am mad.

		% ---
		\paragraph{Mosiah 13:5} % *MOSIAH-13-5 
		% ---
			\label{MOSIAH-13-5}

			Now it came to pass after Abinadi had spoken these words
			that the people of king Noah durst not lay their hands on him,
			for the Spirit of the Lord was upon him.
			And his face shone with exceeding luster even as Moses' did
			while in the mount of Sinai while speaking with the Lord.

		% ---
		\paragraph{Mosiah 13:6} % *MOSIAH-13-6 
		% ---
			\label{MOSIAH-13-6}

			And he spake with power and authority from God.
			And he continued his words, saying:

		% ---
		\paragraph{Mosiah 13:7} % *MOSIAH-13-7 
		% ---
			\label{MOSIAH-13-7}

			Ye see that ye have not power to slay me.
			Therefore I finish my message.
			Yea, and I perceive that it cuts you to your hearts
			because I tell you the truth concerning your iniquities.

		% ---
		\paragraph{Mosiah 13:8} % *MOSIAH-13-8 
		% ---
			\label{MOSIAH-13-8}

			Yea, and my words fill you with wonder and amazement and with anger.

		% ---
		\paragraph{Mosiah 13:9} % *MOSIAH-13-9 
		% ---
			\label{MOSIAH-13-9}

			But I finish my message,
			and then it matters not whither I go,
			if it so be that I am saved.

		% ---
		\paragraph{Mosiah 13:10} % *MOSIAH-13-10 
		% ---
			\label{MOSIAH-13-10}

			But this much I tell you:
			What you do with me after this
			shall be as a type and a shadow of things which is to come.

		% ---
		\paragraph{Mosiah 13:11} % *MOSIAH-13-11 
		% ---
			\label{MOSIAH-13-11}

			And now I read unto you the remainder of the commandments of God,
			for I perceive that they are not written in your hearts.
			I perceive that ye have studied and taught iniquity the most part of your lives.

		% ---
		\paragraph{Mosiah 13:12} % *MOSIAH-13-12 
		% ---
			\label{MOSIAH-13-12}

			And now ye remember that I said unto you:
			Thou shalt not make unto thee any graven image,
			or any likeness of things which is in heaven above,
			or which is in the earth beneath,
			or which is in the water under the earth.

		% ---
		\paragraph{Mosiah 13:13} % *MOSIAH-13-13 
		% ---
			\label{MOSIAH-13-13}

			And again, thou shalt not bow down thyself unto them nor serve them,
			for I the Lord thy God am a jealous God,
			visiting the iniquities of the fathers upon the children
			unto the third and fourth generation of them that hate me,

		% ---
		\paragraph{Mosiah 13:14} % *MOSIAH-13-14 
		% ---
			\label{MOSIAH-13-14}

			and shewing mercy unto thousands of them
			that love me and keep my commandments.

		% ---
		\paragraph{Mosiah 13:15} % *MOSIAH-13-15 
		% ---
			\label{MOSIAH-13-15}

			Thou shalt not take the name of the Lord thy God in vain,
			for the Lord will not hold him guiltless that taketh his name in vain.

		% ---
		\paragraph{Mosiah 13:16} % *MOSIAH-13-16 
		% ---
			\label{MOSIAH-13-16}

			Remember the sabbath day, to keep it holy.

		% ---
		\paragraph{Mosiah 13:17} % *MOSIAH-13-17 
		% ---
			\label{MOSIAH-13-17}

			Six days shalt thou labor and do all thy work;

		% ---
		\paragraph{Mosiah 13:18} % *MOSIAH-13-18 
		% ---
			\label{MOSIAH-13-18}

			but the seventh day, the sabbath of the Lord thy God,
			thou shalt not do any work, thou nor thy son nor thy daughter,
			thy manservant nor thy maidservant nor thy cattle
			nor thy stranger that is within thy gates.

		% ---
		\paragraph{Mosiah 13:19} % *MOSIAH-13-19 
		% ---
			\label{MOSIAH-13-19}

			For in six days the Lord made heaven and earth and the sea and all that in them is;
			wherefore the Lord blessed the sabbath day and hallowed it.

		% ---
		\paragraph{Mosiah 13:20} % *MOSIAH-13-20 
		% ---
			\label{MOSIAH-13-20}

			Honor thy father and thy mother,
			that thy days may be long upon the land
			which the Lord thy God giveth thee.

		% ---
		\paragraph{Mosiah 13:21} % *MOSIAH-13-21 
		% ---
			\label{MOSIAH-13-21}

			Thou shalt not kill.

		% ---
		\paragraph{Mosiah 13:22} % *MOSIAH-13-22 
		% ---
			\label{MOSIAH-13-22}

			Thou shalt not commit adultery.
			Thou shalt not steal.

		% ---
		\paragraph{Mosiah 13:23} % *MOSIAH-13-23 
		% ---
			\label{MOSIAH-13-23}

			Thou shalt not bear false witness against thy neighbor.

		% ---
		\paragraph{Mosiah 13:24} % *MOSIAH-13-24 
		% ---
			\label{MOSIAH-13-24}

			Thou shalt not covet thy neighbor's house;
			thou shalt not covet thy neighbor's wife,
			nor his manservant nor his maidservant nor his ox nor his ass,
			nor any thing that is thy neighbor's.

		% ---
		\paragraph{Mosiah 13:25} % *MOSIAH-13-25 
		% ---
			\label{MOSIAH-13-25}

			And it came to pass that
			after Abinadi had made an end of these sayings
			that he said unto them:
			Have ye taught this people that they should observe to do all these things,
			for to keep these commandments?

		% ---
		\paragraph{Mosiah 13:26} % *MOSIAH-13-26 
		% ---
			\label{MOSIAH-13-26}

			I say unto you: Nay.
			For if ye had,
			the Lord would not have caused me to come forth
			and to prophesy evil concerning this people.

		% ---
		\paragraph{Mosiah 13:27} % *MOSIAH-13-27 
		% ---
			\label{MOSIAH-13-27}

			And now ye have said that salvation cometh by the law of Moses.
			I say unto you that it is expedient that ye should keep the law of Moses as yet;
			but I say unto you that the time shall come
			when it shall no more be expedient to keep the law of Moses.

		% ---
		\paragraph{Mosiah 13:28} % *MOSIAH-13-28 
		% ---
			\label{MOSIAH-13-28}

			And moreover I say unto you
			that salvation doth not come by the law alone.
			And were it not for the atonement
			which God himself shall make for the sins and iniquities of his people
			that they must unavoidably perish,
			notwithstanding the law of Moses.

		% ---
		\paragraph{Mosiah 13:29} % *MOSIAH-13-29 
		% ---
			\label{MOSIAH-13-29}

			And now I say unto you that it was expedient
			that there should be a law given to the children of Israel,
			yea, even a very strict law.
			For they were a stiffnecked people,
			quick to do iniquity and slow to remember the Lord their God.

		% ---
		\paragraph{Mosiah 13:30} % *MOSIAH-13-30 
		% ---
			\label{MOSIAH-13-30}

			Therefore there was a law given them,
			yea, a law of performances and of ordinances,
			a law which they were to observe strictly from day to day
			to keep them in remembrance of God and their duty towards him.

		% ---
		\paragraph{Mosiah 13:31} % *MOSIAH-13-31 
		% ---
			\label{MOSIAH-13-31}

			But behold, I say unto you that all these things were types of things to come.

		% ---
		\paragraph{Mosiah 13:32} % *MOSIAH-13-32 
		% ---
			\label{MOSIAH-13-32}

			And now, did they understand the law?
			I say unto you:
			Nay, they did not all understand the law---
			and this because of the hardness of their hearts.
			For they understood not that there could not any man be saved
			except it were through the redemption of God.

		% ---
		\paragraph{Mosiah 13:33} % *MOSIAH-13-33 
		% ---
			\label{MOSIAH-13-33}

			For behold, did not Moses prophesy unto them
			concerning the coming of the Messiah
			and that God should redeem his people?
			Yea, and even all the prophets which have prophesied ever since the world began,
			have they not spoken more or less concerning these things?

		% ---
		\paragraph{Mosiah 13:34} % *MOSIAH-13-34 
		% ---
			\label{MOSIAH-13-34}

			Have they not said
			that God himself should come down among the children of men
			and take upon him the form of man
			and go forth in mighty power upon the face of the earth?

		% ---
		\paragraph{Mosiah 13:35} % *MOSIAH-13-35 
		% ---
			\label{MOSIAH-13-35}

			Yea, and have they not said also
			that he should bring to pass the resurrection of the dead
			and that he himself should be oppressed and afflicted?
			
	% -----
	\subsubsection{Mosiah 14} % *MOSIAH-14
	% -----
		\label{MOSIAH-14}
		
		% ---
		\paragraph{Mosiah 14:1} % *MOSIAH-14-1 
		% ---
			\label{MOSIAH-14-1}

			Yea, even doth not Isaiah say:
			Who hath believed our report?
			And to whom is the arm of the Lord revealed?

		% ---
		\paragraph{Mosiah 14:2} % *MOSIAH-14-2 
		% ---
			\label{MOSIAH-14-2}

			For he shall grow up before him
			as a tender plant and as a root out of dry ground.
			He hath no form nor comeliness.
			And when we shall see him,
			there is no beauty that we should desire him.

		% ---
		\paragraph{Mosiah 14:3} % *MOSIAH-14-3 
		% ---
			\label{MOSIAH-14-3}

			He is despised and rejected of men,
			a man of sorrows and acquainted with grief.
			And we hid, as it were, our faces from him.
			He was despised and we esteemed him not.

		% ---
		\paragraph{Mosiah 14:4} % *MOSIAH-14-4 
		% ---
			\label{MOSIAH-14-4}

			Surely he hath borne our griefs and carried our sorrows.
			Yet we did esteem him stricken,
			smitten of God and afflicted.

		% ---
		\paragraph{Mosiah 14:5} % *MOSIAH-14-5 
		% ---
			\label{MOSIAH-14-5}

			But he was wounded for our transgressions.
			He was bruised for our iniquities.
			The chastisement of our peace was upon him,
			and with his stripes we are healed.

		% ---
		\paragraph{Mosiah 14:6} % *MOSIAH-14-6 
		% ---
			\label{MOSIAH-14-6}

			All we like sheep have gone astray.
			We have turned every one to his own way.
			And the Lord hath laid on him the iniquities of us all.

		% ---
		\paragraph{Mosiah 14:7} % *MOSIAH-14-7 
		% ---
			\label{MOSIAH-14-7}

			He was oppressed and he was afflicted,
			yet he opened not his mouth.
			He is brought as a lamb to the slaughter;
			and as a sheep before her shearers is dumb,
			so he opened not his mouth.

		% ---
		\paragraph{Mosiah 14:8} % *MOSIAH-14-8 
		% ---
			\label{MOSIAH-14-8}

			He was taken from prison and from judgment.
			And who shall declare his generation?
			For he was cut off out of the land of the living.
			For the transgressions of my people was he stricken.

		% ---
		\paragraph{Mosiah 14:9} % *MOSIAH-14-9 
		% ---
			\label{MOSIAH-14-9}

			And he made his grave with the wicked,
			and with the rich in his death---
			because he had done no evil,
			neither was any deceit in his mouth.

		% ---
		\paragraph{Mosiah 14:10} % *MOSIAH-14-10 
		% ---
			\label{MOSIAH-14-10}

			Yet it pleased the Lord to bruise him.
			He hath put him to grief.
			When thou shalt make his soul an offering for sin,
			he shall see his seed.
			He shall prolong his days,
			and the pleasure of the Lord shall prosper in his hand.

		% ---
		\paragraph{Mosiah 14:11} % *MOSIAH-14-11 
		% ---
			\label{MOSIAH-14-11}

			He shall see of the travail of his soul and shall be satisfied.
			By his knowledge shall my righteous servant justify many,
			for he shall bear their iniquities.

		% ---
		\paragraph{Mosiah 14:12} % *MOSIAH-14-12 
		% ---
			\label{MOSIAH-14-12}

			Therefore will I divide him a portion with the great,
			and he shall divide the spoil with the strong---
			because he hath poured out his soul unto death
			and he was numbered with the transgressors.
			And he bare the sins of many
			and made intercession for the transgressors.
			
	% -----
	\subsubsection{Mosiah 15} % *MOSIAH-15
	% -----
		\label{MOSIAH-15}
		
		% ---
		\paragraph{Mosiah 15:1} % *MOSIAH-15-1 
		% ---
			\label{MOSIAH-15-1}

			And now Abinadi saith unto them:
			I would that ye should understand
			that God himself shall come down among the children of men
			and shall redeem his people.

		% ---
		\paragraph{Mosiah 15:2} % *MOSIAH-15-2 
		% ---
			\label{MOSIAH-15-2}

			And because he dwelleth
				\footnote{%
					% \label{}%
					See ATV 1335-1336 for a discussion of the missing determiner ``the'' in ``in flesh.'' The argument for not adding it back in, which there are good reasons for, is based in \hyperref[JOHN-1-14]{John 1:14}, but ATV concludes, ``Consequently, the critical text will maintain the phrase ``in flesh'', but with the understanding that it could well have originally read ``in \textbf{the} flesh''...this difficult but not impossible reading may have originally read ``he dwelleth in \textbf{the} flesh''.''
					}%
			in flesh,
			he shall be called the Son of God;
			and having subjected the flesh to the will of the Father,
			being the Father and the Son,

		% ---
		\paragraph{Mosiah 15:3} % *MOSIAH-15-3 
		% ---
			\label{MOSIAH-15-3}

			the Father because he was conceived by the power of God
			and the Son because of the flesh,
			thus becoming the Father and
				\footnote{%
					% \label{}%
					See ATV PDF 1337 for another discussion, similar to the one about the previous verse, on the missing determiner ``the'' in ``Father and Son.''
					}%
			Son%
				\footnote{%
					% \label{}%
					Compare \hyperref[DC93-4]{DC 93:4}, which \textit{might} be an explanation of the same thing. See also \hyperref[DC93-12]{DC 93:12--14}, especially 14.
					}

		% ---
		\paragraph{Mosiah 15:4} % *MOSIAH-15-4 
		% ---
			\label{MOSIAH-15-4}

			---and they are one God,
			yea, the very Eternal Father of heaven and of earth---

		% ---
		\paragraph{Mosiah 15:5} % *MOSIAH-15-5 
		% ---
			\label{MOSIAH-15-5}

			and thus the flesh becoming subject to the spirit,%
				\footnote{%
					\label{MOSIAH-15-5-spirit}%
					This word is not capitalized in the Yale version, but is in the current vresion of the BOM. I think this is important and that it should not be capitalized---I don't think it refers to the Holy Spirit or Holy Ghost, but instead to whatever it is that makes Christ have the powers of the Father while dwelling on Earth. It refers to whatever makes a person have a ``spiritual body'' after the resurrection (\hyperref[DC88-27]{DC 88:27}).
					}
			or the Son to the Father, being one God,
			suffereth temptation and yieldeth not to the temptation,
			but suffereth himself to be mocked and scourged
			and cast out and disowned by his people.

		% % ---
		\paragraph{Mosiah 15:6} % *MOSIAH-15-6 
		% ---
			\label{MOSIAH-15-6}

			And after all this
			and after working many mighty miracles among the children of men,
			he shall be led
			---yea, even as Isaiah said,
			as a sheep before the shearer is dumb,
			so he opened not his mouth---

		% ---
		\paragraph{Mosiah 15:7} % *MOSIAH-15-7 
		% ---
			\label{MOSIAH-15-7}
			
			yea, even so he shall be led, crucified, and slain,
			the flesh becoming subject even
				\footnote{%
					% \label{}%
					For ``unto death,'' see \hyperref[PHILIPPIANS-2-8]{Philippians 2:8},
					}%
			unto death,
			the will of the Son being swallowed up in the will%
				\footnote{%
					% \label{MOSIAH-15-7-will}%
					See (\hyperref[3NEPHI-27-13]{3 Nephi 27:13}).
					}
			of the Father.

		% ---
		\paragraph{Mosiah 15:8} % *MOSIAH-15-8 
		% ---
			\label{MOSIAH-15-8}

			And thus God breaketh the bands of death,
			having gained the victory over death,
			giving the Son power to make intercession for the children of men,

		% ---
		\paragraph{Mosiah 15:9} % *MOSIAH-15-9 
		% ---
			\label{MOSIAH-15-9}

			having ascended into heaven,
			having the bowels of mercy,
			being filled with compassion toward the children of men,
			standing betwixt them and justice,
			having broken the bands of death,
			having taken upon himself their iniquity and their transgressions,
			having redeemed them and satisfied the demands of justice.

		% ---
		\paragraph{Mosiah 15:10} % *MOSIAH-15-10 
		% ---
			\label{MOSIAH-15-10}

			And now I say unto you:
			Who shall declare his generation?
			Behold, I say unto you that
			when his soul has been made an offering for sin,
			he shall see his seed.
			And now, what say ye?
			And who shall be his seed?

		% ---
		\paragraph{Mosiah 15:11} % *MOSIAH-15-11 
		% ---
			\label{MOSIAH-15-11}

			Behold, I say unto you that
			whosoever hath heard the words of the prophets,
			yea, all the holy prophets which have prophesied
			concerning the coming of the Lord,
			I say unto you that all those who hath hearkened unto their words
			and believed that the Lord would redeem his people
			and have looked forward to that day for a remission of their sins,
			I say unto you that these are his seed---
			or they are heirs of the kingdom of God.

		% ---
		\paragraph{Mosiah 15:12} % *MOSIAH-15-12 
		% ---
			\label{MOSIAH-15-12}

			For these are they whose sins he hath borne;
			these are they for whom he hath died,
			to redeem them from their transgressions.
			And now, are they not his seed?

		% ---
		\paragraph{Mosiah 15:13} % *MOSIAH-15-13 
		% ---
			\label{MOSIAH-15-13}

			Yea, and are not the prophets,
			every one that has opened his mouth to prophesy
			that has not fallen into transgression
			---I mean all the holy prophets ever since the world began---
			I say unto you that they are his seed.

		% ---
		\paragraph{Mosiah 15:14} % *MOSIAH-15-14 
		% ---
			\label{MOSIAH-15-14}

			And these are they which hath published peace,
			that hath brought good tidings of good,
			that hath published salvation,
			that saith unto Zion:
			Thy God reigneth!

		% ---
		\paragraph{Mosiah 15:15} % *MOSIAH-15-15 
		% ---
			\label{MOSIAH-15-15}

			And O how beautiful upon the mountains were their feet!

		% ---
		\paragraph{Mosiah 15:16} % *MOSIAH-15-16 
		% ---
			\label{MOSIAH-15-16}

			And again, how beautiful upon the mountains are the feet of those
			that are still publishing peace!

		% ---
		\paragraph{Mosiah 15:17} % *MOSIAH-15-17 
		% ---
			\label{MOSIAH-15-17}

			And again, how beautiful upon the mountains are the feet of those
			who shall hereafter publish peace,
			yea, from this time henceforth and forever!

		% ---
		\paragraph{Mosiah 15:18} % *MOSIAH-15-18 
		% ---
			\label{MOSIAH-15-18}

			And behold, I say unto you:
			This is not all,
			for O how beautiful upon the mountains are the feet of him
			that bringeth good tidings,
			that is the founder of peace,
			yea, even the Lord who hath redeemed his people,
			yea, him who hath granted salvation unto his people.

		% ---
		\paragraph{Mosiah 15:19} % *MOSIAH-15-19 
		% ---
			\label{MOSIAH-15-19}

			For were it not for the redemption which he hath made for his people,
			which was prepared from the foundation of the world,
			I say unto you
			---were it not for this---
			that all mankind must have perished.

		% ---
		\paragraph{Mosiah 15:20} % *MOSIAH-15-20 
		% ---
			\label{MOSIAH-15-20}

			But behold, the bands of death shall be broken.
			And the Son reigneth and hath power over the dead;
			therefore he bringeth to pass the resurrection of the dead.

		% ---
		\paragraph{Mosiah 15:21} % *MOSIAH-15-21 
		% ---
			\label{MOSIAH-15-21}

			And there cometh a resurrection,
			even a first resurrection,
			yea, even a resurrection of those
			that have been and which are and which shall be,
			even until the resurrection of Christ,
			for so shall he be called.

		% ---
		\paragraph{Mosiah 15:22} % *MOSIAH-15-22 
		% ---
			\label{MOSIAH-15-22}

			And now the resurrection of all the prophets
			and all those that have believed in their words
			---or all those that have kept the commandments of God---
			these shall come forth in the first resurrection;
			therefore they are the first resurrection.

		% ---
		\paragraph{Mosiah 15:23} % *MOSIAH-15-23 
		% ---
			\label{MOSIAH-15-23}

			They are raised to dwell with God, who hath redeemed them.
			Thus they have eternal life through Christ,
			who hath broken the bands of death.

		% ---
		\paragraph{Mosiah 15:24} % *MOSIAH-15-24 
		% ---
			\label{MOSIAH-15-24}

			And there are those who have part in the first resurrection,
			and these are they that have died before Christ came,
			in their ignorance,
			not having salvation declared unto them.
			And thus the Lord bringeth about the restoration of these,
			and they have a part in the first resurrection or hath eternal life,
			being redeemed by the Lord.

		% ---
		\paragraph{Mosiah 15:25} % *MOSIAH-15-25 
		% ---
			\label{MOSIAH-15-25}

			And little children also hath eternal life.

		% ---
		\paragraph{Mosiah 15:26} % *MOSIAH-15-26 
		% ---
			\label{MOSIAH-15-26}

			But behold and fear and tremble before God
			---for ye had ought to tremble---
			for the Lord redeemeth none such
			that rebelleth against him and dieth in their sins
			---yea, even all those that have perished in their sins,
			ever since the world began---
			that have willfully rebelled against God,
			that have known the commandments of God and would not keep them.
			These are they that have no part in the first resurrection.

		% ---
		\paragraph{Mosiah 15:27} % *MOSIAH-15-27 
		% ---
			\label{MOSIAH-15-27}

			Therefore had ye not ought to tremble?
			For salvation cometh to none such,
			for the Lord hath redeemed none such.
			Yea, neither can the Lord redeem such,
			for he cannot deny himself;
			for he cannot deny justice when it hath its claim.

		% ---
		\paragraph{Mosiah 15:28} % *MOSIAH-15-28 
		% ---
			\label{MOSIAH-15-28}

			And now I say unto you that the time shall come
			that the salvation of the Lord shall be declared
			to every nation, kindred, tongue, and people.

		% ---
		\paragraph{Mosiah 15:29} % *MOSIAH-15-29 
		% ---
			\label{MOSIAH-15-29}

			Yea, Lord, thy watchmen shall lift up their voice;
			with the voice together shall they sing.
			For they shall see eye to eye when the Lord shall bring again Zion.

		% ---
		\paragraph{Mosiah 15:30} % *MOSIAH-15-30 
		% ---
			\label{MOSIAH-15-30}

			Break forth into joy!
			Sing together, ye waste places of Jerusalem!
			For the Lord hath comforted his people;
			he hath redeemed Jerusalem.

		% ---
		\paragraph{Mosiah 15:31} % *MOSIAH-15-31 
		% ---
			\label{MOSIAH-15-31}

			The Lord hath made bare his holy arm in the eyes of all the nations,
			and all the ends of the earth shall see the salvation of our God.
			
	% -----
	\subsubsection{Mosiah 16} % *MOSIAH-16 
	% -----
		\label{MOSIAH-16}
		
		% ---
		\paragraph{Mosiah 16:1} % *MOSIAH-16-1 
		% ---
			\label{MOSIAH-16-1}

			And now it came to pass that after Abinadi had spoken these words,
			he stretched forth his hands and said:
			The time shall come when all shall see the salvation of the Lord,
			when every nation, kindred, tongue, and people shall see eye to eye
			and shall confess before God that his judgments are just.

		% ---
		\paragraph{Mosiah 16:2} % *MOSIAH-16-2 
		% ---
			\label{MOSIAH-16-2}

			And then shall the wicked be cast out,
			and they shall have cause to howl and weep and wail and gnash their teeth---
			and this because they would not hearken unto the voice of the Lord.
			Therefore the Lord redeemeth them not.

		% ---
		\paragraph{Mosiah 16:3} % *MOSIAH-16-3 
		% ---
			\label{MOSIAH-16-3}

			For they are carnal and devilish;
			and the devil hath power over them,
			yea, even that old serpent that did beguile our first parents,
			which was the cause of their fall,
			which was the cause of all mankind's becoming carnal, sensual, devilish,
			knowing evil from good,
			subjecting themselves to the devil.

		% ---
		\paragraph{Mosiah 16:4} % *MOSIAH-16-4 
		% ---
			\label{MOSIAH-16-4}

			Thus all mankind were lost.
			And behold, they would have been endlessly lost
			were it not that God redeemed his people from their lost and fallen state.

		% ---
		\paragraph{Mosiah 16:5} % *MOSIAH-16-5 
		% ---
			\label{MOSIAH-16-5}

			But remember that he that persists in his own carnal nature
			and goes on in the ways of sin and rebellion against God,
			he remaineth in his fallen state,
			and the devil hath all power over him.
			Therefore he is as though there was no redemption made,
			being an enemy to God;
			and also is the devil an enemy to God.

		% ---
		\paragraph{Mosiah 16:6} % *MOSIAH-16-6 
		% ---
			\label{MOSIAH-16-6}

			And now if Christ had not come into the world
			---speaking of things to come as though they had already come---
			there could have been no redemption.

		% ---
		\paragraph{Mosiah 16:7} % *MOSIAH-16-7 
		% ---
			\label{MOSIAH-16-7}

			And if Christ had not risen from the dead
			or broken the bands of death
			---that the grave should have no victory
			and that death should have no sting---
			there could have been no resurrection.

		% ---
		\paragraph{Mosiah 16:8} % *MOSIAH-16-8 
		% ---
			\label{MOSIAH-16-8}

			But there is a resurrection.
			Therefore the grave hath no victory,
			and the sting of death is swallowed up in Christ.

		% ---
		\paragraph{Mosiah 16:9} % *MOSIAH-16-9 
		% ---
			\label{MOSIAH-16-9}

			He is the light and the life of the world,
			yea, a light that is endless that can never be darkened;
			yea, and also a life which is endless, that there can be no more death.

		% ---
		\paragraph{Mosiah 16:10} % *MOSIAH-16-10 
		% ---
			\label{MOSIAH-16-10}

			Even this mortal shall put on immortality,
			and this corruption shall put on incorruption
			and shall be brought to stand before the bar of God
			to be judged of him according to their works,
			whether they be good or whether they be evil:

		% ---
		\paragraph{Mosiah 16:11} % *MOSIAH-16-11 
		% ---
			\label{MOSIAH-16-11}

			if they be good, to the resurrection of endless life and happiness;
			and if they be evil, to the resurrection of endless damnation,
			being delivered up to the devil,
			who hath subjected them---which is damnation---

		% ---
		\paragraph{Mosiah 16:12} % *MOSIAH-16-12 
		% ---
			\label{MOSIAH-16-12}

			having gone according to their own carnal wills and desires,
			having never called upon the Lord
			while the arms of mercy was extended towards them
			---for the arms of mercy was extended towards them and they would not---
			they being warned of their iniquities,
			and yet they would not depart from them.
			And they were commanded to repent,
			and yet they would not repent.

		% ---
		\paragraph{Mosiah 16:13} % *MOSIAH-16-13 
		% ---
			\label{MOSIAH-16-13}

			And now had ye not ought to tremble and repent of your sins?
			---and remember, only in and through Christ ye can be saved.

		% ---
		\paragraph{Mosiah 16:14} % *MOSIAH-16-14 
		% ---
			\label{MOSIAH-16-14}

			Therefore if ye teach the law of Moses,
			also teach that it is a shadow of those things which are to come.

		% ---
		\paragraph{Mosiah 16:15} % *MOSIAH-16-15 
		% ---
			\label{MOSIAH-16-15}

			Teach them that redemption cometh through Christ the Lord,
			which is the very Eternal Father.
			Amen.

	% -----
	\subsubsection{Mosiah 17} % *MOSIAH-17 
	% -----
		\label{MOSIAH-17}
		
		% ---
		\paragraph{Mosiah 17:1} % *MOSIAH-17-1 
		% ---
			\label{MOSIAH-17-1}

			And now it came to pass that when Abinadi had finished these sayings
			that the king commanded that the priests should take him
			and cause that he should be put to death.

		% ---
		\paragraph{Mosiah 17:2} % *MOSIAH-17-2 
		% ---
			\label{MOSIAH-17-2}

			But there was one among them whose name was Alma,
			he also being a descendant of Nephi,
			and he was a young man.
			And he believed the words which Abinadi had spoken,
			for he knew concerning the iniquity which Abinadi had testified against them.
			Therefore he began to plead with the king
			that he would not be angry with Abinadi
			but suffer that he might depart in peace.

		% ---
		\paragraph{Mosiah 17:3} % *MOSIAH-17-3 
		% ---
			\label{MOSIAH-17-3}

			But the king was more wroth and caused
			that Alma should be cast out from among them
			and sent his servants after him that they might slay him.

		% ---
		\paragraph{Mosiah 17:4} % *MOSIAH-17-4 
		% ---
			\label{MOSIAH-17-4}

			But he fled from before them and hid himself that they found him not.
			And he being concealed for many days
			did write all the words which Abinadi had spoken.

		% ---
		\paragraph{Mosiah 17:5} % *MOSIAH-17-5 
		% ---
			\label{MOSIAH-17-5}

			And it came to pass that the king caused
			that his guards should surround Abinadi and take him;
			and they bound him and cast him into prison.

		% ---
		\paragraph{Mosiah 17:6} % *MOSIAH-17-6 
		% ---
			\label{MOSIAH-17-6}

			And after three days, having counseled with his priests,
			he caused that he should again be brought before him.

		% ---
		\paragraph{Mosiah 17:7} % *MOSIAH-17-7 
		% ---
			\label{MOSIAH-17-7}

			And he said unto him:
			Abinadi, we have found an accusation against thee,
			and thou art worthy of death.

		% ---
		\paragraph{Mosiah 17:8} % *MOSIAH-17-8 
		% ---
			\label{MOSIAH-17-8}

			For thou hast said that God himself should come down among the children of men.
			And now for this cause thou shalt be put to death
			unless thou wilt recall all the words
			which thou hast spoken evil concerning me and my people.

		% ---
		\paragraph{Mosiah 17:9} % *MOSIAH-17-9 
		% ---
			\label{MOSIAH-17-9}

			Now Abinadi saith unto him:
			I say unto you:
			I will not recall the words
			which I have spoken unto you concerning this people,
			for they are true.
			And that ye may know of their surety,
			I have suffered myself that I have fallen into your hands.

		% ---
		\paragraph{Mosiah 17:10} % *MOSIAH-17-10 
		% ---
			\label{MOSIAH-17-10}

			Yea, and I will suffer even unto death.
			And I will not recall my words,
			and they shall stand as a testimony against you.
			And if ye slay me, ye will shed innocent blood.
			And this shall also stand as a testimony against you at the last day.

		% ---
		\paragraph{Mosiah 17:11} % *MOSIAH-17-11 
		% ---
			\label{MOSIAH-17-11}

			And now king Noah was about to release him, for he feared his word;
			for he feared that the judgments of God would come upon him.

		% ---
		\paragraph{Mosiah 17:12} % *MOSIAH-17-12 
		% ---
			\label{MOSIAH-17-12}

			But the priests lifted up their voices against him
			and began to accuse him, saying:
			He hath reviled the king!
			Therefore the king was stirred up in anger against him,
			and he delivered him up that he might be slain.

		% ---
		\paragraph{Mosiah 17:13} % *MOSIAH-17-13 
		% ---
			\label{MOSIAH-17-13}

			And it came to pass that they took him and bound him
			and scorched his skin with fagots,
			yea, even unto death.

		% ---
		\paragraph{Mosiah 17:14} % *MOSIAH-17-14 
		% ---
			\label{MOSIAH-17-14}

			And now when the flames began to scorch him,
			he cried unto them, saying:

		% ---
		\paragraph{Mosiah 17:15} % *MOSIAH-17-15 
		% ---
			\label{MOSIAH-17-15}

			Behold, even as ye have done unto me,
			so shall it come to pass that
			thy seed shall cause that many shall suffer the pains that I do suffer,
			even the pains of death by fire---
			and this because they believe in the salvation of the Lord their God.

		% ---
		\paragraph{Mosiah 17:16} % *MOSIAH-17-16 
		% ---
			\label{MOSIAH-17-16}

			And it will come to pass that
			ye shall be afflicted with all manner of diseases because of your iniquities.

		% ---
		\paragraph{Mosiah 17:17} % *MOSIAH-17-17 
		% ---
			\label{MOSIAH-17-17}

			Yea, and ye shall be smitten on every hand
			and shall be driven and scattered to and fro,
			even as a wild flock is driven by wild and ferocious beasts.

		% ---
		\paragraph{Mosiah 17:18} % *MOSIAH-17-18 
		% ---
			\label{MOSIAH-17-18}

			And in that day ye shall be hunted,
			and ye shall be taken by the hand of your enemies.
			And then ye shall suffer as I suffer the pains of death by fire.

		% ---
		\paragraph{Mosiah 17:19} % *MOSIAH-17-19 
		% ---
			\label{MOSIAH-17-19}

			Thus God executeth vengeance upon those that destroy his people.
			O God, receive my soul!

		% ---
		\paragraph{Mosiah 17:20} % *MOSIAH-17-20 
		% ---
			\label{MOSIAH-17-20}

			And now when Abinadi had said these words,
			he fell, having suffered death by fire,
			yea, having been put to death
			because he would not deny the commandments of God,
			having sealed the truth of his words by his death.
			
	% -----
	\subsubsection{Mosiah 18} % *MOSIAH-18
	% -----
		\label{MOSIAH-18}
		
		% ---
		\paragraph{Mosiah 18:1} % *MOSIAH-18-1 
		% ---
			\label{MOSIAH-18-1}

			And now it came to pass that
			Alma, who had fled from the servants of king Noah,
			repented of his sins and iniquities
			and went about privately among the people
			and began to teach the words of Abinadi,

		% ---
		\paragraph{Mosiah 18:2} % *MOSIAH-18-2 
		% ---
			\label{MOSIAH-18-2}

			yea, concerning that which was to come,
			and also concerning the resurrection of the dead,
			and the redemption of the people which was to be brought to pass
			through the power and sufferings and death of Christ
			and his resurrection and ascension into heaven.

		% ---
		\paragraph{Mosiah 18:3} % *MOSIAH-18-3 
		% ---
			\label{MOSIAH-18-3}

			And as many as would hear his word he did teach;
			and he taught them privately,
			that it might not come to the knowledge of the king.
			And many did believe his words.

		% ---
		\paragraph{Mosiah 18:4} % *MOSIAH-18-4 
		% ---
			\label{MOSIAH-18-4}

			And it came to pass that
			as many as did believe him did go forth to a place which was called Mormon,
			having received its name from the king,
			being in the borders of the land,
			having been infested by times or at seasons by wild beasts.%
				\footnote{%
					\label{MOSIAH-18-4-NOTE-mormon}%
					This is the first occurrence of the name ``Mormon'' in the BOM (there are three occurrences in Words of Mormon, but those happen later in time). In 1843, JS wrote a letter to the editor of the TS which argued that the name ``Mormon'' means ``more good'' (see \hyperref[EMS-1843-JS-CIR-20-MAY-1843]{here}). His argument seems to be that the word \textit{mon} in Egyptian (he doesn't say what era of Egyptian, but he mentions the ``reformed Egyptian'' thing from the BOM) means ``good,'' and then you add the English ``more,'' contracted to ``mor,'' and you get ``Mormon,'' for ``more good.'' This is a bit baffling---he's combining an alleged ancient word with a contemporary English one to give the meaning of a name that appeared on the gold plates. It's a pretty suspect methodology. And there's another reason to think it's wrong. The place names in the BOM often have significance in describing the place they are names of, and we read in this verse and around here that the place of Mormon had a ``fountain of pure water'' but also may have had infestations of wild animals. It doesn't seem like the kind of place that would be named ``more good.''
					
					There is a long discussion of the name, and of the letter and its authorship, in the \href{https://onoma.lib.byu.edu/index.php/MORMON}{onomasticon}. They point out that JS may not have been the author of the letter, although it's certainly the kind of linguistic connection he loved to make.
					
					Hinckley has an interesting conference talk from October 1990 on this topic, called ``Mormon Should Mean `More Good'.'' There he also discusses the difference between the nickname ``Mormon'' and the actual name of the church, which Nelson picked up on later.
					}

		% ---
		\paragraph{Mosiah 18:5} % *MOSIAH-18-5 
		% ---
			\label{MOSIAH-18-5}

			Now there was in Mormon a fountain of pure water;
			and Alma resorted thither,
			there being near the water a thicket of small trees
			where he did hide himself in the daytime from the searches of the king.

		% ---
		\paragraph{Mosiah 18:6} % *MOSIAH-18-6 
		% ---
			\label{MOSIAH-18-6}

			And it came to pass that
			as many as believed him went thither to hear his words.

		% ---
		\paragraph{Mosiah 18:7} % *MOSIAH-18-7 
		% ---
			\label{MOSIAH-18-7}

			And it came to pass after many days
			there were a goodly number gathered together
			to the place of Mormon to hear the words of Alma;
			yea, all were gathered together that believed on his word to hear him.
			And he did teach them and did preach unto them
			repentance and redemption and faith on the Lord.

		% ---
		\paragraph{Mosiah 18:8} % *MOSIAH-18-8 
		% ---
			\label{MOSIAH-18-8}

			And it came to pass that he said unto them:
			Behold, here is the waters of Mormon,
			for thus were they called.
			And now as ye are desirous to come into the fold of God
			and to be called his people
			and are willing to bear one another's burdens,
			that they may be light,

		% ---
		\paragraph{Mosiah 18:9} % *MOSIAH-18-9 
		% ---
			\label{MOSIAH-18-9}

			yea, and are willing to mourn with those that mourn,%
				\footnote{%
					% \label{}%
					See \hyperref[DC30-6]{DC 30:6}, an item of instruction given to Peter Whitmer: ``And be you afflicted in all his afflictions, ever lifting up your heart unto me in prayer and faith, for his and your deliverance.'' See also \hyperref[ROMANS-12-15]{Romans 12:15}, \hyperref[2CORINTHIANS-1-4]{2 Corinthians 1:4}, 
					}
			yea, and comfort those that stand in need of comfort,
			and to stand as witnesses of God at all times and in all things
			and in all places that ye may be in,
			even until death,
			that ye may be redeemed of God
			and be numbered with those of the first resurrection,
			that ye may have eternal life---

		% ---
		\paragraph{Mosiah 18:10} % *MOSIAH-18-10 
		% ---
			\label{MOSIAH-18-10}

			now I say unto you,
			if this be the desires of your hearts,
			what have you against being baptized in the name of the Lord,
			as a witness before him
			that ye have entered into a covenant with him,
			that ye will serve him and keep his commandments,
			that he may pour out his Spirit more abundantly upon you?

		% ---
		\paragraph{Mosiah 18:11} % *MOSIAH-18-11 
		% ---
			\label{MOSIAH-18-11}

			And now when the people had heard these words,
			they clapped their hands for joy and exclaimed:
			This is the desires of our hearts!

		% ---
		\paragraph{Mosiah 18:12} % *MOSIAH-18-12 
		% ---
			\label{MOSIAH-18-12}

			And now it came to pass that Alma took Helam,
			he being one of the first,
			and went and stood forth in the water and cried, saying:
			O Lord, pour out thy Spirit upon thy servant,
			that he may do this work with holiness of heart.

		% ---
		\paragraph{Mosiah 18:13} % *MOSIAH-18-13 
		% ---
			\label{MOSIAH-18-13}

			And when he had said these words,
			the Spirit of the Lord was upon him and he said:
			Helam, I baptize thee,
			having authority from the Almighty God,
			as a testimony that ye have entered into a covenant to serve him
			until you are dead as to the mortal body;
			and may the Spirit of the Lord be poured out upon you,
			and may he grant unto you eternal life through the redemption of Christ,
			which he hath prepared from the foundation of the world.

		% ---
		\paragraph{Mosiah 18:14} % *MOSIAH-18-14 
		% ---
			\label{MOSIAH-18-14}

			And after Alma had said these words,
			both Alma and Helam was buried in the water.
			And they arose and came forth out of the water rejoicing,
			being filled with the Spirit.

		% ---
		\paragraph{Mosiah 18:15} % *MOSIAH-18-15 
		% ---
			\label{MOSIAH-18-15}

			And again Alma took another
			and went forth a second time into the water
			and baptized him according to the first,
			only he did not bury himself again in the water.

		% ---
		\paragraph{Mosiah 18:16} % *MOSIAH-18-16 
		% ---
			\label{MOSIAH-18-16}

			And after this manner he did baptize every one
			that went forth to the place of Mormon
			---and they were in number about two hundred and four souls---
			yea, and they were baptized in the waters of Mormon
			and were filled with the grace of God.

		% ---
		\paragraph{Mosiah 18:17} % *MOSIAH-18-17 
		% ---
			\label{MOSIAH-18-17}

			And they were called the church of God, or the church of Christ,
			from that time forward.
			And it came to pass that
			whosoever was baptized by the power and authority of God,
			they was added to his church.

		% ---
		\paragraph{Mosiah 18:18} % *MOSIAH-18-18 
		% ---
			\label{MOSIAH-18-18}

			And it came to pass that Alma having authority from God ordained priests,
			even one priest to every fifty of their number did he ordain to preach unto them
			and to teach them concerning the things pertaining to the kingdom of God.

		% ---
		\paragraph{Mosiah 18:19} % *MOSIAH-18-19 
		% ---
			\label{MOSIAH-18-19}

			And he commanded them that they should teach nothing
			save it were the things which he had taught
			and which had been spoken by the mouth of the holy prophets.

		% ---
		\paragraph{Mosiah 18:20} % *MOSIAH-18-20 
		% ---
			\label{MOSIAH-18-20}

			Yea, even he commanded them that they should preach nothing
			save it were repentance and faith on the Lord, who hath redeemed his people.

		% ---
		\paragraph{Mosiah 18:21} % *MOSIAH-18-21 
		% ---
			\label{MOSIAH-18-21}

			And he commanded them that there should be no contention one with another,
			but that they should look forward with one eye,
			having one faith and one baptism,
			having their hearts knit together in unity and in love one towards another.

		% ---
		\paragraph{Mosiah 18:22} % *MOSIAH-18-22 
		% ---
			\label{MOSIAH-18-22}

			And thus he commanded them to preach.
			And thus they became the children of God.

		% ---
		\paragraph{Mosiah 18:23} % *MOSIAH-18-23 
		% ---
			\label{MOSIAH-18-23}

			And he commanded them
			that they should observe the sabbath day and keep it holy,
			and also every day they should give thanks to the Lord their God.

		% ---
		\paragraph{Mosiah 18:24} % *MOSIAH-18-24 
		% ---
			\label{MOSIAH-18-24}

			And he also commanded them that the priests which he had ordained
			should labor with their own hands for their support.

		% ---
		\paragraph{Mosiah 18:25} % *MOSIAH-18-25 
		% ---
			\label{MOSIAH-18-25}

			And there was one day in every week that was set apart
			that they should gather themselves together to teach the people
			and to worship the Lord their God,
			and also as often as it was in their power to assemble themselves together.

		% ---
		\paragraph{Mosiah 18:26} % *MOSIAH-18-26 
		% ---
			\label{MOSIAH-18-26}

			And the priests was not to depend upon the people for their support,
			but for their labor they were to receive the grace of God,
			that they might wax strong in the Spirit,
			having the knowledge of God,
			that they might teach with power and authority from God.

		% ---
		\paragraph{Mosiah 18:27} % *MOSIAH-18-27 
		% ---
			\label{MOSIAH-18-27}

			And again Alma commanded
			that the people of the church should impart of their substance,
			every one according to that which he hath:
			if he have more abundantly, he should impart more abundantly;
			and he that hath but little, but little should be required;
			and to him that hath not should be given.

		% ---
		\paragraph{Mosiah 18:28} % *MOSIAH-18-28 
		% ---
			\label{MOSIAH-18-28}

			And thus they should impart of their substance
			of their own free will and good desires towards God
			to those priests that stood in need,
			yea, and to every needy, naked soul.

		% ---
		\paragraph{Mosiah 18:29} % *MOSIAH-18-29 
		% ---
			\label{MOSIAH-18-29}

			And this he said unto them,
			having been commanded of God.
			And they did walk uprightly before God,
			imparting to one another both temporally and spiritually
			according to their needs and their wants.

		% ---
		\paragraph{Mosiah 18:30} % *MOSIAH-18-30 
		% ---
			\label{MOSIAH-18-30}

			And now it came to pass that all this was done in Mormon,
			yea, by the waters of Mormon,
			in the forest that was near the waters of Mormon,
			yea, the place of Mormon, the waters of Mormon, the forest of Mormon.
			How beautiful are they to the eyes of them
			who there came to the knowledge of their Redeemer!
			Yea, and how blessed are they,
			for they shall sing to his praise forever.

		% ---
		\paragraph{Mosiah 18:31} % *MOSIAH-18-31 
		% ---
			\label{MOSIAH-18-31}

			And these things were done in the borders of the land
			that they might not come to the knowledge of the king.

		% ---
		\paragraph{Mosiah 18:32} % *MOSIAH-18-32 
		% ---
			\label{MOSIAH-18-32}

			But behold, it came to pass that
			the king having discovered a movement among the people
			sent his servants to watch them.
			Therefore on the day that they were assembling themselves together
			to hear the word of the Lord,
			they were discovered unto the king.

		% ---
		\paragraph{Mosiah 18:33} % *MOSIAH-18-33 
		% ---
			\label{MOSIAH-18-33}

			And now the king saith
			that Alma was a stirring up the people to a rebellion against him.
			Therefore he sent his army to destroy them.

		% ---
		\paragraph{Mosiah 18:34} % *MOSIAH-18-34 
		% ---
			\label{MOSIAH-18-34}

			And it came to pass that
			Alma and the people of the Lord were apprised of the coming of the king's army.
			Therefore they took their tents and their families and departed into the wilderness.

		% ---
		\paragraph{Mosiah 18:35} % *MOSIAH-18-35 
		% ---
			\label{MOSIAH-18-35}

			And they were in number about four hundred and fifty souls.
			
	% -----
	\subsubsection{Mosiah 19} % *MOSIAH-19 
	% -----
		\label{MOSIAH-19}
		
		% ---
		\paragraph{Mosiah 19:1} % *MOSIAH-19-1 
		% ---
			\label{MOSIAH-19-1}

			And it came to pass that the army of the king returned,
			having searched in vain for the people of the Lord.

		% ---
		\paragraph{Mosiah 19:2} % *MOSIAH-19-2 
		% ---
			\label{MOSIAH-19-2}

			And now behold, the forces of the king were small, having been reduced.
			And there began to be a division among the remainder of the people.

		% ---
		\paragraph{Mosiah 19:3} % *MOSIAH-19-3 
		% ---
			\label{MOSIAH-19-3}

			And the lesser part began to breathe out threatenings against the king,
			and there began to be a great contention among them.

		% ---
		\paragraph{Mosiah 19:4} % *MOSIAH-19-4 
		% ---
			\label{MOSIAH-19-4}

			And now there was a man among them whose name was Gideon.
			And he being a strong man and an enemy to the king,
			therefore he drew his sword and swore in his wrath
			that he would slay the king.

		% ---
		\paragraph{Mosiah 19:5} % *MOSIAH-19-5 
		% ---
			\label{MOSIAH-19-5}

			And it came to pass that he fought with the king.
			And when the king saw that he was about to overpower him,
			he fled and ran and got upon the tower which was near the temple.

		% ---
		\paragraph{Mosiah 19:6} % *MOSIAH-19-6 
		% ---
			\label{MOSIAH-19-6}

			And Gideon pursued after him
			and was about to get upon the tower to slay the king.
			And the king cast his eyes round about towards the land of Shemlon---
			and behold, the army of the Lamanites were within the borders of the land!

		% ---
		\paragraph{Mosiah 19:7} % *MOSIAH-19-7 
		% ---
			\label{MOSIAH-19-7}

			And now the king cried out in the anguish of his soul, saying:
			Gideon, spare me!
			For the Lamanites are upon us and they will destroy them---
			yea, they will destroy my people.

		% ---
		\paragraph{Mosiah 19:8} % *MOSIAH-19-8 
		% ---
			\label{MOSIAH-19-8}

			And now the king was not so much concerned about his people
			as he was about his own life.
			Nevertheless Gideon did spare his life.

		% ---
		\paragraph{Mosiah 19:9} % *MOSIAH-19-9 
		% ---
			\label{MOSIAH-19-9}

			And the king commanded the people
			that they should flee before the Lamanites,
			and he himself did go before them.
			And they did flee into the wilderness with their women and their children.

		% ---
		\paragraph{Mosiah 19:10} % *MOSIAH-19-10 
		% ---
			\label{MOSIAH-19-10}

			And it came to pass that the Lamanites did pursue them
			and did overtake them and began to slay them.

		% ---
		\paragraph{Mosiah 19:11} % *MOSIAH-19-11 
		% ---
			\label{MOSIAH-19-11}

			Now it came to pass that the king commanded them
			that all the men should leave their wives and their children
			and flee before the Lamanites.

		% ---
		\paragraph{Mosiah 19:12} % *MOSIAH-19-12 
		% ---
			\label{MOSIAH-19-12}

			Now there were many that would not leave them
			but had rather stay and perish with them.
			And the rest left their wives and their children and fled.

		% ---
		\paragraph{Mosiah 19:13} % *MOSIAH-19-13 
		% ---
			\label{MOSIAH-19-13}

			And it came to pass that
			those that tarried with their wives and their children caused
			that their fair daughters should stand forth and plead with the Lamanites
			that they would not slay them.

		% ---
		\paragraph{Mosiah 19:14} % *MOSIAH-19-14 
		% ---
			\label{MOSIAH-19-14}

			And it came to pass that the Lamanites had compassion on them,
			for they were charmed with the beauty of their women.

		% ---
		\paragraph{Mosiah 19:15} % *MOSIAH-19-15 
		% ---
			\label{MOSIAH-19-15}

			Therefore the Lamanites did spare their lives and took them captives
			and carried them back to the land of Nephi and granted unto them
			that they might possess the land under the conditions
			that they would deliver up the king Noah into the hands of the Lamanites
			and deliver up their property,
			even one half of all they possessed:
			one half of their gold and their silver and all their precious things.
			And thus they should pay tribute to the king of the Lamanites from year to year.

		% ---
		\paragraph{Mosiah 19:16} % *MOSIAH-19-16 
		% ---
			\label{MOSIAH-19-16}

			And now there was one of the sons of the king
			among those that was taken captive
			whose name was Limhi.

		% ---
		\paragraph{Mosiah 19:17} % *MOSIAH-19-17 
		% ---
			\label{MOSIAH-19-17}

			And now Limhi was desirous that his father should not be destroyed.
			Nevertheless Limhi was not ignorant of the iniquities of his father,
			he himself being a just man.

		% ---
		\paragraph{Mosiah 19:18} % *MOSIAH-19-18 
		% ---
			\label{MOSIAH-19-18}

			And it came to pass that Gideon sent men into the wilderness secretly
			to search for the king and those that was with him.
			And it came to pass that they met the people in the wilderness,
			all save the king and his priests.

		% ---
		\paragraph{Mosiah 19:19} % *MOSIAH-19-19 
		% ---
			\label{MOSIAH-19-19}

			Now they had sworn in their hearts
			that they would return to the land of Nephi;
			and if their wives and their children were slain,
			and also those that had tarried with them,
			that they would seek revenge and also perish with them.

		% ---
		\paragraph{Mosiah 19:20} % *MOSIAH-19-20 
		% ---
			\label{MOSIAH-19-20}

			And the king commanded them that they should not return,
			and they were angry with the king
			and caused that he should suffer even unto death by fire.

		% ---
		\paragraph{Mosiah 19:21} % *MOSIAH-19-21 
		% ---
			\label{MOSIAH-19-21}

			And they were about to take the priests also to put them to death,
			and they fled before them.

		% ---
		\paragraph{Mosiah 19:22} % *MOSIAH-19-22 
		% ---
			\label{MOSIAH-19-22}

			And it came to pass that
			they were about to return to the land of Nephi,
			and they met the men of Gideon.
			And the men of Gideon told them
			of all that had happened to their wives and their children
			and that the Lamanites had granted unto them
			that they might possess the land
			by paying a tribute to the Lamanites of one half of all they possessed.

		% ---
		\paragraph{Mosiah 19:23} % *MOSIAH-19-23 
		% ---
			\label{MOSIAH-19-23}

			And the people told the men of Gideon that they had slain the king,
			and his priests had fled from them farther into the wilderness.

		% ---
		\paragraph{Mosiah 19:24} % *MOSIAH-19-24 
		% ---
			\label{MOSIAH-19-24}

			And it came to pass that after they had ended the sermon
			that they returned to the land of Nephi,
			rejoicing because their wives and their children were not slain;
			and they told Gideon what they had done to the king.

		% ---
		\paragraph{Mosiah 19:25} % *MOSIAH-19-25 
		% ---
			\label{MOSIAH-19-25}

			And it came to pass that the king of the Lamanites made an oath unto them
			that his people should not slay them.

		% ---
		\paragraph{Mosiah 19:26} % *MOSIAH-19-26 
		% ---
			\label{MOSIAH-19-26}

			And also Limhi, being the son of the king,
			having the kingdom conferred upon him by the people,
			made oath unto the king of the Lamanites
			that his people should pay tribute unto him,
			even one half of all they possessed.

		% ---
		\paragraph{Mosiah 19:27} % *MOSIAH-19-27 
		% ---
			\label{MOSIAH-19-27}

			And it came to pass that Limhi began to establish the kingdom
			and to establish peace among his people.

		% ---
		\paragraph{Mosiah 19:28} % *MOSIAH-19-28 
		% ---
			\label{MOSIAH-19-28}

			And the king of the Lamanites set guards round about the land,
			that he might keep the people of Limhi in the land,
			that they might not depart into the wilderness.
			And he did support his guards out of the tribute
			which he did receive from the Nephites.

		% ---
		\paragraph{Mosiah 19:29} % *MOSIAH-19-29 
		% ---
			\label{MOSIAH-19-29}

			And now king Limhi did have continual peace
			in his kingdom for the space of two years,
			that the Lamanites did not molest them nor seek to destroy them.

	% -----
	\subsubsection{Mosiah 20} % *MOSIAH-20 
	% -----
		\label{MOSIAH-20}
		
		% ---
		\paragraph{Mosiah 20:1} % *MOSIAH-20-1 
		% ---
			\label{MOSIAH-20-1}

			Now there was a place in Shemlon
			where the daughters of the Lamanites did gather themselves together
			for to sing and to dance and to make themselves merry.

		% ---
		\paragraph{Mosiah 20:2} % *MOSIAH-20-2 
		% ---
			\label{MOSIAH-20-2}

			And it came to pass that there was one day
			a small number of them gathered together to sing and to dance.

		% ---
		\paragraph{Mosiah 20:3} % *MOSIAH-20-3 
		% ---
			\label{MOSIAH-20-3}

			And now the priests of king Noah being ashamed to return to the city of Nephi,
			yea, and also fearing that the people would slay them,
			therefore they durst not return to their wives and their children.

		% ---
		\paragraph{Mosiah 20:4} % *MOSIAH-20-4 
		% ---
			\label{MOSIAH-20-4}

			And having tarried in the wilderness
			and having discovered the daughters of the Lamanites,
			they laid and watched them.

		% ---
		\paragraph{Mosiah 20:5} % *MOSIAH-20-5 
		% ---
			\label{MOSIAH-20-5}

			And when there were but few of them gathered together to dance,
			they came forth out of their secret places
			and took them and carried them into the wilderness,
			yea, twenty and four of the daughters of the Lamanites
			they carried into the wilderness.

		% ---
		\paragraph{Mosiah 20:6} % *MOSIAH-20-6 
		% ---
			\label{MOSIAH-20-6}

			And it came to pass that when the Lamanites found
			that their daughters had been missing,
			they were angry with the people of Limhi;
			for they thought it was the people of Limhi.

		% ---
		\paragraph{Mosiah 20:7} % *MOSIAH-20-7 
		% ---
			\label{MOSIAH-20-7}

			Therefore they sent their armies forth
			---yea, even the king himself went before his people---
			and they went up to the land of Nephi to destroy the people of Limhi.

		% ---
		\paragraph{Mosiah 20:8} % *MOSIAH-20-8 
		% ---
			\label{MOSIAH-20-8}

			And now Limhi had discovered them from the tower,
			even all their preparations for war did he discover.
			Therefore he gathered his people together
			and laid wait for them in the fields and in the forests.

		% ---
		\paragraph{Mosiah 20:9} % *MOSIAH-20-9 
		% ---
			\label{MOSIAH-20-9}

			And it came to pass that when the Lamanites had come up
			that the people of Limhi began to fall upon them from their waiting places
			and began to slay them.

		% ---
		\paragraph{Mosiah 20:10} % *MOSIAH-20-10 
		% ---
			\label{MOSIAH-20-10}

			And it came to pass that the battle became exceeding sore,
			for they fought like lions for their prey.

		% ---
		\paragraph{Mosiah 20:11} % *MOSIAH-20-11 
		% ---
			\label{MOSIAH-20-11}

			And it came to pass that
			the people of Limhi began to drive the Lamanites before them.
			Yet they were not half so numerous as the Lamanites,
			but they fought for their lives and for their wives and for their children.
			Therefore they exerted themselves and like dragons did they fight.

		% ---
		\paragraph{Mosiah 20:12} % *MOSIAH-20-12 
		% ---
			\label{MOSIAH-20-12}

			And it came to pass that they found the king of the Lamanites
			among the number of their dead,
			yet he was not dead,
			having been wounded and left upon the ground,
			so speedy was the flight of his people.

		% ---
		\paragraph{Mosiah 20:13} % *MOSIAH-20-13 
		% ---
			\label{MOSIAH-20-13}

			And they took him and bound up his wounds
			and brought him before Limhi and said:
			Behold, here is the king of the Lamanites;
			he having received a wound hath fallen among their dead,
			and they have left him.
			And behold, we have brought him before you.
			And now let us slay him!

		% ---
		\paragraph{Mosiah 20:14} % *MOSIAH-20-14 
		% ---
			\label{MOSIAH-20-14}

			But Limhi saith unto them:
			Ye shall not slay him but bring him hither that I may see him.
			And they brought him, and Limhi saith unto him:
			What cause have ye to come up to war against my people?
			Behold, my people have not broken the oath that I made unto you.
			Therefore why should ye break the oath which ye made unto my people?

		% ---
		\paragraph{Mosiah 20:15} % *MOSIAH-20-15 
		% ---
			\label{MOSIAH-20-15}

			And now the king said:
			I have broken the oath
			because thy people did carry away the daughters of my people.
			Therefore in my anger I did cause my people
			to come up to war against thy people.

		% ---
		\paragraph{Mosiah 20:16} % *MOSIAH-20-16 
		% ---
			\label{MOSIAH-20-16}

			Now Limhi had heard nothing concerning this matter.
			Therefore he saith:
			I will search among my people;
			and whosoever hath done this thing shall perish.
			Therefore he caused a search to be made among his people.

		% ---
		\paragraph{Mosiah 20:17} % *MOSIAH-20-17 
		% ---
			\label{MOSIAH-20-17}

			Now when Gideon had heard these things,
			he being the king's captain,
			he went forth and said unto the king:
			I pray thee forbear and do not search this people;
			and lay not this thing to their charge.

		% ---
		\paragraph{Mosiah 20:18} % *MOSIAH-20-18 
		% ---
			\label{MOSIAH-20-18}

			For do ye not remember the priests of thy father,
			which this people sought to destroy?
			And are they not in the wilderness?
			And is it not they which have stolen the daughters of the Lamanites?

		% ---
		\paragraph{Mosiah 20:19} % *MOSIAH-20-19 
		% ---
			\label{MOSIAH-20-19}

			And now behold and tell the king of these things,
			that he may tell his people, that they may be pacified towards us.
			For behold, they are already preparing to come against us.
			And behold also, there are but few of us.

		% ---
		\paragraph{Mosiah 20:20} % *MOSIAH-20-20 
		% ---
			\label{MOSIAH-20-20}

			And behold, they come with their numerous hosts;
			and except the king doth pacify them towards us, we must perish.

		% ---
		\paragraph{Mosiah 20:21} % *MOSIAH-20-21 
		% ---
			\label{MOSIAH-20-21}

			For are not the words of Abinadi fulfilled which he prophesied against us---
			and all this because we would not hearken unto the word of the Lord
			and turn from our iniquities?

		% ---
		\paragraph{Mosiah 20:22} % *MOSIAH-20-22 
		% ---
			\label{MOSIAH-20-22}

			And now let us pacify the king,
			and we fulfill the oath which we have made unto him.
			For it is better that we should be in bondage
			than that we should lose our lives.
			Therefore let us put a stop to the shedding of so much blood.

		% ---
		\paragraph{Mosiah 20:23} % *MOSIAH-20-23 
		% ---
			\label{MOSIAH-20-23}

			And now Limhi told the king all the things
			concerning his father and the priests that had fled into the wilderness
			and attributed the carrying away of their daughters to them.

		% ---
		\paragraph{Mosiah 20:24} % *MOSIAH-20-24 
		% ---
			\label{MOSIAH-20-24}

			And it came to pass that the king was pacified towards his people,
			and he said unto them:
			Let us go forth to meet my people without arms;
			and I swear unto you with an oath
			that my people shall not slay thy people.

		% ---
		\paragraph{Mosiah 20:25} % *MOSIAH-20-25 
		% ---
			\label{MOSIAH-20-25}

			And it came to pass that they followed the king
			and went forth without arms to meet the Lamanites.
			And it came to pass that they did meet the Lamanites,
			and the king of the Lamanites did bow himself down before them
			and did plead in behalf of the people of Limhi.

		% ---
		\paragraph{Mosiah 20:26} % *MOSIAH-20-26 
		% ---
			\label{MOSIAH-20-26}

			And when the Lamanites saw the people of Limhi,
			that they were without arms,
			that they had compassion on them and were pacified towards them
			and returned with their king in peace to their own land.

	% -----
	\subsubsection{Mosiah 21} % *MOSIAH-21 
	% -----
		\label{MOSIAH-21}
		
		% ---
		\paragraph{Mosiah 21:1} % *MOSIAH-21-1 
		% ---
			\label{MOSIAH-21-1}

			And it came to pass that Limhi and his people returned to the city of Nephi
			and began to dwell in the land again in peace.

		% ---
		\paragraph{Mosiah 21:2} % *MOSIAH-21-2 
		% ---
			\label{MOSIAH-21-2}

			And it came to pass that after many days
			the Lamanites began again to be stirred up in anger against the Nephites,
			and they began to come into the borders of the land round about.

		% ---
		\paragraph{Mosiah 21:3} % *MOSIAH-21-3 
		% ---
			\label{MOSIAH-21-3}

			Now they durst not slay them because of the oath
			which their king had made unto Limhi,
			but they would smite them on their cheeks
			and exercise authority over them
			and began to put heavy burdens upon their backs
			and drive them as they would a dumb ass.

		% ---
		\paragraph{Mosiah 21:4} % *MOSIAH-21-4 
		% ---
			\label{MOSIAH-21-4}

			Yea, all this was done that the word of the Lord might be fulfilled.

		% ---
		\paragraph{Mosiah 21:5} % *MOSIAH-21-5 
		% ---
			\label{MOSIAH-21-5}

			And now the afflictions of the Nephites was great.
			And there was no way that they could deliver themselves out of their hands,
			for the Lamanites had surrounded them on every side.

		% ---
		\paragraph{Mosiah 21:6} % *MOSIAH-21-6 
		% ---
			\label{MOSIAH-21-6}

			And it came to pass that
			the people began to murmur with the king because of their afflictions.
			And they began to be desirous to go against them to battle,
			and they did afflct the king sorely with their complaints.
			Therefore he granted unto them that they should do according to their desires.

		% ---
		\paragraph{Mosiah 21:7} % *MOSIAH-21-7 
		% ---
			\label{MOSIAH-21-7}

			And they gathered themselves together again and put on their armor
			and went forth against the Lamanites to drive them out of their land.

		% ---
		\paragraph{Mosiah 21:8} % *MOSIAH-21-8 
		% ---
			\label{MOSIAH-21-8}

			And it came to pass that the Lamanites did beat them
			and drove them back and slew many of them.

		% ---
		\paragraph{Mosiah 21:9} % *MOSIAH-21-9 
		% ---
			\label{MOSIAH-21-9}

			And now there was a great mourning and lamentation among the people of Limhi,
			the widow a mourning for her husband,
			the son and the daughter a mourning for their father,
			and the brothers for their brethren.

		% ---
		\paragraph{Mosiah 21:10} % *MOSIAH-21-10 
		% ---
			\label{MOSIAH-21-10}

			Now there were a great many widows in the land,
			and they did cry mightily from day to day,
			for a great fear of the Lamanites had come upon them.

		% ---
		\paragraph{Mosiah 21:11} % *MOSIAH-21-11 
		% ---
			\label{MOSIAH-21-11}

			And it came to pass that
			their continual cries did stir up the remainder of the people of Limhi
			to anger against the Lamanites.
			And they went again to battle,
			but they were driven back again,
			suffering much loss.

		% ---
		\paragraph{Mosiah 21:12} % *MOSIAH-21-12 
		% ---
			\label{MOSIAH-21-12}

			Yea, they went again, even the third time,
			and suffered in the like manner.
			And those that were not slain returned again to the city of Nephi.

		% ---
		\paragraph{Mosiah 21:13} % *MOSIAH-21-13 
		% ---
			\label{MOSIAH-21-13}

			And they did humble themselves even to the dust,
			subjecting themselves to the yoke of bondage,
			submitting themselves to be smitten
			and to be driven to and fro and burdened
			according to the desires of their enemies.

		% ---
		\paragraph{Mosiah 21:14} % *MOSIAH-21-14 
		% ---
			\label{MOSIAH-21-14}

			And they did humble themselves even in the depths of humility.
			And they did cry mightily to God,
			yea, even all the day long did they cry unto their God
			that he would deliver them out of their afflictions.

		% ---
		\paragraph{Mosiah 21:15} % *MOSIAH-21-15 
		% ---
			\label{MOSIAH-21-15}

			And now the Lord was slow to hear their cry because of their iniquities.
			Nevertheless the Lord did hear their cries
			and began to soften the hearts of the Lamanites,
			that they began to ease their burdens.
			Yet the Lord did not see fit to deliver them out of bondage.

		% ---
		\paragraph{Mosiah 21:16} % *MOSIAH-21-16 
		% ---
			\label{MOSIAH-21-16}

			And it came to pass that they began to prosper by degrees in the land
			and began to raise grain more abundantly and flocks and herds,
			that they did not suffer with hunger.

		% ---
		\paragraph{Mosiah 21:17} % *MOSIAH-21-17 
		% ---
			\label{MOSIAH-21-17}

			Now there was a great number of women, more than there was of men.
			Therefore king Limhi commanded
			that every man should impart to the support of the widows and their children,
			that they might not perish with hunger.
			And this they did because of the greatness of their number that had been slain.

		% ---
		\paragraph{Mosiah 21:18} % *MOSIAH-21-18 
		% ---
			\label{MOSIAH-21-18}

			Now the people of Limhi kept together in a body
			as much as it was possible
			to secure their grain and their flocks.

		% ---
		\paragraph{Mosiah 21:19} % *MOSIAH-21-19 
		% ---
			\label{MOSIAH-21-19}

			And the king himself did not trust his person without the walls of the city
			unless he took his guards with him,
			fearing that he might by some means fall into the hands of the Lamanites.

		% ---
		\paragraph{Mosiah 21:20} % *MOSIAH-21-20 
		% ---
			\label{MOSIAH-21-20}

			And he caused that his people should watch the land round about,
			that by some means they might take those priests that fled into the wilderness,
			which had stolen the daughters of the Lamanites,
			and that had caused such a great destruction to come upon them.

		% ---
		\paragraph{Mosiah 21:21} % *MOSIAH-21-21 
		% ---
			\label{MOSIAH-21-21}

			For they were desirous to take them that they might punish them,
			for they had come into the land of Nephi by night
			and carried off of their grain and many of their precious things.
			Therefore they laid wait for them.

		% ---
		\paragraph{Mosiah 21:22} % *MOSIAH-21-22 
		% ---
			\label{MOSIAH-21-22}

			And it came to pass that there was no more disturbance
			between the Lamanites and the people of Limhi,
			even until the time that Ammon and his brethren came into the land.

		% ---
		\paragraph{Mosiah 21:23} % *MOSIAH-21-23 
		% ---
			\label{MOSIAH-21-23}

			And the king having been without the gates of the city with his guard,
			he discovered Ammon and his brethren;
			and supposing them to be priests of Noah,
			therefore he caused that they should be taken and bound and cast into prison.
			And had they been the priests of Noah,
			he would have caused that they should be put to death.

		% ---
		\paragraph{Mosiah 21:24} % *MOSIAH-21-24 
		% ---
			\label{MOSIAH-21-24}

			But when he found that they were not,
			but that they were his brethren and had come from the land of Zarahemla,
			he was filled with exceeding great joy.

		% ---
		\paragraph{Mosiah 21:25} % *MOSIAH-21-25 
		% ---
			\label{MOSIAH-21-25}

			Now king Limhi had sent previous to the coming of Ammon
			a small number of men to search for the land of Zarahemla,
			but they could not find it;
			and they were lost in the wilderness.

		% ---
		\paragraph{Mosiah 21:26} % *MOSIAH-21-26 
		% ---
			\label{MOSIAH-21-26}

			Nevertheless they did find a land which had been peopled,
			yea, a land which was covered with dry bones,
			yea, a land which had been peopled and which had been destroyed.
			And they having supposed it to be the land of Zarahemla
			returned to the land of Nephi,
			having arrived in the borders of the land
			not many days before the coming of Ammon.

		% ---
		\paragraph{Mosiah 21:27} % *MOSIAH-21-27 
		% ---
			\label{MOSIAH-21-27}

			And they brought a record with them,
			even a record of the people whose bones they had found;
			and they were engraven on plates of ore.

		% ---
		\paragraph{Mosiah 21:28} % *MOSIAH-21-28 
		% ---
			\label{MOSIAH-21-28}

			And now Limhi was again filled with joy
			on learning from the mouth of Ammon
			that king Benjamin had a gift from God
			whereby he could interpret such engravings
			yea, and Ammon also did rejoice.

		% ---
		\paragraph{Mosiah 21:29} % *MOSIAH-21-29 
		% ---
			\label{MOSIAH-21-29}

			Yet Ammon and his brethren were filled with sorrow
			because so many of his brethren had been slain,

		% ---
		\paragraph{Mosiah 21:30} % *MOSIAH-21-30 
		% ---
			\label{MOSIAH-21-30}

			and also that king Noah and his priests had caused the people
			to commit so many sins and iniquities against God.
			And they also did mourn for the death of Abinadi
			and also for the departure of Alma and the people that went with him,
			who had formed a church of God
			through the strength and power of God
			and faith on the words which had been spoken by Abinadi.

		% ---
		\paragraph{Mosiah 21:31} % *MOSIAH-21-31 
		% ---
			\label{MOSIAH-21-31}

			Yea, they did mourn for their departure,
			for they knew not whither they had fled.
			Now they would have gladly joined with them,
			for they themselves had entered into a covenant with God
			to serve him and keep his commandments.

		% ---
		\paragraph{Mosiah 21:32} % *MOSIAH-21-32 
		% ---
			\label{MOSIAH-21-32}

			And now since the coming of Ammon,
			king Limhi had also entered into a covenant with God,
			and also many of his people,
			to serve him and keep his commandments.

		% ---
		\paragraph{Mosiah 21:33} % *MOSIAH-21-33 
		% ---
			\label{MOSIAH-21-33}

			And it came to pass that
			king Limhi and many of his people was desirous to be baptized,
			but there was none in the land that had authority from God.
			And Ammon declined doing this thing,
			considering himself an unworthy servant.

		% ---
		\paragraph{Mosiah 21:34} % *MOSIAH-21-34 
		% ---
			\label{MOSIAH-21-34}

			Therefore they did not at that time form themselves into a church,
			waiting upon the Spirit of the Lord.
			Now they were desirous to become even as Alma and his brethren,
			which had fled into the wilderness.

		% ---
		\paragraph{Mosiah 21:35} % *MOSIAH-21-35 
		% ---
			\label{MOSIAH-21-35}

			They were desirous to be baptized as a witness and a testimony
			that they were willing to serve God with all their hearts.
			Nevertheless they did prolong the time;
			and an account of their baptism shall be given hereafter.

		% ---
		\paragraph{Mosiah 21:36} % *MOSIAH-21-36 
		% ---
			\label{MOSIAH-21-36}

			And now all the study of Ammon and his people and king Limhi and his people
			was to deliver themselves out of the hands of the Lamanites and from bondage.

	% -----
	\subsubsection{Mosiah 22} % *MOSIAH-22 
	% -----
		\label{MOSIAH-22}
		
		% ---
		\paragraph{Mosiah 22:1} % *MOSIAH-22-1 
		% ---
			\label{MOSIAH-22-1}

			And now it came to pass that
			Ammon and king Limhi began to consult with the people
			how they should deliver themselves out of bondage.
			And even they did cause
			that all the people should gather themselves together;
			and this they did that they might have the voice of the people concerning the matter.

		% ---
		\paragraph{Mosiah 22:2} % *MOSIAH-22-2 
		% ---
			\label{MOSIAH-22-2}

			And it came to pass that they could find no way
			to deliver themselves out of bondage
			except it were to take their women and children
			and their flocks and their herds and their tents
			and depart into the wilderness,
			for the Lamanites being so numerous
			that it was impossible for the people of Limhi to contend with them,
			thinking to deliver themselves out of bondage by the sword.

		% ---
		\paragraph{Mosiah 22:3} % *MOSIAH-22-3 
		% ---
			\label{MOSIAH-22-3}

			Now it came to pass that
			Gideon went forth and stood before the king and said unto him:
			Now, O king, thou hast hitherto hearkened unto my words many times
			when we have been contending with our brethren the Lamanites.

		% ---
		\paragraph{Mosiah 22:4} % *MOSIAH-22-4 
		% ---
			\label{MOSIAH-22-4}

			And now, O king, if thou hast not found me to be an unprofitable servant,
			or if thou hast hitherto listened to my words in any degree
			and they have been of service to thee,
			even so I desire that thou wouldst listen to my words at this time;
			and I will be thy servant and deliver this people out of bondage.

		% ---
		\paragraph{Mosiah 22:5} % *MOSIAH-22-5 
		% ---
			\label{MOSIAH-22-5}

			And the king granted unto him that he might speak,
			and Gideon saith unto him:

		% ---
		\paragraph{Mosiah 22:6} % *MOSIAH-22-6 
		% ---
			\label{MOSIAH-22-6}

			Behold the back pass through the back wall on the back side of the city.
			The Lamanites, or the guards of the Lamanites, by night are drunken.
			Therefore let us send a proclamation among all this people
			that they gather together their flocks and herds,
			that they may drive them into the wilderness by night.

		% ---
		\paragraph{Mosiah 22:7} % *MOSIAH-22-7 
		% ---
			\label{MOSIAH-22-7}

			And I will go according to thy command
			and pay the last tribute of wine to the Lamanites,
			and they will be drunken;
			and we will pass through the secret pass on the left of their camp
			when they are drunken and asleep.

		% ---
		\paragraph{Mosiah 22:8} % *MOSIAH-22-8 
		% ---
			\label{MOSIAH-22-8}

			Thus we will depart
			with our women and our children, our flocks and our herds,
			into the wilderness;
			and we will travel around the land of Shilom.

		% ---
		\paragraph{Mosiah 22:9} % *MOSIAH-22-9 
		% ---
			\label{MOSIAH-22-9}

			And it came to pass that the king hearkened unto the words of Gideon.

		% ---
		\paragraph{Mosiah 22:10} % *MOSIAH-22-10 
		% ---
			\label{MOSIAH-22-10}

			And it came to pass that king Limhi caused
			that his people should gather their flocks together.
			And he sent the tribute of wine to the Lamanites;
			and he also sent more wine as a present unto them,
			and they did drink freely of the wine which king Limhi did send unto them.

		% ---
		\paragraph{Mosiah 22:11} % *MOSIAH-22-11 
		% ---
			\label{MOSIAH-22-11}

			And it came to pass that the people of king Limhi did depart by night
			into the wilderness with their flocks and their herds.
			And they went round about the land of Shilom in the wilderness
			and bent their course towards the land of Zarahemla,
			being led by Ammon and his brethren.

		% ---
		\paragraph{Mosiah 22:12} % *MOSIAH-22-12 
		% ---
			\label{MOSIAH-22-12}

			And they had taken all their gold and silver
			and their precious things which they could carry,
			and also their provisions,
			with them into the wilderness,
			and they pursued their journey.

		% ---
		\paragraph{Mosiah 22:13} % *MOSIAH-22-13 
		% ---
			\label{MOSIAH-22-13}

			And after being many days in the wilderness,
			they arrived in the land of Zarahemla,
			and joined his people and became his subjects.

		% ---
		\paragraph{Mosiah 22:14} % *MOSIAH-22-14 
		% ---
			\label{MOSIAH-22-14}

			And it came to pass that Mosiah received them with joy,
			and he also received their records and also the records
			which had been found by the people of Limhi.

		% ---
		\paragraph{Mosiah 22:15} % *MOSIAH-22-15 
		% ---
			\label{MOSIAH-22-15}

			And now it came to pass when the Lamanites had found
			that the people of Limhi had departed out of the land by night
			that they sent an army into the wilderness to pursue them.

		% ---
		\paragraph{Mosiah 22:16} % *MOSIAH-22-16 
		% ---
			\label{MOSIAH-22-16}

			And after they had pursued them two days,
			they could no longer follow their tracks;
			therefore they were lost in the wilderness.

	% -----
	\subsubsection{Mosiah 23} % *MOSIAH-23 
	% -----
		\label{MOSIAH-23}
		
		An account of Alma and the people of the Lord,
		which was driven into the wilderness by the people of king Noah.
		
		% ---
		\paragraph{Mosiah 23:1} % *MOSIAH-23-1 
		% ---
			\label{MOSIAH-23-1}

			Now Alma having been warned of the Lord
			that the armies of king Noah would come upon them
			and had made it known to his people,
			therefore they gathered together their flocks and took of their grain
			and departed into the wilderness before the armies of king Noah.

		% ---
		\paragraph{Mosiah 23:2} % *MOSIAH-23-2 
		% ---
			\label{MOSIAH-23-2}

			And the Lord did strengthen them,
			that the people of king Noah could not overtake them to destroy them.

		% ---
		\paragraph{Mosiah 23:3} % *MOSIAH-23-3 
		% ---
			\label{MOSIAH-23-3}

			And it came to pass that
			they fled eight days' journey into the wilderness;

		% ---
		\paragraph{Mosiah 23:4} % *MOSIAH-23-4 
		% ---
			\label{MOSIAH-23-4}

			and they came to a land,
			yea, even a very beautiful and pleasant land,
			a land of pure water.

		% ---
		\paragraph{Mosiah 23:5} % *MOSIAH-23-5 
		% ---
			\label{MOSIAH-23-5}

			And it came to pass that they pitched their tents
			and began to till the ground and began to build buildings etc.;
			yea, they were industrious and did labor exceedingly.

		% ---
		\paragraph{Mosiah 23:6} % *MOSIAH-23-6 
		% ---
			\label{MOSIAH-23-6}

			And it came to pass that the people were desirous
			that Alma should be their king,
			for he was beloved by his people.

		% ---
		\paragraph{Mosiah 23:7} % *MOSIAH-23-7 
		% ---
			\label{MOSIAH-23-7}

			But he saith unto them:
			Behold, it is not expedient that we should have a king,
			for thus saith the Lord:
			Ye shall not esteem one flesh above another,
			or one man shall not think himself above another.
			Therefore I say unto you:
			It is not expedient that ye should have a king;

		% ---
		\paragraph{Mosiah 23:8} % *MOSIAH-23-8 
		% ---
			\label{MOSIAH-23-8}

			nevertheless, if it were possible
			that ye could always have just men to be your kings,
			it would be well for you to have a king---

		% ---
		\paragraph{Mosiah 23:9} % *MOSIAH-23-9 
		% ---
			\label{MOSIAH-23-9}

			but remember the iniquity of king Noah and his priests.
			And I myself was caught in a snare
			and did many things which was abominable in the sight of the Lord,
			which caused me sore repentance.

		% ---
		\paragraph{Mosiah 23:10} % *MOSIAH-23-10 
		% ---
			\label{MOSIAH-23-10}

			Nevertheless, after much tribulation the Lord did hear my cries
			and did answer my prayers
			and hath made me an instrument in his hands
			in bringing so many of you to a knowledge of his truth.

		% ---
		\paragraph{Mosiah 23:11} % *MOSIAH-23-11 
		% ---
			\label{MOSIAH-23-11}

			Nevertheless in this I do not glory,
			for I am unworthy to glory of myself.

		% ---
		\paragraph{Mosiah 23:12} % *MOSIAH-23-12 
		% ---
			\label{MOSIAH-23-12}

			And now I say unto you:
			As you have been oppressed by king Noah
			and have been in bondage to him and his priests
			and have been brought into iniquity by them,
			therefore ye were bound with the bands of iniquity.

		% ---
		\paragraph{Mosiah 23:13} % *MOSIAH-23-13 
		% ---
			\label{MOSIAH-23-13}

			And now as ye have been delivered by the power of God out of these bonds,
			yea, even out of the hands of king Noah and his people,
			and also from the bonds of iniquity,
			even so I desire that ye should stand fast in this liberty
			wherewith ye have been made free
			and that ye trust no man to be a king over you---

		% ---
		\paragraph{Mosiah 23:14} % *MOSIAH-23-14 
		% ---
			\label{MOSIAH-23-14}

			and also trusting no one to be your teachers nor your ministers
			except he be a man of God,
			walking in his ways and keeping his commandments.

		% ---
		\paragraph{Mosiah 23:15} % *MOSIAH-23-15 
		% ---
			\label{MOSIAH-23-15}

			Thus did Alma teach his people,
			that every man should love his neighbor as himself,
			that there should be no contention among them.

		% ---
		\paragraph{Mosiah 23:16} % *MOSIAH-23-16 
		% ---
			\label{MOSIAH-23-16}

			And now Alma was their high priest,
			he being the founder of their church.

		% ---
		\paragraph{Mosiah 23:17} % *MOSIAH-23-17 
		% ---
			\label{MOSIAH-23-17}

			And it came to pass that
			none received authority to preach or to teach
			except it were by him from God;
			therefore he consecrated all their priests and all their teachers,
			and none were consecrated except it were just men.

		% ---
		\paragraph{Mosiah 23:18} % *MOSIAH-23-18 
		% ---
			\label{MOSIAH-23-18}

			Therefore they did watch over their people
			and did nourish them with things pertaining to righteousness.

		% ---
		\paragraph{Mosiah 23:19} % *MOSIAH-23-19 
		% ---
			\label{MOSIAH-23-19}

			And it came to pass that
			they began to prosper exceedingly in the land,
			and they called the land Helam.

		% ---
		\paragraph{Mosiah 23:20} % *MOSIAH-23-20 
		% ---
			\label{MOSIAH-23-20}

			And it came to pass that
			they did multiply and prosper exceedingly in the land of Helam.
			And they built a city which they called the city of Helam.

		% ---
		\paragraph{Mosiah 23:21} % *MOSIAH-23-21 
		% ---
			\label{MOSIAH-23-21}

			Nevertheless the Lord seeth fit to chasten his people;
			yea, he trieth their patience and their faith.

		% ---
		\paragraph{Mosiah 23:22} % *MOSIAH-23-22 
		% ---
			\label{MOSIAH-23-22}

			Nevertheless, whosoever putteth his trust in him,
			the same shall be lifted up at the last day;
			yea, and thus it was with this people.

		% ---
		\paragraph{Mosiah 23:23} % *MOSIAH-23-23 
		% ---
			\label{MOSIAH-23-23}

			For behold, I will shew unto you
			that they were brought into bondage,
			and none could deliver them but the Lord their God,
			yea, even the God of Abraham and of Isaac and of Jacob.

		% ---
		\paragraph{Mosiah 23:24} % *MOSIAH-23-24 
		% ---
			\label{MOSIAH-23-24}

			And it came to pass that he did deliver them;
			and he did shew forth his mighty power unto them,
			and great was their rejoicings.

		% ---
		\paragraph{Mosiah 23:25} % *MOSIAH-23-25 
		% ---
			\label{MOSIAH-23-25}

			For behold, it came to pass that
			while they were in the land of Helam,
			yea, in the city of Helam,
			while tilling the land round about,
			behold, an army of the Lamanites were in the borders of the land.

		% ---
		\paragraph{Mosiah 23:26} % *MOSIAH-23-26 
		% ---
			\label{MOSIAH-23-26}

			Now it came to pass that the brethren of Alma fled from their fields
			and gathered themselves together into the city of Helam,
			and they were much frightened because of the appearance of the Lamanites.

		% ---
		\paragraph{Mosiah 23:27} % *MOSIAH-23-27 
		% ---
			\label{MOSIAH-23-27}

			But Alma went forth and stood among them
			and exhorted them that they should not be frightened,
			but that they should remember the Lord their God
			and he would deliver them.

		% ---
		\paragraph{Mosiah 23:28} % *MOSIAH-23-28 
		% ---
			\label{MOSIAH-23-28}

			Therefore they hushed their fears
			and began to cry unto the Lord
			that he would soften the hearts of the Lamanites,
			that they would spare them and their wives and children.

		% ---
		\paragraph{Mosiah 23:29} % *MOSIAH-23-29 
		% ---
			\label{MOSIAH-23-29}

			And it came to pass that the Lord did soften the hearts of the Lamanites.
			And Alma and his brethren went forth
			and delivered themselves up into their hands,
			and the Lamanites took possession of the land of Helam.

		% ---
		\paragraph{Mosiah 23:30} % *MOSIAH-23-30 
		% ---
			\label{MOSIAH-23-30}

			Now the armies of the Lamanites,
			which had followed after the people of king Limhi,
			had been lost in the wilderness for many days.

		% ---
		\paragraph{Mosiah 23:31} % *MOSIAH-23-31 
		% ---
			\label{MOSIAH-23-31}

			And behold, they had found those priests of king Noah
			in a place which they called Amulon,
			and they had began to possess the land of Amulon
			and had began to till the ground.

		% ---
		\paragraph{Mosiah 23:32} % *MOSIAH-23-32 
		% ---
			\label{MOSIAH-23-32}

			Now the name of the leader of those priests was Amulon.

		% ---
		\paragraph{Mosiah 23:33} % *MOSIAH-23-33 
		% ---
			\label{MOSIAH-23-33}

			And it came to pass that Amulon did plead with the Lamanites,
			and he also sent forth their wives,
			which was the daughters of the Lamanites,
			to plead with their brethren
			that they should not destroy their husbands.

		% ---
		\paragraph{Mosiah 23:34} % *MOSIAH-23-34 
		% ---
			\label{MOSIAH-23-34}

			And it came to pass that
			the Lamanites had compassion on Amulon and his brethren
			and did not destroy them because of their wives.

		% ---
		\paragraph{Mosiah 23:35} % *MOSIAH-23-35 
		% ---
			\label{MOSIAH-23-35}

			And Amulon and his brethren did join the Lamanites,
			and they were traveling in the wilderness in search of the land of Nephi
			when they discovered the land of Helam,
			which was possessed by Alma and his brethren.

		% ---
		\paragraph{Mosiah 23:36} % *MOSIAH-23-36 
		% ---
			\label{MOSIAH-23-36}

			And it came to pass that the Lamanites promised unto Alma and his brethren
			that if they would shew them the way which led to the land of Nephi
			that they would grant unto them their lives and their liberty.

		% ---
		\paragraph{Mosiah 23:37} % *MOSIAH-23-37 
		% ---
			\label{MOSIAH-23-37}

			But it came to pass that
			after Alma had shewn them the way that led to the land of Nephi,
			the Lamanites would not keep their promise,
			but they set guards round about the land of Helam over Alma and his brethren.

		% ---
		\paragraph{Mosiah 23:38} % *MOSIAH-23-38 
		% ---
			\label{MOSIAH-23-38}

			And the remainder of them went to the land of Nephi;
			and a part of them returned to the land of Helam
			and also brought with them the wives and the children of the guards
			which had been left in the land.

		% ---
		\paragraph{Mosiah 23:39} % *MOSIAH-23-39 
		% ---
			\label{MOSIAH-23-39}

			And the king of the Lamanites had granted unto Amulon
			that he should be a king and a ruler over his people
			which was in the land of Helam;
			nevertheless he should have no power to do any thing
			contrary to the will of the king of the Lamanites.

	% -----
	\subsubsection{Mosiah 24} % *MOSIAH-24 
	% -----
		\label{MOSIAH-24}
		
		% ---
		\paragraph{Mosiah 24:1} % *MOSIAH-24-1 
		% ---
			\label{MOSIAH-24-1}

			And it came to pass that Amulon did gain favor
			in the eyes of the king of the Lamanites;
			therefore the king of the Lamanites granted unto him and his brethren
			that they should be appointed teachers over his people,
			yea, even over the people
			which was in the land of Shemlon and the land of Shilom,
			and in the land of Amulon.

		% ---
		\paragraph{Mosiah 24:2} % *MOSIAH-24-2 
		% ---
			\label{MOSIAH-24-2}

			For the Lamanites had taken possession of all these lands;
			therefore the king of the Lamanites had appointed kings over all these lands.

		% ---
		\paragraph{Mosiah 24:3} % *MOSIAH-24-3 
		% ---
			\label{MOSIAH-24-3}

			And now the name of the king of the Lamanites was Laman,
			being called after the name of his father;
			and therefore he was called king Laman.
			And he was king over a numerous people.

		% ---
		\paragraph{Mosiah 24:4} % *MOSIAH-24-4 
		% ---
			\label{MOSIAH-24-4}

			And he appointed teachers of the brethren of Amulon
			in every land which was possessed by his people.
			And thus the language of Nephi began to be taught
			among all the people of the Lamanites.

		% ---
		\paragraph{Mosiah 24:5} % *MOSIAH-24-5 
		% ---
			\label{MOSIAH-24-5}

			And they were a people friendly one with another.
			Nevertheless they knew not God;
			neither did the brethren of Amulon teach them any thing
			concerning the Lord their God,
			neither the law of Moses,
			nor did they teach them the words of Abinadi.

		% ---
		\paragraph{Mosiah 24:6} % *MOSIAH-24-6 
		% ---
			\label{MOSIAH-24-6}

			But they taught them that they should keep their record
			and that they might write one to another.

		% ---
		\paragraph{Mosiah 24:7} % *MOSIAH-24-7 
		% ---
			\label{MOSIAH-24-7}

			And thus the Lamanites began to increase in riches
			and began to trade one with another and wax great
			and began to be a cunning and a wise people as to the wisdom of the world,
			yea, a very cunning people,
			delighting in all manner of wickedness and plunder,
			except it were among their own brethren.

		% ---
		\paragraph{Mosiah 24:8} % *MOSIAH-24-8 
		% ---
			\label{MOSIAH-24-8}

			And now it came to pass that
			Amulon began to exercise authority over Alma and his brethren
			and began to persecute him
			and cause that his children should persecute their children.

		% ---
		\paragraph{Mosiah 24:9} % *MOSIAH-24-9 
		% ---
			\label{MOSIAH-24-9}

			For Amulon knew Alma,
			that he had been one of the king's priests
			and that it was he that believed the words of Abinadi
			and was driven out before the king,
			and therefore he was wroth with him;
			for he was subject to king Laman,
			yet he exercised authority over them
			and put tasks upon them and put taskmasters over them.

		% ---
		\paragraph{Mosiah 24:10} % *MOSIAH-24-10 
		% ---
			\label{MOSIAH-24-10}

			And it came to pass that so great was their afflictions
			that they began to cry mightily to God.

		% ---
		\paragraph{Mosiah 24:11} % *MOSIAH-24-11 
		% ---
			\label{MOSIAH-24-11}

			And it came to pass that Amulon commanded them
			that they should stop their cries
			and put guards over them to watch them,
			that whosoever should be found calling upon God should be put to death.

		% ---
		\paragraph{Mosiah 24:12} % *MOSIAH-24-12 
		% ---
			\label{MOSIAH-24-12}

			And it came to pass that
			Alma and his people did not raise their voices to the Lord their God
			but did pour out their hearts to him;
			and he did know the thoughts of their hearts.

		% ---
		\paragraph{Mosiah 24:13} % *MOSIAH-24-13 
		% ---
			\label{MOSIAH-24-13}

			And it came to pass that
			the voice of the Lord came to them in their afflictions, saying:
			Lift up your heads and be of good comfort,
			for I know of the covenant which ye have made unto me.
			And I will covenant with this my people and deliver them out of bondage.

		% ---
		\paragraph{Mosiah 24:14} % *MOSIAH-24-14 
		% ---
			\label{MOSIAH-24-14}

			And I will also ease the burdens which is put upon your shoulders,
			that even you cannot feel them upon your backs,
			even while you are in bondage.
			And this will I do that ye may stand as witnesses for me hereafter,
			and that ye may know of a surety
			that I the Lord God do visit my people in their afflictions.

		% ---
		\paragraph{Mosiah 24:15} % *MOSIAH-24-15 
		% ---
			\label{MOSIAH-24-15}

			And now it came to pass that
			the burdens which was laid upon Alma and his brethren were made light;
			yea, the Lord did strengthen them
			that they could bear up their burdens with ease,
			and they did submit cheerfully and with patience to all the will of the Lord.

		% ---
		\paragraph{Mosiah 24:16} % *MOSIAH-24-16 
		% ---
			\label{MOSIAH-24-16}

			And it came to pass that so great was their faith and their patience
			that the voice of the Lord came unto them again, saying:
			Be of good comfort,
			for on the morrow I will deliver thee out of bondage.

		% ---
		\paragraph{Mosiah 24:17} % *MOSIAH-24-17 
		% ---
			\label{MOSIAH-24-17}

			And he saith unto Alma:
			Thou shalt go before this people,
			and I will go with thee and deliver this people out of bondage.

		% ---
		\paragraph{Mosiah 24:18} % *MOSIAH-24-18 
		% ---
			\label{MOSIAH-24-18}

			Now it came to pass that
			Alma and his people in the nighttime gathered their flocks together
			and also of their grain,
			yea, even all the nighttime were they gathering their flocks together.

		% ---
		\paragraph{Mosiah 24:19} % *MOSIAH-24-19 
		% ---
			\label{MOSIAH-24-19}

			And in the morning the Lord caused a deep sleep to come upon the Lamanites;
			yea, and all their taskmasters were in a profound sleep.

		% ---
		\paragraph{Mosiah 24:20} % *MOSIAH-24-20 
		% ---
			\label{MOSIAH-24-20}

			And it came to pass that Alma and his people departed into the wilderness.
			And when they had traveled all day, they pitched their tents in a valley;
			and they called the name of the valley Alma
			because he led their way in the wilderness.

		% ---
		\paragraph{Mosiah 24:21} % *MOSIAH-24-21 
		% ---
			\label{MOSIAH-24-21}

			Yea, and in the valley of Alma they poured out their thanks to God
			because he had been merciful unto them
			and eased their burdens and had delivered them out of bondage;
			for they were in bondage and none could deliver them
			except it were the Lord their God.

		% ---
		\paragraph{Mosiah 24:22} % *MOSIAH-24-22 
		% ---
			\label{MOSIAH-24-22}

			And they gave thanks to God;
			yea, all their men and all their women and all their children that could speak
			lifted their voices in the praises of their God.

		% ---
		\paragraph{Mosiah 24:23} % *MOSIAH-24-23 
		% ---
			\label{MOSIAH-24-23}

			And now the Lord said unto Alma:
			Haste thee and get thou and this people out of this land,
			for the Lamanites have awoke and doth pursue thee.
			Therefore get thee out of this land;
			and I will stop the Lamanites in this valley,
			that they come no further in pursuit of this people.

		% ---
		\paragraph{Mosiah 24:24} % *MOSIAH-24-24 
		% ---
			\label{MOSIAH-24-24}

			And it came to pass that they departed out of the valley
			and took their journey into the wilderness.

		% ---
		\paragraph{Mosiah 24:25} % *MOSIAH-24-25 
		% ---
			\label{MOSIAH-24-25}

			And it came to pass that after they had been in the wilderness twelve days,
			they arrived to the land of Zarahemla;
			and king Mosiah did also receive them with joy.

	% -----
	\subsubsection{Mosiah 25} % *MOSIAH-25 
	% -----
		\label{MOSIAH-25}
		
		% ---
		\paragraph{Mosiah 25:1} % *MOSIAH-25-1 
		% ---
			\label{MOSIAH-25-1}

			And now king Mosiah caused that all the people should be gathered together.

		% ---
		\paragraph{Mosiah 25:2} % *MOSIAH-25-2 
		% ---
			\label{MOSIAH-25-2}

			Now there were not so many of the children of Nephi,
			or so many of those which were descendants of Nephi,
			as there were of the people of Zarahemla,
			which was a descendant of Muloch
			and those which came with him into the wilderness.

		% ---
		\paragraph{Mosiah 25:3} % *MOSIAH-25-3 
		% ---
			\label{MOSIAH-25-3}

			And there were not so many
			of the people of Nephi and of the people of Zarahemla
			as there was of the Lamanites;
			yea, they were not half so numerous.

		% ---
		\paragraph{Mosiah 25:4} % *MOSIAH-25-4 
		% ---
			\label{MOSIAH-25-4}

			And now all the people of Nephi was assembled together,
			and also all the people of Zarahemla;
			and they were gathered together in two bodies.

		% ---
		\paragraph{Mosiah 25:5} % *MOSIAH-25-5 
		% ---
			\label{MOSIAH-25-5}

			And it came to pass that Mosiah did read
			and caused to be read the records of Zeniff to his people;
			yea, he read the records of the people of Zeniff
			from the time they left the land of Zarahemla until the time they returned again.

		% ---
		\paragraph{Mosiah 25:6} % *MOSIAH-25-6 
		% ---
			\label{MOSIAH-25-6}

			And he also read the account of Alma and his brethren and all their afflictions.
			And he also read the account of Ammon and his brethren and all their afflictions
			from the time they left the land of Zarahemla until the time they returned again.

		% ---
		\paragraph{Mosiah 25:7} % *MOSIAH-25-7 
		% ---
			\label{MOSIAH-25-7}

			And now when Mosiah had made an end of reading the records,
			his people which tarried in the land was struck with wonder and amazement,

		% ---
		\paragraph{Mosiah 25:8} % *MOSIAH-25-8 
		% ---
			\label{MOSIAH-25-8}

			for they knew not what to think.
			For when they beheld those that had been delivered out of bondage,
			they were filled with exceeding great joy.

		% ---
		\paragraph{Mosiah 25:9} % *MOSIAH-25-9 
		% ---
			\label{MOSIAH-25-9}

			And again, when they thought of their brethren
			which had been slain by the Lamanites,
			they were filled with sorrow and even shed many tears of sorrow.

		% ---
		\paragraph{Mosiah 25:10} % *MOSIAH-25-10 
		% ---
			\label{MOSIAH-25-10}

			And again, when they thought of the \hyperref[IMMEDIATE]{immediate}%
				\footnote{%
					% \label{}%
					See also \hyperref[MOSIAH-25-10]{Mosiah 25:10}, which looks like a use of ``immediate'' that's more clearly the ``acting without a medium'' way of using the word. See \hyperref[MOSIAH-2-24]{Mosiah 2:24}.
					}
			goodness of God
			and his power in delivering Alma and his brethren
			out of the hands of the Lamanites and of bondage,
			they did raise their voices and gave thanks to God.

		% ---
		\paragraph{Mosiah 25:11} % *MOSIAH-25-11 
		% ---
			\label{MOSIAH-25-11}

			And again, when they thought upon the Lamanites,
			which was their brethren,
			of their sinful and polluted state,
			they were filled with pain and anguish for the welfare of their souls.

		% ---
		\paragraph{Mosiah 25:12} % *MOSIAH-25-12 
		% ---
			\label{MOSIAH-25-12}

			And it came to pass that
			when those which were the children of Amulon and his brethren,
			which had taken to wife the daughters of the Lamanites
			---they were displeased with the conduct of their fathers
			and they would no longer be called by the names of their fathers---
			therefore they took upon themselves the name of Nephi,
			that they might be called the children of Nephi
			and be numbered among those which were called Nephites.

		% ---
		\paragraph{Mosiah 25:13} % *MOSIAH-25-13 
		% ---
			\label{MOSIAH-25-13}

			And now all the people of Zarahemla were numbered with the Nephites---
			and this because the kingdom had been conferred upon none
			but those which were descendants of Nephi.

		% ---
		\paragraph{Mosiah 25:14} % *MOSIAH-25-14 
		% ---
			\label{MOSIAH-25-14}

			And now it came to pass that
			when Mosiah had made an end of speaking and reading to the people,
			he desired that Alma should also speak to the people.

		% ---
		\paragraph{Mosiah 25:15} % *MOSIAH-25-15 
		% ---
			\label{MOSIAH-25-15}

			And it came to pass that Alma did speak unto them
			when they were assembled together in large bodies;
			and he went from one body to another,
			preaching unto the people repentance and faith on the Lord.

		% ---
		\paragraph{Mosiah 25:16} % *MOSIAH-25-16 
		% ---
			\label{MOSIAH-25-16}

			And he did exhort the people of Limhi and his brethren
			---all those that had been delivered out of bondage---
			that they should remember that it was the Lord that did deliver them.

		% ---
		\paragraph{Mosiah 25:17} % *MOSIAH-25-17 
		% ---
			\label{MOSIAH-25-17}

			And it came to pass that
			after Alma had taught the people many things
			and had made an end of speaking to them
			that king Limhi was desirous that he might be baptized.
			And all his people were desirous that they might be baptized also.

		% ---
		\paragraph{Mosiah 25:18} % *MOSIAH-25-18 
		% ---
			\label{MOSIAH-25-18}

			Therefore Alma did go forth into the water and did baptize them;
			yea, he did baptize them
			after the manner he did his brethren in the waters of Mormon.
			Yea, and as many as he did baptize did belong to the church of God---
			and this because of their belief on the words of Alma.

		% ---
		\paragraph{Mosiah 25:19} % *MOSIAH-25-19 
		% ---
			\label{MOSIAH-25-19}

			And it came to pass that king Mosiah granted unto Alma
			that he might establish churches throughout all the land of Zarahemla,
			and gave him power to ordain priests and teachers over every church.

		% ---
		\paragraph{Mosiah 25:20} % *MOSIAH-25-20 
		% ---
			\label{MOSIAH-25-20}

			Now this was done because there was so many people
			that they could not be all governed by one teacher,
			neither could they all hear the word of God in one assembly.

		% ---
		\paragraph{Mosiah 25:21} % *MOSIAH-25-21 
		% ---
			\label{MOSIAH-25-21}

			Therefore they did assemble themselves together in different bodies,
			being called churches,
			every church having their priests and their teachers,
			and every priest preaching the word
			according as it was delivered to him by the mouth of Alma.

		% ---
		\paragraph{Mosiah 25:22} % *MOSIAH-25-22 
		% ---
			\label{MOSIAH-25-22}

			And thus notwithstanding there being many churches,
			they were all one church,
			yea, even the church of God;
			for there was nothing preached in all the churches
			except it were repentance and faith in God.

		% ---
		\paragraph{Mosiah 25:23} % *MOSIAH-25-23 
		% ---
			\label{MOSIAH-25-23}

			And now there was seven churches in the land of Zarahemla.
			And it came to pass that
			whosoever was desirous to take upon them the name of Christ or of God,
			they did join the churches of God;

		% ---
		\paragraph{Mosiah 25:24} % *MOSIAH-25-24 
		% ---
			\label{MOSIAH-25-24}

			and they were called the people of God.
			And the Lord did pour out his Spirit upon them,
			and they were blessed and prospered in the land.

	% -----
	\subsubsection{Mosiah 26} % *MOSIAH-26 
	% -----
		\label{MOSIAH-26}
		
		% ---
		\paragraph{Mosiah 26:1} % *MOSIAH-26-1 
		% ---
			\label{MOSIAH-26-1}
			
			Now it came to pass that there was many of the rising generation
			that could not understand the words of king Benjamin,
			being little children at the time he spake unto his people;
			and they did not believe the tradition of their fathers.

		% ---
		\paragraph{Mosiah 26:2} % *MOSIAH-26-2 
		% ---
			\label{MOSIAH-26-2}

			They did not believe what had been said concerning the resurrection of the dead,
			neither did they believe concerning the coming of Christ.

		% ---
		\paragraph{Mosiah 26:3} % *MOSIAH-26-3 
		% ---
			\label{MOSIAH-26-3}

			And now because of their unbelief
			they could not understand the word of God;
			and their hearts were hardened.

		% ---
		\paragraph{Mosiah 26:4} % *MOSIAH-26-4 
		% ---
			\label{MOSIAH-26-4}

			And they would not be baptized,
			neither would they join the church.
			And they were a separate people as to their faith and remained so ever after,
			even in their carnal and sinful state,
			for they would not call upon the Lord their God.

		% ---
		\paragraph{Mosiah 26:5} % *MOSIAH-26-5 
		% ---
			\label{MOSIAH-26-5}

			And now in the reign of Mosiah
			they were not half so numerous as the people of God,
			but because of the dissensions among the brethren
			they became more numerous.

		% ---
		\paragraph{Mosiah 26:6} % *MOSIAH-26-6 
		% ---
			\label{MOSIAH-26-6}

			For it came to pass that
			they did deceive many with their flattering words
			which were in the church
			and did cause them to commit many sins.
			Therefore it became expedient
			that those who committed sin that was in the church
			should be admonished by the church.

		% ---
		\paragraph{Mosiah 26:7} % *MOSIAH-26-7 
		% ---
			\label{MOSIAH-26-7}

			And it came to pass that they were brought before the priests
			and delivered up unto the priests by the teachers;
			and the priests brought them before Alma, which was the high priest

		% ---
		\paragraph{Mosiah 26:8} % *MOSIAH-26-8 
		% ---
			\label{MOSIAH-26-8}

			---now king Mosiah had given Alma the authority over the church---

		% ---
		\paragraph{Mosiah 26:9} % *MOSIAH-26-9 
		% ---
			\label{MOSIAH-26-9}

			and it came to pass that Alma did know concerning them,
			for there were many witnesses against them;
			yea, the people stood and testified of their iniquity in abundance.

		% ---
		\paragraph{Mosiah 26:10} % *MOSIAH-26-10 
		% ---
			\label{MOSIAH-26-10}

			Now there had not any such thing happened before in the church.
			Therefore Alma was troubled in his spirit,
			and he caused that they should be brought before the king.

		% ---
		\paragraph{Mosiah 26:11} % *MOSIAH-26-11 
		% ---
			\label{MOSIAH-26-11}

			And he saith unto the king:
			Behold, here are many which we have brought before thee
			which are accused of their brethren.
			Yea, and they have been taken in divers iniquities,
			and they do not repent of their iniquities.
			Therefore we have brought them before thee,
			that thou may judge them according to their crimes.

		% ---
		\paragraph{Mosiah 26:12} % *MOSIAH-26-12 
		% ---
			\label{MOSIAH-26-12}

			But king Mosiah saith unto Alma:
			Behold, I judge them not.
			Therefore I deliver them into thy hands to be judged.

		% ---
		\paragraph{Mosiah 26:13} % *MOSIAH-26-13 
		% ---
			\label{MOSIAH-26-13}

			And now the spirit of Alma was again troubled.
			And he went and inquired of the Lord what he should do concerning this matter,
			for he feared that he should do wrong in the sight of God.

		% ---
		\paragraph{Mosiah 26:14} % *MOSIAH-26-14 
		% ---
			\label{MOSIAH-26-14}

			And it came to pass that after he had poured out his whole soul to God,
			the voice of the Lord came to him, saying:

		% ---
		\paragraph{Mosiah 26:15} % *MOSIAH-26-15 
		% ---
			\label{MOSIAH-26-15}

			Blessed art thou Alma.
			And blessed are they which were baptized in the waters of Mormon.
			Thou art blessed because of thy exceeding faith
			in the words alone of my servant Abinadi.

		% ---
		\paragraph{Mosiah 26:16} % *MOSIAH-26-16 
		% ---
			\label{MOSIAH-26-16}

			And blessed are they because of their exceeding faith
			in the words alone which thou hast spoken unto them.

		% ---
		\paragraph{Mosiah 26:17} % *MOSIAH-26-17 
		% ---
			\label{MOSIAH-26-17}

			And blessed art thou because thou hast established a church among this people.
			And they shall be established, and they shall be my people.

		% ---
		\paragraph{Mosiah 26:18} % *MOSIAH-26-18 
		% ---
			\label{MOSIAH-26-18}

			Yea, blessed is this people, which is willing to bear my name;
			for in my name shall they be called, and they are mine.

		% ---
		\paragraph{Mosiah 26:19} % *MOSIAH-26-19 
		% ---
			\label{MOSIAH-26-19}

			And because thou hast inquired of me concerning the transgressor,
			thou art blessed.

		% ---
		\paragraph{Mosiah 26:20} % *MOSIAH-26-20 
		% ---
			\label{MOSIAH-26-20}

			Thou art my servant,
			and I covenant with thee that thou shalt have eternal life.
			And thou shalt serve me and go forth in my name
			and shall gather together my sheep.

		% ---
		\paragraph{Mosiah 26:21} % *MOSIAH-26-21 
		% ---
			\label{MOSIAH-26-21}

			And he that will hear my voice shall be my sheep;
			and him shall ye receive into the church,
			and him will I also receive.

		% ---
		\paragraph{Mosiah 26:22} % *MOSIAH-26-22 
		% ---
			\label{MOSIAH-26-22}

			For behold, this is my church.
			Whosoever that is baptized shall be baptized unto repentance;
			and whosoever ye receive shall believe in my name,
			and him will I freely forgive.

		% ---
		\paragraph{Mosiah 26:23} % *MOSIAH-26-23 
		% ---
			\label{MOSIAH-26-23}

			For it is I that taketh upon me the sins of the world,
			for it is I that hath created them.
			And it is I that granteth unto him that believeth
			in the end a place at my right hand.

		% ---
		\paragraph{Mosiah 26:24} % *MOSIAH-26-24 
		% ---
			\label{MOSIAH-26-24}

			For behold, in my name are they called;
			and if they know me,
			they shall come forth and shall have a place eternally at my right hand.

		% ---
		\paragraph{Mosiah 26:25} % *MOSIAH-26-25 
		% ---
			\label{MOSIAH-26-25}

			And it shall come to pass that when the second trump shall sound,
			then shall they that never knew me come forth
			and shall stand before me.

		% ---
		\paragraph{Mosiah 26:26} % *MOSIAH-26-26 
		% ---
			\label{MOSIAH-26-26}

			And then shall they know that I am the Lord their God,
			that I am their Redeemer,
			but they would not be redeemed.

		% ---
		\paragraph{Mosiah 26:27} % *MOSIAH-26-27 
		% ---
			\label{MOSIAH-26-27}

			And then will I confess unto them that I never knew them,
			and they shall depart into everlasting fire prepared for the devil and his angels.

		% ---
		\paragraph{Mosiah 26:28} % *MOSIAH-26-28 
		% ---
			\label{MOSIAH-26-28}

			Therefore I say unto you
			that he that will not hear my voice,
			the same shall ye not receive into my church,
			for him I will not receive at the last day.

		% ---
		\paragraph{Mosiah 26:29} % *MOSIAH-26-29 
		% ---
			\label{MOSIAH-26-29}

			Therefore I say unto you: Go;
			and whosoever transgresseth against me,
			him shall ye judge according to the sins which he hath committed.
			And if he confess his sins before thee and me
			and repenteth in the sincerity of his heart,
			him shall ye forgive;
			and I will forgive him also.

		% ---
		\paragraph{Mosiah 26:30} % *MOSIAH-26-30 
		% ---
			\label{MOSIAH-26-30}

			Yea, and as often as my people repent
			will I forgive them their trespasses against me.

		% ---
		\paragraph{Mosiah 26:31} % *MOSIAH-26-31 
		% ---
			\label{MOSIAH-26-31}

			And ye shall also forgive one another your trespasses.
			For verily I say unto you:
			He that forgiveth not his neighbor's trespasses when he saith that he repenteth,
			the same hath brought himself under condemnation.

		% ---
		\paragraph{Mosiah 26:32} % *MOSIAH-26-32 
		% ---
			\label{MOSIAH-26-32}

			Now I say unto you: Go;
			and whosoever will not repent of his sins,
			the same shall not be numbered among my people.
			And this shall be observed from this time forward.

		% ---
		\paragraph{Mosiah 26:33} % *MOSIAH-26-33 
		% ---
			\label{MOSIAH-26-33}

			And it came to pass,
			when Alma had heard these words,
			he wrote them down
			that he might have them
			and that he might judge the people of that church
			according to the commandments of God.

		% ---
		\paragraph{Mosiah 26:34} % *MOSIAH-26-34 
		% ---
			\label{MOSIAH-26-34}

			And it came to pass that Alma went
			and judged those that had been taken in iniquity,
			according to the word of the Lord.

		% ---
		\paragraph{Mosiah 26:35} % *MOSIAH-26-35 
		% ---
			\label{MOSIAH-26-35}

			And whosoever repented of their sins and did confess them,
			them he did number among the people of the church.

		% ---
		\paragraph{Mosiah 26:36} % *MOSIAH-26-36 
		% ---
			\label{MOSIAH-26-36}

			And them that would not confess their sins and repent of their iniquity,
			the same were not numbered among the people of the church;
			and their names were blotted out.

		% ---
		\paragraph{Mosiah 26:37} % *MOSIAH-26-37 
		% ---
			\label{MOSIAH-26-37}

			And it came to pass that Alma did regulate all the affairs of the church.
			And they began again to have peace
			and to prosper exceedingly in the affairs of the church,
			walking circumspectly before God,
			receiving many and baptizing many.

		% ---
		\paragraph{Mosiah 26:38} % *MOSIAH-26-38 
		% ---
			\label{MOSIAH-26-38}

			And now all these things did Alma and his fellow laborers do
			which were over the church,
			walking in all diligence,
			teaching the word of God in all things,
			suffering all manner of afflictions,
			being persecuted by all those who did not belong to the church of God.

		% ---
		\paragraph{Mosiah 26:39} % *MOSIAH-26-39 
		% ---
			\label{MOSIAH-26-39}

			And they did admonish their brethren;
			and they were also admonished every one by the word of God
			according to his sins or to the sins which he had committed,
			being commanded of God to pray without ceasing
			and to give thanks in all things.

	% -----
	\subsubsection{Mosiah 27} % *MOSIAH-27 
	% -----
		\label{MOSIAH-27}
		
		% ---
		\paragraph{Mosiah 27:1} % *MOSIAH-27-1 
		% ---
			\label{MOSIAH-27-1}

			And now it came to pass that
			the persecutions which was inflicted on the church by the unbelievers became so great
			that the church began to murmur
			and complain to their leaders concerning the matter;
			and they did complain to Alma.
			And Alma laid the case before their king Mosiah,
			and Mosiah consulted with his priests.

		% ---
		\paragraph{Mosiah 27:2} % *MOSIAH-27-2 
		% ---
			\label{MOSIAH-27-2}

			And it came to pass that
			king Mosiah sent a proclamation throughout the land round about
			that there should not any unbeliever persecute any of those
			which belonged to the church of God.

		% ---
		\paragraph{Mosiah 27:3} % *MOSIAH-27-3 
		% ---
			\label{MOSIAH-27-3}

			And there was a strict command throughout all the churches
			that there should be no persecutions among them,
			that there should be an equality among all men,

		% ---
		\paragraph{Mosiah 27:4} % *MOSIAH-27-4 
		% ---
			\label{MOSIAH-27-4}

			that they should let no pride nor haughtiness disturb their peace,
			that every man should esteem his neighbor as himself,
			laboring with their own hands for their support.

		% ---
		\paragraph{Mosiah 27:5} % *MOSIAH-27-5 
		% ---
			\label{MOSIAH-27-5}

			Yea, and all their priests and teachers should labor with their own hands
			for their support in all cases save it were in sickness or in much want;
			and doing these things, they did abound in the grace of God.

		% ---
		\paragraph{Mosiah 27:6} % *MOSIAH-27-6 
		% ---
			\label{MOSIAH-27-6}

			And there began to be much peace again in the land.
			And the people began to be very numerous
			and began to scatter abroad upon the face of the earth,
			yea, on the north and on the south, on the east and on the west,
			building large cities and villages in all quarters of the land.

		% ---
		\paragraph{Mosiah 27:7} % *MOSIAH-27-7 
		% ---
			\label{MOSIAH-27-7}

			And the Lord did visit them and prosper them,
			and they became a large and a wealthy people.

		% ---
		\paragraph{Mosiah 27:8} % *MOSIAH-27-8 
		% ---
			\label{MOSIAH-27-8}

			Now the sons of Mosiah was numbered among the unbelievers;
			and also one of the sons of Alma was numbered among them,
			he being called Alma after his father.
			Nevertheless he became a very wicked and an idolatrous man;
			and he was a man of many words and did speak much flattery to the people.
			Therefore he led many of the people to do after the manner of his iniquities.

		% ---
		\paragraph{Mosiah 27:9} % *MOSIAH-27-9 
		% ---
			\label{MOSIAH-27-9}

			And he became a great hinderment to the prosperity of the church of God,
			stealing away the hearts of the people,
			causing much dissension among the people,
			giving a chance for the enemy of God to exercise his power over them.

		% ---
		\paragraph{Mosiah 27:10} % *MOSIAH-27-10 
		% ---
			\label{MOSIAH-27-10}

			And now it came to pass that
			while he was going about to destroy%
				\footnote{%
					% \label{}%
					Compare Paul in \hyperref[GALATIANS-1-23]{Galatians 1:23}.
					}
			the church of God
			---for he did go about secretly with the sons of Mosiah,
			seeking to destroy the church
			and to lead astray the people of the Lord
			contrary to the commandments of God or even the king---

		% ---
		\paragraph{Mosiah 27:11} % *MOSIAH-27-11 
		% ---
			\label{MOSIAH-27-11}

			and as I said unto you,
			as they were going about rebelling against God,
			behold, the angel of the Lord appeared unto them,
			and he descended as it were in a cloud.
			And he spake as it were with a voice of thunder,
			which caused the earth to shake upon which they stood.

		% ---
		\paragraph{Mosiah 27:12} % *MOSIAH-27-12 
		% ---
			\label{MOSIAH-27-12}

			And so great was their astonishment that they fell to the earth
			and understood not the words which he spake unto them.

		% ---
		\paragraph{Mosiah 27:13} % *MOSIAH-27-13 
		% ---
			\label{MOSIAH-27-13}

			Nevertheless he cried again, saying:
			Alma, arise and stand forth.
			For why persecutest thou the church of God?
			For the Lord hath said:
			This is my church and I will establish it;
			and nothing shall overthrow it save it is the transgression of my people.

		% ---
		\paragraph{Mosiah 27:14} % *MOSIAH-27-14 
		% ---
			\label{MOSIAH-27-14}

			And again the angel saith:
			Behold, the Lord hath heard the prayers of his people
			and also the prayers of his servant Alma, which is thy father.
			For he hath prayed with much faith concerning thee,
			that thou mightest be brought to the knowledge of the truth.
			Therefore for this purpose have I come
			to convince thee of the power and authority of God,
			that the prayers of his servants might be answered according to their faith.

		% ---
		\paragraph{Mosiah 27:15} % *MOSIAH-27-15 
		% ---
			\label{MOSIAH-27-15}

			And now behold, can ye dispute the power of God?
			For behold, doth not my voice shake the earth?
			And can ye not also behold me before you?
			And I am sent from God.

		% ---
		\paragraph{Mosiah 27:16} % *MOSIAH-27-16 
		% ---
			\label{MOSIAH-27-16}

			Now I say unto thee: Go;
			and remember the captivity of thy fathers
			in the land of Helam and in the land of Nephi,
			and remember how great things he hath done for them.
			For they were in bondage and he hath delivered them.
			And now I say unto thee Alma:
			Go thy way and seek to destroy the church no more,
			that their prayers may be answered.
			And this even if thou wilt of thyself be cast off!

		% ---
		\paragraph{Mosiah 27:17} % *MOSIAH-27-17 
		% ---
			\label{MOSIAH-27-17}

			And now it came to pass that
			these were the last words which the angel spake unto Alma,
			and he departed.

		% ---
		\paragraph{Mosiah 27:18} % *MOSIAH-27-18 
		% ---
			\label{MOSIAH-27-18}

			And now Alma and those that were with him fell again to the earth,
			for great was their astonishment.
			For with their own eyes they had beheld an angel of the Lord,
			and his voice was as thunder which shook the earth.
			And they knew that there was nothing save the power of God
			that could shake the earth and cause it to tremble
			as though it would part asunder.

		% ---
		\paragraph{Mosiah 27:19} % *MOSIAH-27-19 
		% ---
			\label{MOSIAH-27-19}

			And now the astonishment of Alma was so great
			that he became dumb,
			that he could not open his mouth;
			yea, and he became weak,
			even that he could not move his hands.
			Therefore he was taken by those that were with him and carried helpless,
			even until he was laid before his father.

		% ---
		\paragraph{Mosiah 27:20} % *MOSIAH-27-20 
		% ---
			\label{MOSIAH-27-20}

			And they rehearsed unto his father all that had happened unto them.
			And his father rejoiced,
			for he knew that it was the power of God.

		% ---
		\paragraph{Mosiah 27:21} % *MOSIAH-27-21 
		% ---
			\label{MOSIAH-27-21}

			And he caused that a multitude should be gathered together,
			that they might witness what the Lord had done for his son,
			and also for those that were with him.

		% ---
		\paragraph{Mosiah 27:22} % *MOSIAH-27-22 
		% ---
			\label{MOSIAH-27-22}

			And he caused that the priests should assemble themselves together.
			And they began to fast and to pray to the Lord their God
			that he would open the mouth of Alma, that he might speak,
			and also that his limbs might receive their strength,
			that the eyes of the people might be opened,
			to see and know of the goodness and glory of God.

		% ---
		\paragraph{Mosiah 27:23} % *MOSIAH-27-23 
		% ---
			\label{MOSIAH-27-23}

			And it came to pass
			after they had fasted and prayed for the space of two days and two nights,
			the limbs of Alma received their strength.
			And he stood up and began to speak unto them,
			bidding them to be of good comfort;

		% ---
		\paragraph{Mosiah 27:24} % *MOSIAH-27-24 
		% ---
			\label{MOSIAH-27-24}

			for, said he, I have repented of my sins
			and have been redeemed of the Lord.
			Behold, I am born of the Spirit.

		% ---
		\paragraph{Mosiah 27:25} % *MOSIAH-27-25 
		% ---
			\label{MOSIAH-27-25}

			And the Lord said unto me:
			Marvel not that all mankind,
			yea, men and women
			---all nations, kindreds, tongues, and people---
			must be born again, yea, born of God,
			changed from their carnal and fallen state to a state of righteousness,
			being redeemed of God,
			becoming his sons and daughters.

		% ---
		\paragraph{Mosiah 27:26} % *MOSIAH-27-26 
		% ---
			\label{MOSIAH-27-26}

			And thus they become new creatures;
			and unless they do this, they can in no wise inherit the kingdom of God.

		% ---
		\paragraph{Mosiah 27:27} % *MOSIAH-27-27 
		% ---
			\label{MOSIAH-27-27}

			I say unto you:
			Unless this be the case, they must be cast off.
			And this I know because I was like to be cast off.

		% ---
		\paragraph{Mosiah 27:28} % *MOSIAH-27-28 
		% ---
			\label{MOSIAH-27-28}

			Nevertheless, after wading through much tribulation,
			repenting
				\footnote{%
					% \label{}%
					For ``nigh unto death,'' see \hyperref[PHILIPPIANS-2-27]{Philippians 2:27} and \hyperref[PHILIPPIANS-2-30]{2:30}. 2:27 also mentions mercy, the same as this verse here. Note that the Greek for ``nigh unto death'' differs between the two Philippians verses.
					}%
			nigh unto death,
			the Lord in mercy hath seen fit to snatch me out of an everlasting burning;
			and I am born of God.

		% ---
		\paragraph{Mosiah 27:29} % *MOSIAH-27-29 
		% ---
			\label{MOSIAH-27-29}

			My soul hath been redeemed from the gall of bitterness and bonds of iniquity.
			I was in the darkest abyss,
			but now I behold the marvelous light of God.
			My soul was racked with eternal torment,
			but I am snatched and my soul is pained no more.

		% ---
		\paragraph{Mosiah 27:30} % *MOSIAH-27-30 
		% ---
			\label{MOSIAH-27-30}

			I rejected my Redeemer
			and denied that which had been spoken of by our fathers.
			But now I know that they may foresee that he will come
			and that he remembereth every creature of his creating
			and he will make himself manifest unto all.

		% ---
		\paragraph{Mosiah 27:31} % *MOSIAH-27-31 
		% ---
			\label{MOSIAH-27-31}

			Yea, every knee shall bow and every tongue confess before him.
			Yea, even at the last day when all men shall stand to be judged of him,
			then shall they confess that he is God;
			then shall they confess who live without God in the world
			that the judgment of an everlasting punishment is just upon them.
			And they shall quake and tremble and shrink
			beneath the glance of his all-searching eye.

		% ---
		\paragraph{Mosiah 27:32} % *MOSIAH-27-32 
		% ---
			\label{MOSIAH-27-32}

			And now it came to pass that
			Alma began from this time forward to teach the people,
			and those which were with Alma at the time the angel appeared unto them,
			traveling round about through all the land,
			publishing to all the people the things which they had heard and seen,
			and preaching the word of God,
			in much tribulation,
			being greatly persecuted by those which were unbelievers,
			being smitten by many of them.

		% ---
		\paragraph{Mosiah 27:33} % *MOSIAH-27-33 
		% ---
			\label{MOSIAH-27-33}

			But notwithstanding all this, they did impart much consolation to the church,
			confirming their faith and exhorting them with long-suffering and much travail
			to keep the commandments of God.

		% ---
		\paragraph{Mosiah 27:34} % *MOSIAH-27-34 
		% ---
			\label{MOSIAH-27-34}

			And four of them were the sons of Mosiah.
			And their names were Ammon and Aaron and Omner and Himni;
			these were the names of the sons of Mosiah.

		% ---
		\paragraph{Mosiah 27:35} % *MOSIAH-27-35 
		% ---
			\label{MOSIAH-27-35}

			And after they had traveled throughout all the land of Zarahemla
			and among all the people which was under the reign of king Mosiah,
			zealously striving to repair all the injuries which they had done to the church,
			confessing all their sins and publishing all the things which they had seen,
			and explaining the prophecies and the scriptures to all who desired to hear them---

		% ---
		\paragraph{Mosiah 27:36} % *MOSIAH-27-36 
		% ---
			\label{MOSIAH-27-36}

			and thus they were instruments in the hands of God
			in bringing many to the knowledge of the truth,
			yea, to the knowledge of their Redeemer.

		% ---
		\paragraph{Mosiah 27:37} % *MOSIAH-27-37 
		% ---
			\label{MOSIAH-27-37}

			And how blessed are they!
			For they did publish peace;
			they did publish good tidings of good,
			and they did declare unto the people that the Lord reigneth.

	% -----
	\subsubsection{Mosiah 28} % *MOSIAH-28 
	% -----
		\label{MOSIAH-28}
		
		% ---
		\paragraph{Mosiah 28:1} % *MOSIAH-28-1 
		% ---
			\label{MOSIAH-28-1}

			Now it came to pass that after the sons of Mosiah had done all these things,
			they took a small number with them and returned to their father the king
			and desired of him that he would grant unto them
			that they might with those whom they had selected go up to the land of Nephi,
			that they might preach the things which they had heard,
			and that they might impart the word of God to their brethren the Lamanites,

		% ---
		\paragraph{Mosiah 28:2} % *MOSIAH-28-2 
		% ---
			\label{MOSIAH-28-2}

			that perhaps they might bring them to the knowledge of the Lord their God
			and convince them of the iniquity of their fathers,
			and that perhaps they might cure them of their hatred towards the Nephites,
			that they might also be brought to rejoice in the Lord their God,
			that they might become friendly to one another,
			and that there should be no more contentions
			in all the land which the Lord their God hath given them.

		% ---
		\paragraph{Mosiah 28:3} % *MOSIAH-28-3 
		% ---
			\label{MOSIAH-28-3}

			Now they were desirous that salvation should be declared to every creature,
			for they could not bear that any human soul should perish;
			yea, even the very thoughts that any soul should endure endless torment
			did cause them to quake and tremble.

		% ---
		\paragraph{Mosiah 28:4} % *MOSIAH-28-4 
		% ---
			\label{MOSIAH-28-4}

			And thus did the Spirit of the Lord work upon them.
			For they were the very vilest of sinners,
			and the Lord saw fit in his infinite mercy to spare them.
			Nevertheless they suffered much anguish of soul because of their iniquities---
			and suffering much fearing that they should be cast off forever.

		% ---
		\paragraph{Mosiah 28:5} % *MOSIAH-28-5 
		% ---
			\label{MOSIAH-28-5}

			And it came to pass that they did plead with their father many days
			that they might go up to the land of Nephi.

		% ---
		\paragraph{Mosiah 28:6} % *MOSIAH-28-6 
		% ---
			\label{MOSIAH-28-6}

			And it came to pass that king Mosiah went and inquired of the Lord
			if he should let his sons go up among the Lamanites to preach the word.

		% ---
		\paragraph{Mosiah 28:7} % *MOSIAH-28-7 
		% ---
			\label{MOSIAH-28-7}

			And the Lord said unto Mosiah:
			Let them go up, for many shall believe on their words;
			and they shall have eternal life.
			And I will deliver thy sons out of the hands of the Lamanites.

		% ---
		\paragraph{Mosiah 28:8} % *MOSIAH-28-8 
		% ---
			\label{MOSIAH-28-8}

			And it came to pass that Mosiah granted
			that they might go and do according to their request.

		% ---
		\paragraph{Mosiah 28:9} % *MOSIAH-28-9 
		% ---
			\label{MOSIAH-28-9}

			And they took their journey into the wilderness
			to go up to preach the word among the Lamanites.
			And I shall give an account of their proceedings hereafter.

		% ---
		\paragraph{Mosiah 28:10} % *MOSIAH-28-10 
		% ---
			\label{MOSIAH-28-10}

			Now king Mosiah had no one to confer the kingdom upon,
			for there was not any of his sons which would accept of the kingdom.

		% ---
		\paragraph{Mosiah 28:11} % *MOSIAH-28-11 
		% ---
			\label{MOSIAH-28-11}

			Therefore he took the records which were engraven upon the plates of brass,
			and also the plates of Nephi,
			and all the things which he had kept and preserved
			according to the commandments of God,
			and after having translated and caused to be written the records
			which were on the plates of gold which had been found by the people of Limhi,
			which was delivered to him by the hand of Limhi---

		% ---
		\paragraph{Mosiah 28:12} % *MOSIAH-28-12 
		% ---
			\label{MOSIAH-28-12}

			and this he done because of the great anxiety of his people,
			for they were desirous beyond measure
			to know concerning those people which had been destroyed.

		% ---
		\paragraph{Mosiah 28:13} % *MOSIAH-28-13 
		% ---
			\label{MOSIAH-28-13}

			And now he translated them by the means of those two stones
			which was fastened into the two rims of a bow.

		% ---
		\paragraph{Mosiah 28:14} % *MOSIAH-28-14 
		% ---
			\label{MOSIAH-28-14}

			Now these things was prepared from the beginning
			and was handed down from generation to generation
			for the purpose of interpreting languages.

		% ---
		\paragraph{Mosiah 28:15} % *MOSIAH-28-15 
		% ---
			\label{MOSIAH-28-15}

			And they have been kept and preserved by the hand of the Lord,
			that he should discover to every creature which should possess the land
			the iniquities and abominations of his people.

		% ---
		\paragraph{Mosiah 28:16} % *MOSIAH-28-16 
		% ---
			\label{MOSIAH-28-16}

			And whosoever has the things is called seer, after the manner of old times.

		% ---
		\paragraph{Mosiah 28:17} % *MOSIAH-28-17 
		% ---
			\label{MOSIAH-28-17}

			Now after Mosiah had finished translating these records,
			behold, it gave an account of the people which was destroyed
			from the time that they were destroyed
			back to the building of the great tower,
			at the time the Lord confounded the language of the people,
			and they were scattered abroad upon the face of all the earth,
			yea, and even from that time until the creation of Adam.

		% ---
		\paragraph{Mosiah 28:18} % *MOSIAH-28-18 
		% ---
			\label{MOSIAH-28-18}

			Now this account did cause the people of Mosiah to mourn exceedingly;
			yea, they were filled with sorrow.
			Nevertheless it gave them much knowledge,
			in the which they did rejoice.

		% ---
		\paragraph{Mosiah 28:19} % *MOSIAH-28-19 
		% ---
			\label{MOSIAH-28-19}

			And this account shall be written hereafter;
			for behold, it is expedient
			that all people should know the things which are written in this account.

		% ---
		\paragraph{Mosiah 28:20} % *MOSIAH-28-20 
		% ---
			\label{MOSIAH-28-20}

			And now as I said unto you that
			after king Mosiah had done these things,
			he took the plates of brass and all the things which he had kept
			and conferred them upon Alma, which was the son of Alma
			---yea, all the records and also the interpreters---
			and conferred them upon him,
			and commanding him that he should keep and preserve them
			and also keep a record of the people,
			handing them down from one generation to another,
			even as they had been handed down from the time that Lehi left Jerusalem.

	% -----
	\subsubsection{Mosiah 29} % *MOSIAH-29 
	% -----
		\label{MOSIAH-29}
		
		% ---
		\paragraph{Mosiah 29:1} % *MOSIAH-29-1 
		% ---
			\label{MOSIAH-29-1}

			Now when Mosiah had done this,
			he sent out through all the land among all the people,
			desiring to know their will concerning who should be their king.

		% ---
		\paragraph{Mosiah 29:2} % *MOSIAH-29-2 
		% ---
			\label{MOSIAH-29-2}

			And it came to pass that the voice of the people came, saying:
			We are desirous that Aaron thy son should be our king and our ruler.

		% ---
		\paragraph{Mosiah 29:3} % *MOSIAH-29-3 
		% ---
			\label{MOSIAH-29-3}

			Now Aaron had gone up to the land of Nephi;
			therefore the king could not confer the kingdom upon him.
			Neither would Aaron take upon him the kingdom,
			neither was any of the sons of Mosiah willing to take upon them the kingdom.

		% ---
		\paragraph{Mosiah 29:4} % *MOSIAH-29-4 
		% ---
			\label{MOSIAH-29-4}

			Therefore king Mosiah sent again among the people;
			yea, even a written word sent he among the people.
			And these were the words that were written, saying:

		% ---
		\paragraph{Mosiah 29:5} % *MOSIAH-29-5 
		% ---
			\label{MOSIAH-29-5}

			Behold, O ye my people or my brethren
			---for I esteem you as such---
			for I desire that ye should consider the cause which ye are called to consider,
			for ye are desirous to have a king.

		% ---
		\paragraph{Mosiah 29:6} % *MOSIAH-29-6 
		% ---
			\label{MOSIAH-29-6}

			Now I declare unto you that he to whom that the kingdom doth rightly belong
			hath declined and will not take upon him the kingdom.

		% ---
		\paragraph{Mosiah 29:7} % *MOSIAH-29-7 
		% ---
			\label{MOSIAH-29-7}

			And now if there should be another appointed in his stead,
			behold, I fear there would rise contentions among you.
			And who knoweth but what my son to whom the kingdom doth belong
			should turn to be angry and draw away a part of this people after him
			---which will cause wars and contentions among you,
			which would be the cause of shedding much blood
			and perverting the way of the Lord ---
			yea, and destroy the souls of much people.

		% ---
		\paragraph{Mosiah 29:8} % *MOSIAH-29-8 
		% ---
			\label{MOSIAH-29-8}

			Now I say unto you:
			Let us be wise and consider these things.
			For we have no right to destroy my son,
			neither should we have any right to destroy another
			if he should be appointed in his stead.

		% ---
		\paragraph{Mosiah 29:9} % *MOSIAH-29-9 
		% ---
			\label{MOSIAH-29-9}

			And if my son should turn again to his pride and vain things,
			he would recall the things which he had said
			and claim his right to the kingdom,
			which would cause him and also this people to commit much sin.

		% ---
		\paragraph{Mosiah 29:10} % *MOSIAH-29-10 
		% ---
			\label{MOSIAH-29-10}

			And now let us be wise and look forward to these things
			and do that which will make for the peace of this people.

		% ---
		\paragraph{Mosiah 29:11} % *MOSIAH-29-11 
		% ---
			\label{MOSIAH-29-11}

			Therefore I will be your king the remainder of my days.
			Nevertheless let us appoint judges to judge this people according to our law,
			and we will newly arrange the affairs of this people.
			For we will appoint wise men to be judges
			that will judge this people according to the commandments of God.

		% ---
		\paragraph{Mosiah 29:12} % *MOSIAH-29-12 
		% ---
			\label{MOSIAH-29-12}

			Now it is better that a man should be judged of God than of man;
			for the judgments of God are always just,
			but the judgments of man are not always just.

		% ---
		\paragraph{Mosiah 29:13} % *MOSIAH-29-13 
		% ---
			\label{MOSIAH-29-13}

			Therefore if it were possible that ye could have just men to be your kings
			which would establish the laws of God
			and judge this people according to his commandments,
			yea, if ye could have men for your kings
			which would do even as my father Benjamin did for this people,
			I say unto you,
			if this could always be the case,
			then it would be expedient
			that ye should always have kings to rule over you.

		% ---
		\paragraph{Mosiah 29:14} % *MOSIAH-29-14 
		% ---
			\label{MOSIAH-29-14}

			And even I myself have labored
			with all the power and faculties which I have possessed
			to teach you the commandments of God
			and to establish peace throughout the land,
			that there should be no wars nor contentions,
			no stealing nor plundering nor murdering nor no manner of iniquity.

		% ---
		\paragraph{Mosiah 29:15} % *MOSIAH-29-15 
		% ---
			\label{MOSIAH-29-15}

			And whosoever hath committed iniquity,
			him have I punished according to the crime which he hath committed
			according to the law which hath been given to us by our fathers.

		% ---
		\paragraph{Mosiah 29:16} % *MOSIAH-29-16 
		% ---
			\label{MOSIAH-29-16}

			Now I say unto you that because all men are not just,
			it is not expedient that ye should have a king or kings to rule over you.

		% ---
		\paragraph{Mosiah 29:17} % *MOSIAH-29-17 
		% ---
			\label{MOSIAH-29-17}

			For behold, how much iniquity doth one wicked king cause to be committed!
			Yea, and what great destruction!

		% ---
		\paragraph{Mosiah 29:18} % *MOSIAH-29-18 
		% ---
			\label{MOSIAH-29-18}

			Yea, remember king Noah, his wickedness and his abominations,
			and also the wickedness and abominations of his people.
			Behold, what great destruction did come upon them!
			And also because of their iniquities they were brought into bondage.

		% ---
		\paragraph{Mosiah 29:19} % *MOSIAH-29-19 
		% ---
			\label{MOSIAH-29-19}

			And were it not for the interposition of their all-wise Creator
			---and this because of their sincere repentance---
			they must have unavoidably remained in bondage until now.

		% ---
		\paragraph{Mosiah 29:20} % *MOSIAH-29-20 
		% ---
			\label{MOSIAH-29-20}

			But behold, he did deliver them
			because they did humble themselves before him;
			and because they cried mightily unto him,
			he did deliver them out of bondage.
			And thus doth the Lord work with his power
			in all cases among the children of men,
			extending the arm of mercy towards them that put their trust in him.

		% ---
		\paragraph{Mosiah 29:21} % *MOSIAH-29-21 
		% ---
			\label{MOSIAH-29-21}

			And behold, now I say unto you:
			Ye cannot dethrone an iniquitous king
			save it be through much contention and the shedding of much blood.

		% ---
		\paragraph{Mosiah 29:22} % *MOSIAH-29-22 
		% ---
			\label{MOSIAH-29-22}

			For behold, he hath his friends in iniquity,
			and he keepeth his guards about him,
			and he teareth up the laws of those
			which have reigned in righteousness before him,
			and he trampleth under his feet the commandments of God.

		% ---
		\paragraph{Mosiah 29:23} % *MOSIAH-29-23 
		% ---
			\label{MOSIAH-29-23}

			And he enacteth laws and sendeth them forth among his people,
			yea, laws after the manner of his own wickedness;
			and whosoever doth not obey his laws,
			he causeth to be destroyed.
			And whosoever doth rebel against him,
			he will send his armies against them to war;
			and if he can, he will destroy them.
			And thus an unrighteous king doth pervert the ways of all righteousness.

		% ---
		\paragraph{Mosiah 29:24} % *MOSIAH-29-24 
		% ---
			\label{MOSIAH-29-24}

			And now behold, I say unto you:
			It is not expedient that such abominations should come upon you.

		% ---
		\paragraph{Mosiah 29:25} % *MOSIAH-29-25 
		% ---
			\label{MOSIAH-29-25}

			Therefore choose you by the voice of this people judges,
			that ye may be judged according to the laws
			which hath been given you by our fathers,
			which are correct
			and which was given them by the hand of the Lord.

		% ---
		\paragraph{Mosiah 29:26} % *MOSIAH-29-26 
		% ---
			\label{MOSIAH-29-26}

			Now it is not common that the voice of the people desireth any thing
			contrary to that which is right,
			but it is common for the lesser part of the people
			to desire that which is not right.
			Therefore this shall ye observe and make it your law,
			to do your business by the voice of the people.

		% ---
		\paragraph{Mosiah 29:27} % *MOSIAH-29-27 
		% ---
			\label{MOSIAH-29-27}

			And if the time cometh that the voice of the people doth choose iniquity,
			then is the time that the judgments of God will come upon you.
			Yea, then is the time he will visit you with great destruction,
			even as he hath hitherto visited this land.

		% ---
		\paragraph{Mosiah 29:28} % *MOSIAH-29-28 
		% ---
			\label{MOSIAH-29-28}

			And now if ye have judges
			and they do not judge you according to the law which has been given,
			ye can cause that he may be judged of a higher judge.

		% ---
		\paragraph{Mosiah 29:29} % *MOSIAH-29-29 
		% ---
			\label{MOSIAH-29-29}

			If your higher judges doth not judge righteous judgments,
			ye shall cause that a small number of your lower judges
			should be gathered together
			and they shall judge your higher judges
			according to the voice of the people.

		% ---
		\paragraph{Mosiah 29:30} % *MOSIAH-29-30 
		% ---
			\label{MOSIAH-29-30}

			And I commanded you to do these things in the fear of the Lord;
			and I commanded you to do these things and that ye have no king,
			that if these people commit sins and iniquities,
			they shall be answered upon their own heads.

		% ---
		\paragraph{Mosiah 29:31} % *MOSIAH-29-31 
		% ---
			\label{MOSIAH-29-31}

			For behold, I say unto you:
			The sins of many people have been caused by the iniquities of their kings;
			therefore their iniquities are answered upon the heads of their kings.

		% ---
		\paragraph{Mosiah 29:32} % *MOSIAH-29-32 
		% ---
			\label{MOSIAH-29-32}

			And now I desire that this unequality should be no more in this land,
			especially among this my people.
			But I desire that this land be a land of liberty
			and every man may enjoy his rights and privileges alike
			so long as the Lord sees fit that we may live and inherit the land,
			yea, even as long as any of our posterity remaineth upon the face of the land.

		% ---
		\paragraph{Mosiah 29:33} % *MOSIAH-29-33 
		% ---
			\label{MOSIAH-29-33}

			And many more things did king Mosiah write unto them,
			unfolding unto them all the trials and troubles of a righteous king,
			yea, all the travails of soul for their people,
			and also all the murmurings of the people to their king;
			and he explained it all unto them.

		% ---
		\paragraph{Mosiah 29:34} % *MOSIAH-29-34 
		% ---
			\label{MOSIAH-29-34}

			And he told them that these things had not ought to be,
			but that the burden should come upon all the people,
			that every man might bear his part.

		% ---
		\paragraph{Mosiah 29:35} % *MOSIAH-29-35 
		% ---
			\label{MOSIAH-29-35}

			And he also unfolded unto them all the disadvantages they labored under
			by having an unrighteous king to rule over them,

		% ---
		\paragraph{Mosiah 29:36} % *MOSIAH-29-36 
		% ---
			\label{MOSIAH-29-36}

			yea, all his iniquities and abominations
			and all the wars and contentions and bloodshed
			and the stealing and the plundering
			and the committing of whoredoms
			and all manner of iniquities which cannot be enumerated,
			telling them that these things ought not to be,
			that they was expressly repugnant to the commandments of God.

		% ---
		\paragraph{Mosiah 29:37} % *MOSIAH-29-37 
		% ---
			\label{MOSIAH-29-37}

			And now it came to pass
			after king Mosiah had sent these things forth among the people,
			they were convinced of the truth of his words.

		% ---
		\paragraph{Mosiah 29:38} % *MOSIAH-29-38 
		% ---
			\label{MOSIAH-29-38}

			Therefore they relinquished their desires for a king
			and became exceedingly anxious
			that every man should have an equal chance throughout all the land;
			yea, and every man expressed a willingness to answer for his own sins.

		% ---
		\paragraph{Mosiah 29:39} % *MOSIAH-29-39 
		% ---
			\label{MOSIAH-29-39}

			Therefore it came to pass that
			they assembled themselves together in bodies throughout the land
			to cast in their voices concerning who should be their judges
			to judge them according to the law which had been given them.
			And they were exceedingly rejoiced
			because of the liberty which had been granted unto them.

		% ---
		\paragraph{Mosiah 29:40} % *MOSIAH-29-40 
		% ---
			\label{MOSIAH-29-40}

			And they did wax strong in love towards Mosiah;
			yea, they did esteem him more than any other man.
			For they did not look upon him as a tyrant who was seeking for gain,
			yea, for that lucre which doth corrupt the soul;
			for he had not exacted riches of them,
			neither had he delighted in the shedding of blood,
			but he had established peace in the land.
			And he had granted unto his people
			that they should be delivered from all manner of bondage.
			Therefore they did esteem him,
			yea, exceedingly beyond measure.

		% ---
		\paragraph{Mosiah 29:41} % *MOSIAH-29-41 
		% ---
			\label{MOSIAH-29-41}

			And it came to pass that they did appoint judges to rule over them,
			or to judge them according to the law;
			and this they done throughout all the land.

		% ---
		\paragraph{Mosiah 29:42} % *MOSIAH-29-42 
		% ---
			\label{MOSIAH-29-42}

			And it came to pass that Alma was appointed to be the chief judge,
			he being also the high priest,
			his father having conferred the office upon him
			and had given him the charge concerning all the affairs of the church.

		% ---
		\paragraph{Mosiah 29:43} % *MOSIAH-29-43 
		% ---
			\label{MOSIAH-29-43}

			And now it came to pass that Alma did walk in the ways of the Lord,
			and he did keep his commandments,
			and he did judge righteous judgments.
			And there was continual peace through the land.

		% ---
		\paragraph{Mosiah 29:44} % *MOSIAH-29-44 
		% ---
			\label{MOSIAH-29-44}

			And thus commenced the reign of the judges throughout all the land of Zarahemla
			among all the people which was called the Nephites;
			and Alma was the first and chief judge.

		% ---
		\paragraph{Mosiah 29:45} % *MOSIAH-29-45 
		% ---
			\label{MOSIAH-29-45}

			And now it came to pass that his father died,
			being eighty and two years old,
			having lived to fulfill the commandments of God.

		% ---
		\paragraph{Mosiah 29:46} % *MOSIAH-29-46 
		% ---
			\label{MOSIAH-29-46}

			And it came to pass that
			Mosiah died also, in the thirty and third year of his reign,
			being sixty and three years old,
			making in the whole five hundred and nine years
			from the time Lehi left Jerusalem.

		% ---
		\paragraph{Mosiah 29:47} % *MOSIAH-29-47 
		% ---
			\label{MOSIAH-29-47}

			And thus ended the reign of the kings over the people of Nephi;
			and thus ended the days of Alma,
			who was the founder of their church.